% block: nonhermitian-systems-heading page:111 type:section
\subsection{\BT{Non-Hermitian Linear Systems}{非 Hermitian 线性方程组}}\label{sec:nonhermitian-linear-systems}
\BT{Let us recall that the CG method and the CR method (or the MINRES method) have the following two favorable properties:}{回顾 CG 方法和 CR 方法 (或 MINRES 方法) 具有的以下两个良好性质:}
\begin{enumerate}
\item\BT{The residual vectors satisfy Ritz-Galerkin approach \eqref{eq:ritz-galerkin-orthogonality} or minimal residual approach \eqref{eq:minimal-residual-problem}.}{残差向量满足 Ritz-Galerkin 途径 \eqref{eq:ritz-galerkin-orthogonality} 或最小残差途径 \eqref{eq:minimal-residual-problem}.}
\item\BT{The residual vectors are generated by short-term recurrences.}{残差向量由短递推生成.}
\end{enumerate}
\BT{As described in Sect.\ \ref{sec:hermitian-linear-systems}, the first property determines the optimality of approximate solutions, and the second property plays an important role in keeping the computational costs and memory requirements constant at each iteration step. Recalling that the Ritz--Galerkin approach \eqref{eq:ritz-galerkin-orthogonality} is for the CG method, and the minimal residual approach \eqref{eq:minimal-residual-problem} is for the CR method and the MINRES method.}{如 Section\ \ref{sec:hermitian-linear-systems} 所述, 第一个性质决定近似解的最优性, 第二个性质在使每次迭代的计算开销和内存需求保持不变方面起重要作用. 回顾 Ritz--Galerkin 途径 \eqref{eq:ritz-galerkin-orthogonality} 对应于 CG 方法, 最小残差途径 \eqref{eq:minimal-residual-problem} 对应于 CR 方法和 MINRES 方法.}
% block: faber-manteuffel-limitation page:112 type:ordinary-paragraph
\BT{Here, a question arises: are there Krylov subspace methods satisfying both the properties for non-Hermitian linear systems? The answer is negative, which was proved in \BookCite{57} and known as the Faber--Manteuffel theorem. More concretely, if matrix $A$ has the form}{这里产生一个问题: 对于非 Hermitian 线性方程组, 是否存在同时满足这两个性质的 Krylov 子空间方法? 答案是否定的, 这已在 \BookCite{57} 中证明, 称为 Faber--Manteuffel 定理. 更具体地, 若矩阵 $A$ 具有形式}
\[A=e^{i\phi}(icI+G)\quad(c\ge0,\ 0\le\phi\le2\pi,\ B^{-1}G^{\mathrm H}B=G),\]
\BT{then algorithms satisfying both properties can be constructed. For details, see \BookCite{57} and Greenbaum's book \BookCite{81}, Chap.6. See also \BookCite{121} and the references therein.}{则可构造同时满足这两个性质的算法. 详见 \BookCite{57} 和 Greenbaum 的著作 \BookCite{81}, Chapter 6. 另见 \BookCite{121} 及其中的参考文献.}
% block: nonhermitian-normal-equations pages:112,113 type:ordinary-paragraph
\BT{For solving non-Hermitian linear systems, one may consider the following transformed linear systems:}{为求解非 Hermitian 线性方程组, 可以考虑如下变换后的线性方程组:}
\begin{equation}\label{eq:nonhermitian-normal-equations}A^{\mathrm H}A\bm x=A^{\mathrm H}\bm b.\end{equation}
\BT{If $A$ is nonsingular, then $A^{\mathrm H}A$ is Hermitian positive definite. Thus the CG method can be applied to the transformed linear systems. However, the condition number of $A^{\mathrm H}A$ may be twice as large as the condition number of the original matrix $A$. Hence, from Theorem\ \ref{thm:cg-condition-number-bound}, the speed of convergence may be slow. Furthermore, if the condition number is very large, then it follows from the error analysis \eqref{eq:relative-perturbation-bound}, the accuracy of an obtained approximate solution may be very low. Note that the error analysis \eqref{eq:relative-perturbation-bound} does not depend on numerical algorithms for solving linear systems. This means that the accuracy of an approximate solution computed by any numerical algorithm may become low.}{若 $A$ 非奇异, 则 $A^{\mathrm H}A$ 为 Hermitian 正定矩阵. 因而可对变换后的线性方程组应用 CG 方法. 然而, $A^{\mathrm H}A$ 的条件数可能为原矩阵 $A$ 的条件数的两倍. 因此, 由 Theorem\ \ref{thm:cg-condition-number-bound}, 收敛速度可能较慢. 此外, 若条件数很大, 则由误差分析 \eqref{eq:relative-perturbation-bound} 可知, 所得近似解的精度可能很低. 注意, 误差分析 \eqref{eq:relative-perturbation-bound} 不依赖于求解线性方程组的数值算法. 这意味着, 由任何数值算法计算的近似解的精度都可能变低.}
\BT{On the other hand, if the condition number $A^{\mathrm H}A$ is not so large, there are cases where applying the CG method to \eqref{eq:nonhermitian-normal-equations} is a method of choice. In fact, if $A$ is a unitary matrix, then $A^{\mathrm H}A=I$, which indicates that the CG method gives the solution within only one iteration. Thus, we can guess that this approach is particularly useful for the case where $A$ is close to a unitary matrix. Well-known algorithms in this approach are the CGNE method and the CGNR method, which will be described in Sect.\ \ref{subsec:cg-normal-equations}.}{另一方面, 若 $A^{\mathrm H}A$ 的条件数不很大, 则有些情况下对 \eqref{eq:nonhermitian-normal-equations} 应用 CG 方法是一种可取的方法. 事实上, 若 $A$ 为酉矩阵, 则 $A^{\mathrm H}A=I$, 这表明 CG 方法只需一次迭代就给出解. 因此可以推测, 当 $A$ 接近酉矩阵时, 这一途径尤其有用. 这一途径中著名的算法为 CGNE 方法和 CGNR 方法, 将在 Section\ \ref{subsec:cg-normal-equations} 中介绍.}
\BT{In what follows, Krylov subspace methods for non-Hermitian linear systems are derived. We will see that the BiCG method, the BiCR method, and the QMR method can be regarded as extensions of the CG method, the CR method, and the MINRES method respectively, and these are based on the bi-Lanczos process (or a bi-Lanczos-like process), and the GMRES method can be regarded as an extension of the MINRES method, and the GMRES method is based on the Arnoldi process. We will also see how the BiCG method has been improved via product-type Krylov subspace methods and the IDR($s$) method.}{下面推导非 Hermitian 线性方程组的 Krylov 子空间方法. 将看到 BiCG 方法, BiCR 方法和 QMR 方法可分别视为 CG 方法, CR 方法和 MINRES 方法的推广, 它们基于双 Lanczos 过程 (或类似双 Lanczos 的过程); GMRES 方法可视为 MINRES 方法的推广, 且 GMRES 方法基于 Arnoldi 过程. 还将看到如何通过乘积型 Krylov 子空间方法和 IDR($s$) 方法改进 BiCG 方法.}
% block: bicg-heading page:113 type:subsection
\subsubsection{\BT{The Bi-Conjugate Gradient (BiCG) Method}{双共轭梯度 (BiCG) 方法}}\label{subsec:bicg-method}
\BT{Similar to the derivation process of the CG method in Sect.\ \ref{subsec:conjugate-gradient}, the BiCG method \BookCite{58}, \BookCite{119} can be derived by the bi-Lanczos process in Section\ \ref{subsec:bi-lanczos-process}.}{与 Section\ \ref{subsec:conjugate-gradient} 中 CG 方法的推导过程类似, BiCG 方法 \BookCite{58}, \BookCite{119} 可由 Section\ \ref{subsec:bi-lanczos-process} 中的双 Lanczos 过程导出.}
\BT{In this subsection, we give a concise way of the derivation to see that the BiCG method for non-Hermitian linear systems can be derived from the preconditioned COCG method (Algorithm\ \ref{alg:pcocg}). First, we consider the following $2N\times2N$ complex symmetric linear system:}{本小节给出一种简洁的推导方式, 以说明非 Hermitian 线性方程组的 BiCG 方法可从预条件 COCG 方法 (Algorithm\ \ref{alg:pcocg}) 导出. 首先考虑下列 $2N\times2N$ 复对称线性方程组:}
\begin{equation}\label{eq:bicg-augmented-system}
\begin{bmatrix}O&A\\A^{\top}&O\end{bmatrix}\begin{bmatrix}\bar{\bm x}^{*}\\\bm x\end{bmatrix}=\begin{bmatrix}\bm b\\\bar{\bm b}^{*}\end{bmatrix},\quad\text{or }\widetilde A\widetilde{\bm x}=\widetilde{\bm b}.
\end{equation}
\BT{We can see from the system \eqref{eq:bicg-augmented-system} that it is mathematically equivalent to}{由方程组 \eqref{eq:bicg-augmented-system} 可看出, 它在数学上等价于}
\[A\bm x=\bm b\quad\text{and}\quad A^{\top}\bar{\bm x}^{*}=\bar{\bm b}^{*}\quad(\Leftrightarrow A^{\mathrm H}\bm x^{*}=\bm b^{*}).\]
\BT{If we apply the algorithm of COCG with the preconditioner}{若采用下列预条件子, 将 COCG 算法}
\begin{equation}\label{eq:bicg-swap-preconditioner}\widetilde M=\begin{bmatrix}O&I\\I&O\end{bmatrix},\quad I:\text{ identity matrix},\end{equation}
\BT{to \eqref{eq:bicg-augmented-system}, then the resulting algorithm at the $n$th iteration step can be written as}{应用于 \eqref{eq:bicg-augmented-system}, 则所得算法在第 $n$ 次迭代时可写为}
% block: bicg-augmented-cocg-recurrences page:114 type:ordinary-paragraph
\[\begin{aligned}
\widetilde{\bm p}_n^{\mathrm{COCG}}&=\widetilde M^{-1}\widetilde{\bm r}_n^{\mathrm{COCG}}+\beta_{n-1}\widetilde{\bm p}_{n-1}^{\mathrm{COCG}},\\
\alpha_n&=\frac{(\overline{\widetilde M^{-1}\widetilde{\bm r}_n^{\mathrm{COCG}}},\widetilde{\bm r}_n^{\mathrm{COCG}})}{(\overline{\widetilde{\bm p}_n^{\mathrm{COCG}}},\widetilde A\widetilde{\bm p}_n^{\mathrm{COCG}})},\\
\widetilde{\bm x}_{n+1}^{\mathrm{COCG}}&=\widetilde{\bm x}_n^{\mathrm{COCG}}+\alpha_n\widetilde{\bm p}_n^{\mathrm{COCG}},\\
\widetilde{\bm r}_{n+1}^{\mathrm{COCG}}&=\widetilde{\bm r}_n^{\mathrm{COCG}}-\alpha_n\widetilde A\widetilde{\bm p}_n^{\mathrm{COCG}},\\
\beta_n&=\frac{(\overline{\widetilde M^{-1}\widetilde{\bm r}_{n+1}^{\mathrm{COCG}}},\widetilde{\bm r}_{n+1}^{\mathrm{COCG}})}{(\overline{\widetilde M^{-1}\widetilde{\bm r}_n^{\mathrm{COCG}}},\widetilde{\bm r}_n^{\mathrm{COCG}})}.
\end{aligned}\]
\BT{Substituting $\widetilde M^{-1}(=\widetilde M)$ of \eqref{eq:bicg-swap-preconditioner} and the vectors}{将 \eqref{eq:bicg-swap-preconditioner} 中的 $\widetilde M^{-1}(=\widetilde M)$ 及向量}
\[\widetilde{\bm x}_n^{\mathrm{COCG}}:=\begin{bmatrix}\bm x_n\\\bar{\bm x}_n^{*}\end{bmatrix},\quad\widetilde{\bm r}_n^{\mathrm{COCG}}:=\begin{bmatrix}\bm r_n\\\bar{\bm r}_n^{*}\end{bmatrix},\quad\widetilde{\bm p}_n^{\mathrm{COCG}}:=\begin{bmatrix}\bar{\bm p}_n^{*}\\\bm p_n\end{bmatrix}\]
\BT{into the previous recurrences, and using the following results:}{代入前述递推式, 并使用下列结果:}
\[\begin{aligned}
(\overline{\widetilde M^{-1}\widetilde{\bm r}_n^{\mathrm{COCG}}},\widetilde{\bm r}_n^{\mathrm{COCG}})&=\begin{bmatrix}\bar{\bm r}_n^{*\top}&\bm r_n^{\top}\end{bmatrix}\begin{bmatrix}\bm r_n\\\bar{\bm r}_n^{*}\end{bmatrix}\\
&=\bar{\bm r}_n^{*\top}\bm r_n+\bm r_n^{\top}\bar{\bm r}_n^{*}\\
&=\bm r_n^{*\mathrm H}\bm r_n+(\bm r_n^{\top}\bar{\bm r}_n^{*})^{\top}\\
&=(\bm r_n^{*},\bm r_n)+\bm r_n^{*\mathrm H}\bm r_n\\
&=2(\bm r_n^{*},\bm r_n)\\
(\overline{\widetilde{\bm p}_n^{\mathrm{COCG}}},\widetilde A\widetilde{\bm p}_n^{\mathrm{COCG}}),&=\begin{bmatrix}((\bar{\bm p}_n^{*})^{\top}\bm p_n)^{\top}\end{bmatrix}\begin{bmatrix}A\bm p_n\\A^{\top}\bar{\bm p}_n^{*}\end{bmatrix}\\
&=(\bm p_n^{*})^{\mathrm H}A\bm p_n+\bm p_n^{\top}A^{\top}\bar{\bm p}_n^{*},\\
&=(\bm p_n^{*})^{\mathrm H}A\bm p_n+(\bm p_n^{\top}A^{\top}\bar{\bm p}_n^{*})^{\top},\\
&=2(\bm p_n^{*},A\bm p_n),
\end{aligned}\]
\BT{we readily obtain the algorithm of BiCG for solving non-Hermitian linear systems described in Algorithm\ \ref{alg:bicg}. If matrix $A$ is real non-symmetric, the aforementioned derivation process corresponds to the derivation in van der Vorst's book \BookCite{196}, Chap.7, where the BiCG method for real non-symmetric linear systems is derived from the CG method.}{就容易得到 Algorithm\ \ref{alg:bicg} 所述的求解非 Hermitian 线性方程组的 BiCG 算法. 若矩阵 $A$ 为实非对称矩阵, 则上述推导过程对应于 van der Vorst 的著作 \BookCite{196}, Chapter 7 中的推导, 其中实非对称线性方程组的 BiCG 方法由 CG 方法导出.}
\BT{It can be seen from Algorithm\ \ref{alg:bicg} that the choice $\bm r_0^{*}=\bm r_0$ reduces to the CG method if the coefficient matrix $A$ is Hermitian.}{由 Algorithm\ \ref{alg:bicg} 可见, 若系数矩阵 $A$ 为 Hermitian 矩阵, 则选择 $\bm r_0^{*}=\bm r_0$ 就退化为 CG 方法.}
% block: bicg-biorthogonality-proposition page:114 type:proposition
\BT{Similar to the CG iterates as shown in Proposition\ \ref{prop:cg-orthogonality}, the BiCG iterates have the following relations:}{与 Proposition\ \ref{prop:cg-orthogonality} 所示的 CG 迭代类似, BiCG 迭代具有下列关系:}
\par\Needspace{3\baselineskip}
\begin{Proposition}\label{prop:bicg-biorthogonality}
\BT{Let $\bm p_n,\bm r_n$ be the vectors in the BiCG method. Then,}{设 $\bm p_n,\bm r_n$ 为 BiCG 方法中的向量. 则,}
\begin{enumerate}\item $(\bm r_i^{*},\bm r_j)=0\quad\text{for }i\ne j$,\item $(\bm p_i^{*},A\bm p_j)=0\quad\text{for }i\ne j$.\end{enumerate}
\end{Proposition}
% block: bicg-algorithm page:115 type:ordinary-paragraph
\begin{algorithm}[H]\caption{\BT{The BiCG method}{BiCG 方法}}\label{alg:bicg}\begin{algorithmic}[1]
\Statex\textbf{Input:} $\bm x_0\in\mathbb C^N,\ \beta_{-1}=0,\ \bm p_{-1}=\bm p_{-1}^{*}=\bm0,\ \bm r_0=\bm b-A\bm x_0$
\Statex\textbf{Input:} \BT{Choose $\bm r_0^{*}\in\mathbb C^N$, e.g., $\bm r_0^{*}=\bm r_0$}{选择 $\bm r_0^{*}\in\mathbb C^N$, 例如 $\bm r_0^{*}=\bm r_0$}
\Statex\textbf{Output:} $\bm x_n$
\For{$n=0,1,2,\ldots$, \BT{until convergence}{直至收敛}}
\State $\bm p_n=\bm r_n+\beta_{n-1}\bm p_{n-1},\quad\bm p_n^{*}=\bm r_n^{*}+\bar\beta_{n-1}\bm p_{n-1}^{*}$
\State $\alpha_n=\frac{(\bm r_n^{*},\bm r_n)}{(\bm p_n^{*},A\bm p_n)}$
\State $\bm x_{n+1}=\bm x_n+\alpha_n\bm p_n$
\State $\bm r_{n+1}=\bm r_n-\alpha_nA\bm p_n,\quad\bm r_{n+1}^{*}=\bm r_n^{*}-\bar\alpha_nA^{\mathrm H}\bm p_n^{*}$
\State $\beta_n=\frac{(\bm r_{n+1}^{*},\bm r_{n+1})}{(\bm r_n^{*},\bm r_n)}$
\EndFor\end{algorithmic}\end{algorithm}
% block: bicg-krylov-corollary page:115 type:corollary
\BT{From Proposition\ \ref{prop:bicg-biorthogonality}, the following properties hold true:}{由 Proposition\ \ref{prop:bicg-biorthogonality}, 下列性质成立:}
\par\Needspace{3\baselineskip}
\begin{Corollary}\label{cor:bicg-krylov-orthogonality}
\BT{Let $\bm p_n,\bm r_n$ be the vectors in the BiCG method. Then,}{设 $\bm p_n,\bm r_n$ 为 BiCG 方法中的向量. 则,}
\begin{enumerate}\item $\bm r_n\perp\mathcal K_n(A^{\mathrm H},\bm r_0^{*})$,\item $A\bm p_n\perp\mathcal K_n(A^{\mathrm H},\bm r_0^{*})$.\end{enumerate}
\end{Corollary}
% block: bicg-alternative-alpha page:115 type:ordinary-paragraph
\BT{In what follows, we give the other computational formulas for $\alpha_n$ and $\beta_n$ of the BiCG method, which will be used to derive the product-type Krylov subspace methods in Section\ \ref{subsec:product-type-krylov}.}{下面给出 BiCG 方法中 $\alpha_n$ 和 $\beta_n$ 的其他计算公式, 它们将用于推导 Section\ \ref{subsec:product-type-krylov} 中的乘积型 Krylov 子空间方法.}
\BT{Since $\bm r_n^{*}\in\mathcal K_{n+1}(A^{\mathrm H},\bm r_0^{*})$, the vector $\bm r_n^{*}$ can be written as}{由于 $\bm r_n^{*}\in\mathcal K_{n+1}(A^{\mathrm H},\bm r_0^{*})$, 向量 $\bm r_n^{*}$ 可写为}
\[\begin{aligned}\bm r_n^{*}&=c_n(A^{\mathrm H})^n\bm r_0^{*}+\underbrace{c_{n-1}(A^{\mathrm H})^{n-1}\bm r_0^{*}+\cdots+c_0\bm r_0^{*}}_{\bm z}\\&=c_n(A^{\mathrm H})^n\bm r_0^{*}+\bm z,\quad\bm z\in\mathcal K_n(A^{\mathrm H},\bm r_0^{*}).\end{aligned}\]
\BT{From the above equation and Corollary\ \ref{cor:bicg-krylov-orthogonality}, it follows that}{由上式及 Corollary\ \ref{cor:bicg-krylov-orthogonality} 得}
\[\begin{aligned}(\bm r_n^{*},\bm r_n)&=(c_n(A^{\mathrm H})^n\bm r_0^{*}+\bm z,\bm r_n)=\bar c_n((A^{\mathrm H})^n\bm r_0^{*},\bm r_n)+(\bm z,\bm r_n)\\&=\bar c_n((A^{\mathrm H})^n\bm r_0^{*},\bm r_n)\\
(\bm r_n^{*},A\bm p_n)&=(c_n(A^{\mathrm H})^n\bm r_0^{*}+\bm z,A\bm p_n)=\bar c_n((A^{\mathrm H})^n\bm r_0^{*},A\bm p_n)+(\bm z,A\bm p_n)\\&=\bar c_n((A^{\mathrm H})^n\bm r_0^{*},A\bm p_n).
\end{aligned}\]
\BT{Thus, $\alpha_n$ of the BiCG method (Algorithm\ \ref{alg:bicg}) can be rewritten as}{因此, BiCG 方法 (Algorithm\ \ref{alg:bicg}) 的 $\alpha_n$ 可改写为}
\begin{equation}\label{eq:bicg-alpha-shadow-power}\alpha_n=\frac{(\bm r_n^{*},\bm r_n)}{(\bm p_n^{*},A\bm p_n)}=\frac{((A^{\mathrm H})^n\bm r_0^{*},\bm r_n)}{((A^{\mathrm H})^n\bm r_0^{*},A\bm p_n)}.\end{equation}
% block: bicg-alternative-beta pages:115,116 type:ordinary-paragraph
\BT{Similarly, we give the other computational formulas for $\beta_n$. From Corollary\ \ref{cor:bicg-krylov-orthogonality}, two vectors $(A^{\mathrm H})^n\bm r_0^{*}\in\mathcal K_{n+1}(A^{\mathrm H},\bm r_0^{*})$ and $A\bm p_{n+1}$ are orthogonal, and thus}{类似地, 给出 $\beta_n$ 的其他计算公式. 由 Corollary\ \ref{cor:bicg-krylov-orthogonality}, 两个向量 $(A^{\mathrm H})^n\bm r_0^{*}\in\mathcal K_{n+1}(A^{\mathrm H},\bm r_0^{*})$ 和 $A\bm p_{n+1}$ 正交, 因此}
\[\begin{aligned}0&=((A^{\mathrm H})^n\bm r_0^{*},A\bm p_{n+1})\\&=((A^{\mathrm H})^n\bm r_0^{*},A\bm r_{n+1}+\beta_nA\bm p_n)\\&=((A^{\mathrm H})^n\bm r_0^{*},A\bm r_{n+1})+\beta_n((A^{\mathrm H})^n\bm r_0^{*},A\bm p_n)\\&=((A^{\mathrm H})^{n+1}\bm r_0^{*},\bm r_{n+1})+\beta_n((A^{\mathrm H})^n\bm r_0^{*},A\bm p_n),\end{aligned}\]
\BT{which immediately leads to}{立即得到}
\[\beta_n=-\frac{((A^{\mathrm H})^{n+1}\bm r_0^{*},\bm r_{n+1})}{((A^{\mathrm H})^n\bm r_0^{*},A\bm p_n)}.\]
\BT{Using \eqref{eq:bicg-alpha-shadow-power} together with the above result yields}{结合 \eqref{eq:bicg-alpha-shadow-power} 与上述结果, 得}
\begin{equation}\label{eq:bicg-beta-shadow-power}\beta_n=-\alpha_n\frac{((A^{\mathrm H})^{n+1}\bm r_0^{*},\bm r_{n+1})}{((A^{\mathrm H})^n\bm r_0^{*},\bm r_n)}.\end{equation}
% block: composite-step-bicg-heading page:116 type:subsection
\subsubsection{\BT{The Composite Step Bi-Conjugate Gradient (CSBiCG) Method}{复合步双共轭梯度 (CSBiCG) 方法}}\label{subsec:composite-step-bicg}
\BT{From the algorithm of the BiCG method (Algorithm\ \ref{alg:bicg}), there are two kinds of possible (near) breakdown:}{由 BiCG 方法的算法 (Algorithm\ \ref{alg:bicg}), 有两类可能的 (近) 中断:}
\begin{enumerate}
\item $\rho_n:=(\bm r_n^{*},\bm r_n)=0$ \BT{for $\bm r_n\ne\bm0$ (Lanczos breakdown),}{而 $\bm r_n\ne\bm0$ (Lanczos 中断),}
\item $\sigma_n:=(\bm p_n^{*},A\bm p_n)=0$ (\BT{pivot breakdown}{主元中断}).
\end{enumerate}
\BT{The Lanczos breakdown can be circumvented by the look-ahead Lanczos process in \BookCite{145}, \BookCite{189}. Bank and Chan focused on the pivot breakdown and proposed the Composite Step Bi-Conjugate Gradient (CSBiCG) method to cure the pivot breakdown \BookCite{16}, \BookCite{17} under the assumption that the Lanczos breakdown does not occur. Below is a brief explanation of the algorithm that is based on \BookCite{17}.}{Lanczos 中断可通过 \BookCite{145}, \BookCite{189} 中的前瞻 Lanczos 过程避免. Bank 和 Chan 关注主元中断, 在假设不发生 Lanczos 中断的条件下, 提出复合步双共轭梯度 (CSBiCG) 方法来消除主元中断 \BookCite{16}, \BookCite{17}. 下面根据 \BookCite{17} 对该算法作简要说明.}
% block: composite-step-auxiliary-vectors pages:116,117 type:ordinary-paragraph
\BT{If $\sigma_n=0$, then $\bm r_{n+1}=\bm r_n-(\rho_n/\sigma_n)A\bm p_n$ and $\bm r_{n+1}^{*}=\bm r_n^{*}-\overline{(\rho_n/\sigma_n)}A^{\mathrm H}\bm p_n^{*}$ cannot be obtained. Consider auxiliary vectors:}{若 $\sigma_n=0$, 就无法得到 $\bm r_{n+1}=\bm r_n-(\rho_n/\sigma_n)A\bm p_n$ 和 $\bm r_{n+1}^{*}=\bm r_n^{*}-\overline{(\rho_n/\sigma_n)}A^{\mathrm H}\bm p_n^{*}$. 考虑辅助向量:}
\[\begin{aligned}\bm z_{n+1}&=\sigma_n\bm r_n-\rho_nA\bm p_n,\\\bm z_{n+1}^{*}&=\bar\sigma_n\bm r_n-\bar\rho_nA^{\mathrm H}\bm p_n^{*}.\end{aligned}\]
\BT{Then, $\bm z_{n+1}$ and $\bm z_{n+1}^{*}$ exist even if $\sigma_n=0$, and $\bm z_{n+1}$ and $\bm z_{n+1}^{*}$ belong to $\mathcal K_{n+2}(A,\bm r_0)$ and $\mathcal K_{n+2}(A^{\mathrm H},\bm r_0^{*})$ respectively. The composite step technique then considers $\bm r_{n+2}$ and $\bm r_{n+2}^{*}$ as follows:}{于是即使 $\sigma_n=0$, $\bm z_{n+1}$ 和 $\bm z_{n+1}^{*}$ 也存在, 且 $\bm z_{n+1}$ 和 $\bm z_{n+1}^{*}$ 分别属于 $\mathcal K_{n+2}(A,\bm r_0)$ 和 $\mathcal K_{n+2}(A^{\mathrm H},\bm r_0^{*})$. 随后复合步技术按下式考虑 $\bm r_{n+2}$ 和 $\bm r_{n+2}^{*}$:}
\[\begin{aligned}\bm r_{n+2}&=\bm r_n-c_1A\bm p_n-c_2A\bm z_{n+1},\\\bm r_{n+2}^{*}&=\bm r_n^{*}-c_1^{*}A^{\mathrm H}\bm p_n^{*}-c_2^{*}A^{\mathrm H}\bm z_{n+1}^{*},\end{aligned}\]
\BT{where $c_1,c_2,c_1^{*},c_2^{*}$ are determined so that}{其中 $c_1,c_2,c_1^{*},c_2^{*}$ 确定为使得}
\[\bm r_{n+2}\perp\operatorname{span}\{\bm p_n^{*},\bm z_{n+1}^{*}\},\qquad\bm r_{n+2}^{*}\perp\operatorname{span}\{\bm p_n,\bm z_{n+1}\},\]
\BT{i.e., $(\bm r_{n+2},\bm p_n^{*})=(\bm r_{n+2},\bm z_{n+1}^{*})=(\bm r_{n+2}^{*},\bm p_n)=(\bm r_{n+2}^{*},\bm z_{n+1})=0$. Similarly, let}{即 $(\bm r_{n+2},\bm p_n^{*})=(\bm r_{n+2},\bm z_{n+1}^{*})=(\bm r_{n+2}^{*},\bm p_n)=(\bm r_{n+2}^{*},\bm z_{n+1})=0$. 类似地, 令}
\[\begin{aligned}\bm p_{n+2}&=\bm r_{n+2}+d_1\bm p_n+d_2\bm z_{n+1},\\\bm p_{n+2}^{*}&=\bm r_{n+2}^{*}+d_1^{*}\bm p_n^{*}+d_2^{*}\bm z_{n+1}^{*},\end{aligned}\]
\BT{and $d_1,d_2,d_1^{*},d_2^{*}$ are determined so that}{并确定 $d_1,d_2,d_1^{*},d_2^{*}$ 使得}
\[A\bm p_{n+2}\perp\operatorname{span}\{\bm p_n^{*},\bm z_{n+1}^{*}\},\qquad A^{\mathrm H}\bm p_{n+2}^{*}\perp\operatorname{span}\{\bm p_n,\bm z_{n+1}\}.\]
\BT{Then, it is shown in \BookCite{17}, Theorem 4.4 that $(\bm r_i^{*},\bm r_j)=(\bm p_i^{*},A\bm p_j)=0$ for $i\ne j$, which means that the composite step residual vector $\bm r_{n+2}$ and $A\bm p_{n+2}$ satisfies}{于是 \BookCite{17}, Theorem 4.4 说明, 当 $i\ne j$ 时有 $(\bm r_i^{*},\bm r_j)=(\bm p_i^{*},A\bm p_j)=0$, 这意味着复合步残差向量 $\bm r_{n+2}$ 和 $A\bm p_{n+2}$ 满足}
\[\bm r_{n+2}\perp\mathcal K_{n+2}(A^{\mathrm H},\bm r_0^{*}),\quad A\bm p_{n+2}\perp\mathcal K_{n+2}(A^{\mathrm H},\bm r_0^{*}).\]
\BT{which are the same properties that the BiCG method has, see Corollary\ \ref{cor:bicg-krylov-orthogonality}. The resulting algorithm is described in Algorithm\ \ref{alg:composite-step-bicg}.}{这些与 BiCG 方法具有的性质相同, 见 Corollary\ \ref{cor:bicg-krylov-orthogonality}. 所得算法见 Algorithm\ \ref{alg:composite-step-bicg}.}
% block: composite-step-choice page:117 type:ordinary-paragraph
\BT{From Algorithm\ \ref{alg:composite-step-bicg}, the number of matrix--vector multiplications is the same as that of the BiCG method. The composite step techniques are used for other Krylov subspace methods \BookCite{35}, \BookCite{113}.}{由 Algorithm\ \ref{alg:composite-step-bicg}, 矩阵向量乘法的次数与 BiCG 方法相同. 复合步技术也用于其他 Krylov 子空间方法 \BookCite{35}, \BookCite{113}.}
\BT{At each step in Algorithm\ \ref{alg:composite-step-bicg}, one needs to choose a $1\times1$ step (BiCG step) or a $2\times2$ step (composite step). In \BookCite{17}, the following choice is proposed:}{在 Algorithm\ \ref{alg:composite-step-bicg} 的每一步, 需要选择 $1\times1$ 步 (BiCG 步) 或 $2\times2$ 步 (复合步). 文献 \BookCite{17} 提出如下选择:}
\begin{algorithmic}[1]
\If{$\|\bm z_{n+1}\|\le\|\bm r_n\||\sigma_n|$}
\State $1\times1$ \BT{step}{步}
\Else
\State $\nu_{n+2}=\|\delta_n\bm r_n-\rho_n^3\zeta_{n+1}\bm q_n-\theta_{n+1}\rho_n^2\bm y_{n+1}\|$
\If{$\nu_{n+2}|\sigma_n|<\|\bm z_{n+1}\||\delta_n|$}
\State $2\times2$ \BT{step}{步}
\Else\State $1\times1$ \BT{step}{步}\EndIf\EndIf
\end{algorithmic}
\BT{This choice leads to not only avoiding (near) pivot breakdowns but also some smoothing of the convergence history in residual norms.}{这一选择不仅避免 (近) 主元中断, 也使残差范数的收敛历程得到一定程度的平滑.}
% block: composite-step-bicg-algorithm page:118 type:ordinary-paragraph
\begin{algorithm}[H]\caption{\BT{The Composite Step BiCG method}{复合步 BiCG 方法}}\label{alg:composite-step-bicg}\begin{algorithmic}[1]
\Statex\textbf{Input:} $\bm x_0\in\mathbb C^N,\ \bm r_0=\bm b-A\bm x_0$
\Statex\textbf{Input:} \BT{Choose $\bm r_0^{*}\in\mathbb C^N$, e.g., $\bm r_0^{*}=\bm r_0$}{选择 $\bm r_0^{*}\in\mathbb C^N$, 例如 $\bm r_0^{*}=\bm r_0$}
\Statex\textbf{Output:} $\bm x_n$
\State $\bm p_0=\bm r_0,\quad\bm p_0^{*}=\bm r_0^{*}$
\State $\bm q_0=A\bm p_0,\quad\bm q_0^{*}=A^{\mathrm H}\bm p_0^{*}$
\State $\rho_0=(\bm r_0^{*},\bm r_0)$
\State $n=0$
\While{\BT{until convergence}{直至收敛}}
\State $\sigma_n=(\bm p_n^{*},\bm q_n)$
\State $\bm z_{n+1}=\sigma_n\bm r_n-\rho_n\bm q_n,\quad\bm z_{n+1}^{*}=\bar\sigma_n\bm r_n^{*}-\bar\rho_n\bm q_n^{*}$
\State $\bm y_{n+1}=A\bm z_{n+1},\quad\bm y_{n+1}^{*}=A^{\mathrm H}\bm z_{n+1}^{*}$
\State $\theta_{n+1}=(\bm z_{n+1}^{*},\bm z_{n+1})$
\State $\zeta_{n+1}=(\bm z_{n+1}^{*},\bm y_{n+1})$
\If{$1\times1$ \BT{step}{步}}
\State $\alpha=\rho_n/\sigma_n$
\State $\rho_{n+1}=\theta_{n+1}/\sigma_n^2$
\State $\beta_{n+1}=\rho_{n+1}/\rho_n$
\State $\bm x_{n+1}=\bm x_n+\alpha_n\bm p_n$
\State $\bm r_{n+1}=\bm r_n-\alpha_n\bm q_n,\quad\bm r_{n+1}^{*}=\bm r_n^{*}-\bar\alpha_n\bm q_n^{*}$
\State $\bm p_{n+1}=\bm z_{n+1}/\sigma_n+\beta_{n+1}\bm p_n,\quad\bm p_{n+1}^{*}=\bm z_{n+1}^{*}/\bar\sigma_n+\bar\beta_{n+1}\bm p_n^{*}$
\State $\bm q_{n+1}=\bm y_{n+1}/\sigma_n+\beta_{n+1}\bm q_n,\quad\bm q_{n+1}^{*}=\bm y_{n+1}^{*}/\bar\sigma_n+\bar\beta_{n+1}\bm q_n^{*}$
\State $n=n+1$
\EndIf
\If{$2\times2$ \BT{step}{步}}
\State $\delta_n=\sigma_n\zeta_{n+1}\rho_n^2-\theta_{n+1}^2$
\State $\alpha_n=\zeta_{n+1}\rho_n^3/\delta_n,\quad\alpha_{n+1}=\theta_{n+1}\rho_n^2/\delta_n$
\State $\bm x_{n+2}=\bm x_n+\alpha_n\bm p_n+\alpha_{n+1}\bm z_{n+1}$
\State $\bm r_{n+2}=\bm r_n-\alpha_n\bm q_n-\alpha_{n+1}\bm y_{n+1},\quad\bm r_{n+2}^{*}=\bm r_n^{*}-\bar\alpha_n\bm q_n^{*}-\bar\alpha_{n+1}\bm y_{n+1}$
\State $\rho_{n+2}=(\bm r_{n+2}^{*},\bm r_{n+2})$
\State $\beta_{n+1}=\rho_{n+2}/\rho_n,\quad\beta_{n+2}=\rho_{n+2}\sigma_n/\theta_{n+1}$
\State $\bm p_{n+2}=\bm r_{n+2}+\beta_{n+1}\bm p_n+\beta_{n+2}\bm z_{n+1},\quad\bm p_{n+2}^{*}=\bm r_{n+2}^{*}+\bar\beta_{n+1}\bm p_n^{*}+\bar\beta_{n+2}\bm z_{n+1}^{*}$
\State $\bm q_{n+2}=A\bm p_{n+2},\quad\bm q_{n+2}^{*}=A^{\mathrm H}\bm p_{n+2}^{*}$
\State $n=n+2$
\EndIf\EndWhile\end{algorithmic}\end{algorithm}
% block: composite-bicg-best-approximation pages:118,119 type:theorem
\BT{The following theorem shows the best approximation result of the (composite step) BiCG method:}{下述定理给出 (复合步) BiCG 方法的最佳逼近结果:}
\par\Needspace{14\baselineskip}
\begin{Theorem}[\BookCite{16}]\label{thm:composite-bicg-best-approximation}
\BT{Consider applying the (composite step) BiCG method to real nonsymmetric linear systems with initial guess $\bm x_0=\bm0$ and $\bm r_0^{*}=\bm r_0$. Suppose that for all $\bm v,\bm w\in\mathbb R^N$, we have}{考虑对实非对称线性方程组应用 (复合步) BiCG 方法, 初始猜测为 $\bm x_0=\bm0$, 且 $\bm r_0^{*}=\bm r_0$. 假设对所有 $\bm v,\bm w\in\mathbb R^N$, 有}
\[|\bm w^{\top}A\bm v|\le\Gamma\|\bm v\|\|\bm w\|,\]
\BT{where $\Gamma$ is a constant independent of $\bm v$ and $\bm w$. Further, suppose that for those steps in the (composite step) BiCG method in which we compute an approximation $\bm x_n$, we have}{其中 $\Gamma$ 为与 $\bm v$ 和 $\bm w$ 无关的常数. 进一步假设, 对于 (复合步) BiCG 方法中计算近似解 $\bm x_n$ 的那些步骤, 有}
\[\inf_{\substack{\bm v\in\mathcal K_n(A,\bm r_0)\\\|\bm v\|=1}}\ \sup_{\substack{\bm w\in\mathcal K_n(A^{\top},\bm r_0^{*})\\\|\bm w\|\le1}}\bm w^{\top}A\bm v\ge\delta_n\ge\delta>0.\]
\BT{Then}{则}
\[\|\bm x-\bm x_n\|\le(1+\Gamma/\delta)\inf_{\bm v\in\mathcal K_n(A,\bm r_0)}\|\bm x-\bm v\|.\]
\end{Theorem}
\BT{Theorem\ \ref{thm:composite-bicg-best-approximation} is a simplified result of Theorem 4.1 in \BookCite{16}. From Theorem\ \ref{thm:composite-bicg-best-approximation}, we see that the (composite step) BiCG method produces an approximate solution close to the best approximate solution in error norm when $\Gamma/\delta$ is close to 0.}{Theorem\ \ref{thm:composite-bicg-best-approximation} 是 \BookCite{16} 中 Theorem 4.1 的简化结果. 由 Theorem\ \ref{thm:composite-bicg-best-approximation} 可见, 当 $\Gamma/\delta$ 接近 0 时, (复合步) BiCG 方法产生在误差范数下接近最佳近似解的近似解.}
% block: bicr-heading page:119 type:subsection
\subsubsection{\BT{The Bi-Conjugate Residual (BiCR) Method}{双共轭残差 (BiCR) 方法}}\label{subsec:bicr-method}
\BT{The BiCR method \BookCite{169} is an extension of the CR method to non-Hermitian linear systems.}{BiCR 方法 \BookCite{169} 是 CR 方法向非 Hermitian 线性方程组的推广.}\footnote{\BT{The method was mentioned in the unpublished paper \BookCite{38} by Chronopulous and Ma in 1989, and Dr. Chronopulous emailed the author about this fact on Apr. 6th, 2020.}{Chronopulous 和 Ma 于 1989 年的未发表论文 \BookCite{38} 中提到了该方法, Chronopulous 博士于 2020 年 4 月 6 日通过电子邮件将这一事实告知作者.}}
\BT{In this subsection, we show that the BiCR method can be derived from the preconditioned COCR method in Algorithm\ \ref{alg:pcocr}. When we apply the algorithm of the COCR method with the preconditioner \eqref{eq:bicg-swap-preconditioner} to \eqref{eq:bicg-augmented-system}, the resulting algorithm at the $n$th iteration step can be written as}{本小节说明 BiCR 方法可由 Algorithm\ \ref{alg:pcocr} 中的预条件 COCR 方法导出. 采用预条件子 \eqref{eq:bicg-swap-preconditioner}, 对 \eqref{eq:bicg-augmented-system} 应用 COCR 算法, 所得算法在第 $n$ 次迭代时可写为}
\[\begin{aligned}
\widetilde{\bm p}_n^{\mathrm{COCR}}&=\widetilde M^{-1}\widetilde{\bm r}_n^{\mathrm{COCR}}+\beta_{n-1}\widetilde{\bm p}_{n-1}^{\mathrm{COCR}},\\
\alpha_n&=\frac{(\overline{\widetilde M^{-1}\widetilde{\bm r}_n^{\mathrm{COCR}}},\widetilde A\widetilde M^{-1}\widetilde{\bm r}_n^{\mathrm{COCR}})}{(\overline{\widetilde A\widetilde{\bm p}_n^{\mathrm{COCR}}},\widetilde M^{-1}\widetilde A\widetilde{\bm p}_n^{\mathrm{COCR}})},\\
\widetilde{\bm x}_{n+1}^{\mathrm{COCR}}&=\widetilde{\bm x}_n^{\mathrm{COCR}}+\alpha_n\widetilde{\bm p}_n^{\mathrm{COCR}},\\
\widetilde{\bm r}_{n+1}^{\mathrm{COCR}}&=\widetilde{\bm r}_n^{\mathrm{COCR}}-\alpha_n\widetilde A\widetilde{\bm p}_n^{\mathrm{COCR}},\\
\beta_n&=\frac{(\overline{\widetilde M^{-1}\widetilde{\bm r}_{n+1}^{\mathrm{COCR}}},\widetilde A\widetilde M^{-1}\widetilde{\bm r}_{n+1}^{\mathrm{COCR}})}{(\overline{\widetilde M^{-1}\widetilde{\bm r}_n^{\mathrm{COCR}}},\widetilde A\widetilde M^{-1}\widetilde{\bm r}_n^{\mathrm{COCR}})}.
\end{aligned}\]
\BT{Substituting $\widetilde M^{-1}(=\widetilde M)$ of \eqref{eq:bicg-swap-preconditioner} and the vectors}{将 \eqref{eq:bicg-swap-preconditioner} 中的 $\widetilde M^{-1}(=\widetilde M)$ 以及向量}
\[\widetilde{\bm x}_n^{\mathrm{COCR}}:=\begin{bmatrix}\bm x_n\\\bar{\bm x}_n^{*}\end{bmatrix},\quad\widetilde{\bm r}_n^{\mathrm{COCR}}:=\begin{bmatrix}\bm r_n\\\bar{\bm r}_n^{*}\end{bmatrix},\quad\widetilde{\bm p}_n^{\mathrm{COCR}}:=\begin{bmatrix}\bar{\bm p}_n^{*}\\\bm p_n\end{bmatrix}\]
\BT{into the previous recurrences, and using the following results:}{代入前述递推式, 并使用下列结果:}
% block: bicr-bilinear-reductions page:120 type:ordinary-paragraph
\[\begin{aligned}
(\overline{\widetilde M^{-1}\widetilde{\bm r}_n^{\mathrm{COCR}}},\widetilde A\widetilde M^{-1}\widetilde{\bm r}_n^{\mathrm{COCR}})&=\begin{bmatrix}\bar{\bm r}_n^{*\top}&\bm r_n^{\top}\end{bmatrix}\begin{bmatrix}A\bm r_n\\A^{\top}\bar{\bm r}_n^{*}\end{bmatrix}\\
&=\bar{\bm r}_n^{*\top}A\bm r_n+\bm r_n^{\top}A^{\top}\bar{\bm r}_n^{*}\\
&=\bm r_n^{*\mathrm H}A\bm r_n+(\bm r_n^{\top}A^{\top}\bar{\bm r}_n^{*})^{\top}\\
&=(\bm r_n^{*},A\bm r_n)+\bm r_n^{*\mathrm H}A\bm r_n\\
&=2(\bm r_n^{*},A\bm r_n)
\end{aligned}\]
\BT{and}{以及}
\[\begin{aligned}
(\overline{\widetilde A\widetilde{\bm p}_n^{\mathrm{COCR}}},\widetilde M^{-1}\widetilde A\widetilde{\bm p}_n^{\mathrm{COCR}})&=\begin{bmatrix}(A\bm p_n)^{\top}&(A^{\top}\bar{\bm p}_n^{*})^{\top}\end{bmatrix}\begin{bmatrix}A^{\top}\bar{\bm p}_n^{*}\\A\bm p_n\end{bmatrix}\\
&=(A\bm p_n)^{\top}A^{\top}\bar{\bm p}_n^{*}+(A^{\top}\bar{\bm p}_n^{*})^{\top}A\bm p_n\\
&=(\bm p_n^{\top}A^{\top}A^{\top}\bar{\bm p}_n^{*})^{\top}+(A^{\mathrm H}\bm p_n^{*})^{\mathrm H}A\bm p_n\\
&=(A^{\mathrm H}\bm p_n^{*})^{\mathrm H}A\bm p_n+(A^{\mathrm H}\bm p_n^{*},A\bm p_n)\\
&=2(A^{\mathrm H}\bm p_n^{*},A\bm p_n),
\end{aligned}\]
\BT{we readily obtain the algorithm of the BiCR method for solving non-Hermitian linear systems.}{就容易得到求解非 Hermitian 线性方程组的 BiCR 算法.}
% block: bicr-algorithm page:120 type:ordinary-paragraph
\begin{algorithm}[H]\caption{\BT{The BiCR method}{BiCR 方法}}\label{alg:bicr}\begin{algorithmic}[1]
\Statex\textbf{Input:} $\bm x_0\in\mathbb C^N,\ \beta_{-1}=0,\ \bm p_{-1}=\bm p_{-1}^{*}=\bm0,\ \bm r_0=\bm b-A\bm x_0$
\Statex\textbf{Input:} \BT{Choose $\bm r_0^{*}\in\mathbb C^N$, e.g., $\bm r_0^{*}=\bm r_0$}{选择 $\bm r_0^{*}\in\mathbb C^N$, 例如 $\bm r_0^{*}=\bm r_0$}
\Statex\textbf{Output:} $\bm x_n$
\For{$n=0,1,2,\ldots$, \BT{until convergence}{直至收敛}}
\State $\bm p_n=\bm r_n+\beta_{n-1}\bm p_{n-1},\quad\bm p_n^{*}=\bm r_n^{*}+\bar\beta_{n-1}\bm p_{n-1}^{*}$
\State $A\bm p_n=A\bm r_n+\beta_{n-1}A\bm p_{n-1}$
\State $\alpha_n=\frac{(\bm r_n^{*},A\bm r_n)}{(A^{\mathrm H}\bm p_n^{*},A\bm p_n)}$
\State $\bm x_{n+1}=\bm x_n+\alpha_n\bm p_n$
\State $\bm r_{n+1}=\bm r_n-\alpha_nA\bm p_n,\quad\bm r_{n+1}^{*}=\bm r_n^{*}-\bar\alpha_nA^{\mathrm H}\bm p_n^{*}$
\State $\beta_n=\frac{(\bm r_{n+1}^{*},A\bm r_{n+1})}{(\bm r_n^{*},A\bm r_n)}$
\EndFor\end{algorithmic}\end{algorithm}
% block: bicr-biorthogonality-proposition page:120 type:proposition
\BT{Similar to Proposition\ \ref{prop:bicg-biorthogonality}, BiCR iterates hold the following properties:}{与 Proposition\ \ref{prop:bicg-biorthogonality} 类似, BiCR 迭代具有以下性质:}
\par\Needspace{3\baselineskip}
\begin{Proposition}\label{prop:bicr-biorthogonality}
\BT{Let $\bm p_n,\bm r_n$ be the vectors in the BiCR method. Then,}{设 $\bm p_n,\bm r_n$ 为 BiCR 方法中的向量. 则,}
\begin{enumerate}\item $(\bm r_i^{*},A\bm r_j)=0\quad\text{for }i\ne j$,\item $(A^{\mathrm H}\bm p_i^{*},A\bm p_j)=0\quad\text{for }i\ne j$.\end{enumerate}
\end{Proposition}
\BT{If the coefficient matrix $A$ is Hermitian or real symmetric, the choice $\bm r_0^{*}=\bm r_0$ reduces to the CR method that generates optimal approximate solutions in terms of residual 2-norms. Thus, if $A$ is close to Hermitian or real symmetric, it is expected that the BiCR method has smooth convergence behavior and generates near minimal residual solutions.}{若系数矩阵 $A$ 为 Hermitian 矩阵或实对称矩阵, 则选择 $\bm r_0^{*}=\bm r_0$ 就退化为 CR 方法, 后者产生残差 2-范数意义下的最优近似解. 因此, 若 $A$ 接近 Hermitian 矩阵或实对称矩阵, 可以预期 BiCR 方法具有平滑的收敛行为, 并产生接近最小残差的解.}
% block: bicr-convergence-comparison page:121 type:ordinary-paragraph
\BT{To see this, we consider nonsymmetric linear systems \eqref{eq:three-dimensional-difference-system} with parameters $N=20$, $a_1=a_2=a_3=1$, two different numbers $b_1=b_2=b_3=0.1$ (near symmetric matrix) and $b_1=b_2=b_3=10$ (far symmetric matrix), and $c=700$. The horizontal axis is the number of matrix--vector multiplications. The stopping criterion used for the BiCG method and the BiCR method is $\|\bm r_n\|/\|\bm b\|\le10^{-10}$.}{为说明这一点, 考虑非对称线性方程组 \eqref{eq:three-dimensional-difference-system}, 参数为 $N=20$, $a_1=a_2=a_3=1$, 两组不同的数值 $b_1=b_2=b_3=0.1$ (接近对称的矩阵) 和 $b_1=b_2=b_3=10$ (远离对称的矩阵), 以及 $c=700$. 横轴为矩阵向量乘法的次数. BiCG 方法和 BiCR 方法所用的停止准则为 $\|\bm r_n\|/\|\bm b\|\le10^{-10}$.}
\BT{The results for $b_1=b_2=b_3=0.1$ (near symmetric matrix) are shown in Fig.\ \ref{fig:bicg-bicr-near-symmetric}. We see from Fig.\ \ref{fig:bicg-bicr-near-symmetric} that the BiCR method shows much smoother convergence behavior than the BiCG method.}{当 $b_1=b_2=b_3=0.1$ (接近对称的矩阵) 时的结果见 Figure\ \ref{fig:bicg-bicr-near-symmetric}. 从 Figure\ \ref{fig:bicg-bicr-near-symmetric} 可见, BiCR 方法表现出的收敛行为比 BiCG 方法平滑得多.}
\BT{The results for $b_1=b_2=b_3=0.1$ (far symmetric matrix) are shown in Fig.\ \ref{fig:bicg-bicr-far-symmetric}. From Fig.\ \ref{fig:bicg-bicr-far-symmetric}, the BiCG method and the BiCR method have similar convergence behavior, and the BiCR method does not show smooth convergence behavior.}{当 $b_1=b_2=b_3=0.1$ (远离对称的矩阵) 时的结果见 Figure\ \ref{fig:bicg-bicr-far-symmetric}. 从 Figure\ \ref{fig:bicg-bicr-far-symmetric} 可见, BiCG 方法和 BiCR 方法具有相似的收敛行为, BiCR 方法未表现出平滑的收敛行为.}
% block: bicr-near-symmetric-figure page:121 type:figure
\begin{figure}[H]\centering\includegraphics[width=.87\linewidth]{figures/bicg-bicr-near-symmetric.png}
\caption{\BT{Convergence histories of BiCG and BiCR for $b_1=b_2=b_3=0.1$}{当 $b_1=b_2=b_3=0.1$ 时 BiCG 和 BiCR 的收敛历程}}\label{fig:bicg-bicr-near-symmetric}\end{figure}
% block: bicr-far-symmetric-figure page:122 type:figure
\begin{figure}[H]\centering\includegraphics[width=.87\linewidth]{figures/bicg-bicr-far-symmetric.png}
\caption{\BT{Convergence histories of BiCG and BiCR for $b_1=b_2=b_3=10$}{当 $b_1=b_2=b_3=10$ 时 BiCG 和 BiCR 的收敛历程}}\label{fig:bicg-bicr-far-symmetric}\end{figure}
% block: qmr-heading pages:121,122 type:subsection
\subsubsection{\BT{The Quasi-Minimal Residual (QMR) Method}{准最小残差 (QMR) 方法}}\label{subsec:qmr-method}
\BT{In this subsection, we describe the idea of the Quasi-Minimal Residual (QMR) method \BookCite{66}. The derivation process of the QMR method is almost the same as that of the QMR\_SYM method in Sect.\ \ref{subsec:qmr-sym}. The main difference is the use not of the complex symmetric Lanczos process but the bi-Lanczos process. Let $V_n$ be the bi-orthogonal basis of $\mathcal K_n(A,\bm r_0)$ via the bi-Lanczos process given in Algorithm\ \ref{alg:bi-lanczos}. Then, the QMR method finds an approximate solution $\bm x_n$ such that $\bm x_n$ lies in the affine space $\bm x_0+\mathcal K_n(A,\bm r_0)$, i.e.,}{本小节介绍准最小残差 (QMR) 方法 \BookCite{66} 的思想. QMR 方法的推导过程与 Section\ \ref{subsec:qmr-sym} 中 QMR\_SYM 方法的推导过程几乎相同. 主要区别是不采用复对称 Lanczos 过程, 而采用双 Lanczos 过程. 设 $V_n$ 为通过 Algorithm\ \ref{alg:bi-lanczos} 给出的双 Lanczos 过程得到的 $\mathcal K_n(A,\bm r_0)$ 的双正交基. 则 QMR 方法寻找近似解 $\bm x_n$, 使 $\bm x_n$ 位于仿射空间 $\bm x_0+\mathcal K_n(A,\bm r_0)$ 中, 即,}
\[\bm x_n=\bm x_0+V_n\bm y_n,\quad\bm y_n\in\mathbb C^n.\]
\BT{The corresponding residual vector is given by}{相应的残差向量为}
\[\bm r_n=\bm r_0-AV_n\bm y_n.\]
\BT{From the matrix form of the bi-Lanczos process \eqref{eq:bilanczos-matrix-relation}, it follows that}{由双 Lanczos 过程的矩阵形式 \eqref{eq:bilanczos-matrix-relation} 得}
\[\bm r_n=\bm r_0-V_{n+1}T_{n+1,n}\bm y_n=V_{n+1}(\beta\bm e_1-T_{n+1,n}\bm y_n),\]
\BT{where $\beta:=\|\bm r_0\|$. The above derivation process is very similar to that of the MINRES method; however, in this case, if we choose $\bm y_n$ such that the norm of the residual is minimized, then it may lead to a large amount of computational costs and memory requirement as described in Sect.\ \ref{subsec:qmr-sym}. Similar to the QMR\_SYM method, the QMR method chooses $\bm y_n$ such that}{其中 $\beta:=\|\bm r_0\|$. 上述推导过程与 MINRES 方法的推导过程很相似; 然而, 在此情形下, 若选择 $\bm y_n$ 使残差范数达到极小值, 则可能导致如 Section\ \ref{subsec:qmr-sym} 所述的大量计算开销和内存需求. 与 QMR\_SYM 方法类似, QMR 方法选择 $\bm y_n$ 使得}
\begin{equation}\label{eq:qmr-projected-minimization}\bm y_n:=\operatorname*{arg\,min}_{\bm y\in\mathbb C^n}\|\beta\bm e_1-T_{n+1,n}\bm y\|.\end{equation}
\BT{The minimization problem can be solved efficiently by Givens rotations. Following the solution of \eqref{eq:minres-projected-least-squares} and Sect.\ \ref{subsec:qmr-sym}, the QMR method is obtained and the algorithm is described in Algorithm\ \ref{alg:qmr}.}{这一极小化问题可由 Givens 旋转高效求解. 按照 \eqref{eq:minres-projected-least-squares} 的求解及 Section\ \ref{subsec:qmr-sym}, 即得到 QMR 方法, 其算法见 Algorithm\ \ref{alg:qmr}.}
% block: qmr-breakdown pages:122,123 type:ordinary-paragraph
\BT{Finally, one should note that Algorithm\ \ref{alg:qmr}, as well as the BiCG method and the BiCR method, has a possibility of breakdown (zero division). In finite precision arithmetic, such breakdown is very rare; however, near breakdown may occur, and this causes numerical instability. Hence, to avoid the (near) breakdown problem, the QMR method in \BookCite{66} uses the look-ahead Lanczos process, which was proposed by Taylor \BookCite{189} and Parlett et al.\ \BookCite{145}.}{最后应注意, Algorithm\ \ref{alg:qmr} 与 BiCG 方法及 BiCR 方法一样, 存在发生中断 (除以零) 的可能性. 在有限精度运算中, 这样的中断很罕见; 但可能发生近中断, 并导致数值不稳定. 因此, 为避免 (近) 中断问题, \BookCite{66} 中的 QMR 方法使用由 Taylor \BookCite{189} 和 Parlett 等 \BookCite{145} 提出的前瞻 Lanczos 过程.}
% block: qmr-algorithm page:123 type:ordinary-paragraph
\begin{algorithm}[H]\caption{\BT{The QMR method}{QMR 方法}}\label{alg:qmr}\begin{algorithmic}[1]
\Statex\textbf{Input:} $\bm x_0\in\mathbb C^N,\ \beta_0=0,\ \gamma_0=0,\ \bm v_0=\bm0,\ \bm w_0=\bm0,\ \bm r_0=\bm b-A\bm x_0$
\Statex\textbf{Output:} $\bm x_n$
\State $\bm g=(\|\bm r_0\|,0,\ldots,0)^{\top},\quad\bm v_1=\bm r_0/\|\bm r_0\|$
\For{$n=1,2,\ldots$}
\State (\BT{Bi-Lanczos process}{双 Lanczos 过程})
\State $\alpha_n=(\bm w_n,A\bm v_n),$
\State $\widetilde{\bm v}_{n+1}=A\bm v_n-\alpha_n\bm v_n-\beta_{n-1}\bm v_{n-1}$
\State $\widetilde{\bm w}_{n+1}=A^{\mathrm H}\bm w_n-\bar\alpha_n\bm w_n-\gamma_{n-1}\bm w_{n-1}$
\State $\gamma_n=\|\widetilde{\bm v}_{n+1}\|_2$
\State $\bm v_{n+1}=\widetilde{\bm v}_{n+1}/\gamma_n$
\State $\beta_n=(\widetilde{\bm w}_{n+1},\bm v_{n+1})$
\State $\bm w_{n+1}=\widetilde{\bm w}_{n+1}/\bar\beta_n$
\State (Givens \BT{rotations}{旋转})
\State \textbf{for} $i=\max\{1,n-2\},\ldots,n-1$ \textbf{do}
\State $\begin{bmatrix}t_{i,n}\\t_{i+1,n}\end{bmatrix}=\begin{bmatrix}c_i&s_i\\-\bar s_i&c_i\end{bmatrix}\begin{bmatrix}t_{i,n}\\t_{i+1,n}\end{bmatrix}$
\State \textbf{end for}
\State $c_n=\frac{|t_{n,n}|}{\sqrt{|t_{n,n}|^2+|t_{n+1,n}|^2}}$
\State $\bar s_n=\frac{t_{n+1,n}}{t_{n,n}}c_n$
\State $t_{n,n}=c_nt_{n,n}+s_nt_{n+1,n}$
\State $t_{n+1,n}=0$
\State $\begin{bmatrix}g_n\\g_{n+1}\end{bmatrix}=\begin{bmatrix}c_n&s_n\\-\bar s_n&c_n\end{bmatrix}\begin{bmatrix}g_n\\0\end{bmatrix}$
\State (\BT{Update}{更新} $\bm x_n$)
\State $\bm p_n=(\bm v_n-t_{n-2,n}\bm p_{n-2},$
\State $\qquad-t_{n-1,n}\bm p_{n-1})/t_{n,n}$
\State $\bm x_n=\bm x_{n-1}+g_n\bm p_n$
\State (\BT{Check convergence}{检查收敛})
\State $A\bm p_n=(A\bm v_n-t_{n-2,n}A\bm p_{n-2}$
\State $\qquad-t_{n-1,n}A\bm p_{n-1})/t_{n,n}$
\State $\bm r_n=\bm r_{n-1}-g_nA\bm p_n$
\State \BT{if $\|\bm r_n\|/\|\bm b\|\le\epsilon$, then stop}{若 $\|\bm r_n\|/\|\bm b\|\le\epsilon$, 则停止}
\EndFor\end{algorithmic}\end{algorithm}
% block: gmres-heading pages:123,124 type:subsection
\subsubsection{\BT{The Generalized Minimal Residual (GMRES) Method}{广义最小残差 (GMRES) 方法}}\label{subsec:gmres-method}
\BT{The derivation process of the GMRES method \BookCite{152} is closely related to the MINRES method in Sect.\ \ref{subsec:minres}. The GMRES method generates $\bm x_n$ that minimizes $\|\bm b-A\bm x_n\|$ over the affine space $\bm x_0+\mathcal K_n(A,\bm r_0)$. This can efficiently be achieved by using the Arnoldi process in Section\ \ref{subsec:arnoldi-process}. Now, we give the derivation process of the GMRES method.}{GMRES 方法 \BookCite{152} 的推导过程与 Section\ \ref{subsec:minres} 中的 MINRES 方法密切相关. GMRES 方法生成使 $\|\bm b-A\bm x_n\|$ 在仿射空间 $\bm x_0+\mathcal K_n(A,\bm r_0)$ 上达到极小值的 $\bm x_n$. 这可通过 Section\ \ref{subsec:arnoldi-process} 中的 Arnoldi 过程高效实现. 下面给出 GMRES 方法的推导过程.}
\BT{Let $V_n$ be an $N$-by-$n$ matrix whose columns are the orthonormal basis vectors of $\mathcal K_n(A,\bm r_0)$ by the Arnoldi process. Then, since $\bm x_n$ lies in the affine space $\bm x_0+\mathcal K_n(A,\bm r_0)$, we have}{设 $V_n$ 为 $N$ 行 $n$ 列的矩阵, 其列为 Arnoldi 过程生成的 $\mathcal K_n(A,\bm r_0)$ 的标准正交基向量. 于是, 由于 $\bm x_n$ 位于仿射空间 $\bm x_0+\mathcal K_n(A,\bm r_0)$ 中, 有}
\[\bm x_n=\bm x_0+V_n\bm y_n,\quad\bm y_n\in\mathbb C^n.\]
\BT{The corresponding residual vector is written as}{相应的残差向量写为}
\[\bm r_n=\bm r_0-AV_n\bm y_n.\]
\BT{From the matrix representation of the Arnoldi process in \eqref{eq:arnoldi-matrix-relation} it follows that}{由 \eqref{eq:arnoldi-matrix-relation} 中 Arnoldi 过程的矩阵表示得}
\begin{equation}\label{eq:gmres-projected-residual}\bm r_n=V_{n+1}(\beta\bm e_1-H_{n+1,n}\bm y_n),\end{equation}
\BT{where $\beta:=\|\bm r_0\|$. Since $V_{n+1}^{\mathrm H}V_{n+1}=I_{n+1}$, i.e., the identity matrix, the 2-norm of the residual vector is given by}{其中 $\beta:=\|\bm r_0\|$. 由于 $V_{n+1}^{\mathrm H}V_{n+1}=I_{n+1}$, 即单位矩阵, 残差向量的 2-范数为}
\[\|\bm r_n\|=\|V_{n+1}(\beta\bm e_1-H_{n+1,n}\bm y_n)\|=\|\beta\bm e_1-H_{n+1,n}\bm y_n\|.\]
\BT{Hence, the 2-norm of the residual vector can be minimized by choosing $\bm y_n$ as follows:}{因此, 按下式选择 $\bm y_n$ 可使残差向量的 2-范数达到极小值:}
\begin{equation}\label{eq:gmres-projected-minimization}\bm y_n:=\operatorname*{arg\,min}_{\bm y\in\mathbb C^n}\|\beta\bm e_1-H_{n+1,n}\bm y\|.\end{equation}
\BT{The vector $\bm y_n$ can be obtained by using Givens rotations. Following the solution of \eqref{eq:minres-projected-least-squares}, the GMRES method is obtained. Now, we describe the algorithm of GMRES in Algorithm\ \ref{alg:gmres}, which is based on the Arnoldi process of modified Gram--Schmidt type (Algorithm\ \ref{alg:arnoldi-modified}).}{向量 $\bm y_n$ 可通过 Givens 旋转得到. 按照 \eqref{eq:minres-projected-least-squares} 的求解就得到 GMRES 方法. 下面在 Algorithm\ \ref{alg:gmres} 中描述 GMRES 算法, 它基于改进 Gram--Schmidt 型 Arnoldi 过程 (Algorithm\ \ref{alg:arnoldi-modified}).}
% block: gmres-algorithm page:124 type:ordinary-paragraph
\begin{algorithm}[H]\caption{\BT{The GMRES method of modified Gram-Schmidt type}{改进 Gram-Schmidt 型 GMRES 方法}}\label{alg:gmres}\begin{algorithmic}[1]
\Statex\textbf{Input:} $\bm x_0\in\mathbb C^N,\ \bm r_0=\bm b-A\bm x_0$
\Statex\textbf{Output:} $\bm x_n$
\State $\bm g=(\|\bm r_0\|,0,\ldots,0)^{\top},\quad\bm v_1=\bm r_0/\|\bm r_0\|$
\For{$n=1,2,\ldots$}
\State (Arnoldi \BT{process}{过程})
\State $\bm t=A\bm v_n$
\For{$i=1,2,\ldots,n$}
\State $h_{i,n}=(\bm v_i,\bm t)$
\State $\bm t=\bm t-h_{i,n}\bm v_i$
\EndFor
\State $h_{n+1,n}=\|\bm t\|$
\State $\bm v_{n+1}=\bm t/h_{n+1,n}$
\State (Givens \BT{rotations}{旋转})
\For{$i=1,2,\ldots,n-1$}
\State $\begin{bmatrix}h_{i,n}\\h_{i+1,n}\end{bmatrix}=\begin{bmatrix}c_i&s_i\\-\bar s_i&c_i\end{bmatrix}\begin{bmatrix}h_{i,n}\\h_{i+1,n}\end{bmatrix}$
\EndFor
\State $c_n=\frac{|h_{n,n}|}{\sqrt{|h_{n,n}|^2+|h_{n+1,n}|^2}}$
\State $\bar s_n=\frac{h_{n+1,n}}{h_{n,n}}c_n$
\State $h_{n,n}=c_nh_{n,n}+s_nh_{n+1,n}$
\State $h_{n+1,n}=0$
\State $\begin{bmatrix}g_n\\g_{n+1}\end{bmatrix}=\begin{bmatrix}c_n&s_n\\-\bar s_n&c_n\end{bmatrix}\begin{bmatrix}g_n\\0\end{bmatrix}$
\State (\BT{Check convergence}{检查收敛})
\If{$|g_{n+1}|/\|\bm b\|\le\epsilon$}
\State $\bm x_n=\bm x_0+V_nH_n^{-1}\bm g$
\EndIf\EndFor\end{algorithmic}\end{algorithm}
% block: gmres-back-substitution page:124 type:ordinary-paragraph
\BT{Note that computing $\bm c:=H_n^{-1}\bm g$ in line 22 of Algorithm\ \ref{alg:gmres} corresponds to solving linear systems of the form $H_n\bm c=\bm g$. After Givens rotations, $H_n$ becomes not a Hessenberg but an upper triangular matrix. Thus the linear systems can easily be solved by the back substitution as described in \eqref{eq:lu-triangular-solves}.}{注意, 在 Algorithm\ \ref{alg:gmres} 第 22 行计算 $\bm c:=H_n^{-1}\bm g$, 对应于求解形如 $H_n\bm c=\bm g$ 的线性方程组. 经过 Givens 旋转后, $H_n$ 不再是 Hessenberg 矩阵, 而变为上三角矩阵. 因此该线性方程组可通过 \eqref{eq:lu-triangular-solves} 所述的回代容易地求解.}
\BT{In what follows, the convergence analysis of the GMRES method is described.}{下面介绍 GMRES 方法的收敛性分析.}
% block: gmres-convergence-proposition pages:124,125 type:proposition
\par\Needspace{3\baselineskip}
\begin{Proposition}\label{prop:gmres-diagonalizable-bound}
\BT{Let $A$ be an $N$-by-$N$ diagonalizable matrix, i.e., $A=X\Lambda X^{-1}$ with $\Lambda=\operatorname{diag}\{\lambda_1,\ldots,\lambda_N\}$, and let $\bm r_n$ be the $m$-step GMRES residual vector. Then}{设 $A$ 为 $N$ 行 $N$ 列的可对角化矩阵, 即 $A=X\Lambda X^{-1}$, 其中 $\Lambda=\operatorname{diag}\{\lambda_1,\ldots,\lambda_N\}$, 设 $\bm r_n$ 为 $m$ 步 GMRES 残差向量. 则}
\[\|\bm r_n\|\le\kappa(X)\epsilon_n\|\bm r_0\|,\]
\BT{where $\kappa(X)=\|X\|\|X^{-1}\|$ and}{其中 $\kappa(X)=\|X\|\|X^{-1}\|$, 且}
\[\epsilon^{(n)}=\min_{p\in\mathbb P_n,p(0)=1}\max_{i=1,\ldots,N}|p(\lambda_i)|.\]
\BT{Here, $\mathbb P_m$ denotes the set of all polynomials of degree at most $m$.}{这里 $\mathbb P_m$ 表示所有次数至多为 $m$ 的多项式构成的集合.}
\end{Proposition}
% block: gmres-convergence-proof page:125 type:proof
\begin{Proof}
\BT{The GMRES method finds an approximate solution $\bm x_n$ such that}{GMRES 方法寻找近似解 $\bm x_n$ 使得}
\[\min_{\bm x_n\in\bm x_0+\mathcal K_n(A,\bm r_0)}\|\bm r_n\|.\]
\BT{Since $\bm r_n\in\mathcal K_{n+1}(A,\bm r_0)$, the GMRES residual vector can be written in the matrix polynomial form $\bm r_n=(I+c_1A+c_2A^2+\cdots+c_nA^n)\bm r_0$, so that $c_i$'s are determined by minimizing $\|\bm r_n\|$, or equivalently $\bm r_n=p_n^{(\mathrm{opt})}(A)\bm r_0$ with $p_n^{(\mathrm{opt})}(0)=1$, and $p_n^{(\mathrm{opt})}$ is determined by minimizing $\|\bm r_n\|$, i.e.,}{由于 $\bm r_n\in\mathcal K_{n+1}(A,\bm r_0)$, GMRES 残差向量可写为矩阵多项式形式 $\bm r_n=(I+c_1A+c_2A^2+\cdots+c_nA^n)\bm r_0$, 其中 $c_i$ 通过极小化 $\|\bm r_n\|$ 确定; 等价地, $\bm r_n=p_n^{(\mathrm{opt})}(A)\bm r_0$, 其中 $p_n^{(\mathrm{opt})}(0)=1$, 且 $p_n^{(\mathrm{opt})}$ 通过极小化 $\|\bm r_n\|$ 确定, 即,}
\[\min_{p\in\mathbb P_n,p(0)=1}\|p_n(A)\bm r_0\|.\]
\BT{Therefore, for any $p\in\mathbb P_n$ and $p(0)=1$, the residual 2-norm $\|\bm r_n\|$ can be bounded by}{因此, 对任意 $p\in\mathbb P_n$ 且 $p(0)=1$, 残差 2-范数 $\|\bm r_n\|$ 有如下界}
\[\begin{aligned}\|\bm r_n\|&=\|p_n^{(\mathrm{opt})}(A)\bm r_0\|\\&\le\|p_n(A)\bm r_0\|\\&=\|Xp_n(\Lambda)X^{-1}\bm r_0\|\\&\le\|X\|\|X^{-1}\|\|p_n(\Lambda)\|\|\bm r_0\|\\&=\kappa(X)\max_{i=1,\ldots,N}|p_n(\lambda_i)|\|\bm r_0\|\end{aligned}\]
\BT{Since the above inequality holds for any $p\in\mathbb P_n$ and $p(0)=1$, we have}{由于上述不等式对任意 $p\in\mathbb P_n$ 且 $p(0)=1$ 都成立, 有}
\[\begin{aligned}\|\bm r_n\|&\le\kappa(X)\left(\min_{p\in\mathbb P_n,p(0)=1}\max_{i=1,\ldots,N}|p(\lambda_i)|\right)\|\bm r_0\|\\&=\kappa(X)\epsilon_n\|\bm r_0\|,\end{aligned}\]
\BT{which concludes the proof.}{证明完毕.}
\par\nopagebreak[4]
\end{Proof}
% block: gmres-eigenvalue-discussion page:125 type:ordinary-paragraph
\BT{From Proposition\ \ref{prop:gmres-diagonalizable-bound}, if the condition number of $X$ is small, the value of $\epsilon_n$ depends highly on the speed of convergence. Further, if the condition number of $X$ is small and eigenvalues are well clustered, the GMRES method shows good convergence behavior. If matrix $A$ is symmetric, then the condition number of $X$ is 1. Therefore the convergence of the GMRES method (or the MINRES method) depends only on the distribution of eigenvalues. By using an upper bound of $\epsilon_n$, the convergence behavior can be estimated without using all the eigenvalues. For the details, see an excellent book by Saad \BookCite{151}, pp.206--207.}{由 Proposition\ \ref{prop:gmres-diagonalizable-bound}, 若 $X$ 的条件数较小, 则 $\epsilon_n$ 的值高度依赖于收敛速度. 进一步, 若 $X$ 的条件数较小且特征值良好聚集, 则 GMRES 方法表现出良好的收敛行为. 若矩阵 $A$ 对称, 则 $X$ 的条件数为 1. 因此 GMRES 方法 (或 MINRES 方法) 的收敛性仅依赖于特征值的分布. 使用 $\epsilon_n$ 的一个上界, 即可不使用全部特征值而估计收敛行为. 详情见 Saad 的优秀著作 \BookCite{151}, 第 206--207 页.}
% block: restarted-gmres-motivation page:126 type:ordinary-paragraph
\BT{Nice features of the GMRES method are that the GMRES method never suffers from breakdown, and the residual 2-norm decreases monotonically and has the optimality regarding the residual 2-norm, unlike Krylov subspace methods based on the BiCG method and the BiCR method. On the other hand, a drawback of the GMRES method is that the computational costs and memory requirement of the GMRES method increase linearly with the number of iterations. For example, if the number of iterations is 10000, then an $N\times10000$ matrix needs to be stored and the size of the least-squares problem in \eqref{eq:gmres-projected-minimization} is $10001\times10000$. To overcome the drawback, a restarted version of the GMRES method is proposed, which is described in Algorithm\ \ref{alg:gmres-restarted}.}{GMRES 方法的良好特点在于, 与基于 BiCG 方法和 BiCR 方法的 Krylov 子空间方法不同, GMRES 方法从不发生中断, 残差 2-范数单调下降, 且具有残差 2-范数意义下的最优性. 另一方面, GMRES 方法的一个缺点是其计算开销和内存需求随迭代次数线性增加. 例如, 若迭代次数为 10000, 则需存储一个 $N\times10000$ 矩阵, 而 \eqref{eq:gmres-projected-minimization} 中最小二乘问题的规模为 $10001\times10000$. 为克服这一缺点, 提出了重启的 GMRES 方法, 见 Algorithm\ \ref{alg:gmres-restarted}.}
\BT{The idea of the restarted GMRES method denoted by GMRES($m$) is that we run the GMRES method until the prescribed maximum iteration number $m$ is reached, and then we restart the GMRES method using the initial guess as the approximate solution produced by the previous GMRES method at the $m$th iteration step.}{记为 GMRES($m$) 的重启 GMRES 方法的思想是, 运行 GMRES 方法直至达到预定最大迭代次数 $m$, 然后以此前 GMRES 方法在第 $m$ 次迭代产生的近似解作为初始猜测, 重新启动 GMRES 方法.}
% block: restarted-gmres-algorithm page:126 type:ordinary-paragraph
\begin{algorithm}[H]\caption{\BT{The GMRES($m$) method}{GMRES($m$) 方法}}\label{alg:gmres-restarted}\begin{algorithmic}[1]
\Statex\textbf{Input:} \BT{Initial guess $\bm x_0\in\mathbb C^N$ and restart number $m$}{初始猜测 $\bm x_0\in\mathbb C^N$ 和重启数 $m$}
\Statex\textbf{Output:} $\bm x_m$
\State\BT{Run $m$ iterations of the GMRES method (Algorithm\ \ref{alg:gmres}) with $\bm x_0$ and produce $\bm x_m$.}{以 $\bm x_0$ 运行 GMRES 方法 (Algorithm\ \ref{alg:gmres}) 的 $m$ 次迭代, 得到 $\bm x_m$.}
\While{\BT{convergence}{收敛}}
\State\BT{Set $\bm x_0=\bm x_m$.}{置 $\bm x_0=\bm x_m$.}
\State\BT{Run $m$ iterations of the GMRES method (Algorithm\ \ref{alg:gmres}) with $\bm x_0$ and produce $\bm x_m$.}{以 $\bm x_0$ 运行 GMRES 方法 (Algorithm\ \ref{alg:gmres}) 的 $m$ 次迭代, 得到 $\bm x_m$.}
\EndWhile\end{algorithmic}\end{algorithm}
% block: gmres-restart-frequency page:126 type:ordinary-paragraph
\BT{Consider the restart frequency $m$ of the GMRES($m$) method as fixed. This means that we apply the $m$-step GMRES method to the linear systems repeatedly until convergence. On the other hand, an unfixed restart frequency means that we run the GMRES($m_1$) and then run the GMRES($m_2$) method, the GMRES($m_3$) method, and so on, until convergence. If $m_1=m_2=\cdots=m$ then this corresponds to the fixed restart frequency.}{考虑 GMRES($m$) 方法的重启频率 $m$ 固定的情形. 这意味着反复对线性方程组应用 $m$ 步 GMRES 方法, 直至收敛. 另一方面, 非固定的重启频率是指先运行 GMRES($m_1$), 然后运行 GMRES($m_2$) 方法, GMRES($m_3$) 方法, 如此继续, 直至收敛. 若 $m_1=m_2=\cdots=m$, 则对应于固定的重启频率.}
\BT{Efficient ways for determining $m_k$ for $k=1,2,\ldots$ are studied in \BookCite{117} and by some authors, e.g., \BookCite{14}, \BookCite{137}, \BookCite{176}, \BookCite{211}.}{文献 \BookCite{117} 及一些作者研究了对 $k=1,2,\ldots$ 确定 $m_k$ 的高效方法, 例如 \BookCite{14}, \BookCite{137}, \BookCite{176}, \BookCite{211}.}
% block: lgmres-augmentation pages:126,127 type:ordinary-paragraph
\BT{Another important development of the GMRES method is augmentation. There are two approaches to the augmentation: one is to use previous approximate solutions, e.g., the LGMRES method \BookCite{15}; another is to use approximate eigenvectors corresponding to small eigenvalues in magnitude, e.g., the GMRES-E method \BookCite{134}. The basic ideas of the LGMRES method and the GMRES-E method are described next.}{GMRES 方法的另一个重要发展是增广. 增广有两种途径: 一种利用此前的近似解, 例如 LGMRES 方法 \BookCite{15}; 另一种利用对应于模较小特征值的近似特征向量, 例如 GMRES-E 方法 \BookCite{134}. 下面介绍 LGMRES 方法和 GMRES-E 方法的基本思想.}
\BT{The group of $m$ iteration steps (e.g., lines 3 and 4 in Algorithm\ \ref{alg:gmres-restarted}) is called a cycle, and the approximate solution of the LGMRES method after the $j$th cycle is denoted by $\bm x_m^{(j)}$. Let $\bm z^{(k)}:=\bm x_m^{(j-1)}-\bm x_m^{(j-2)}$. Then for a given $l$, the LGMRES method finds an approximate solution in the following way:}{一组 $m$ 次迭代 (例如 Algorithm\ \ref{alg:gmres-restarted} 中第 3 和第 4 行) 称为一个循环, 将 LGMRES 方法在第 $j$ 个循环后的近似解记为 $\bm x_m^{(j)}$. 令 $\bm z^{(k)}:=\bm x_m^{(j-1)}-\bm x_m^{(j-2)}$. 则对给定的 $l$, LGMRES 方法按下列方式寻找近似解:}
\[\bm x_m^{(j)}=\bm x_m^{(j-1)}+q_{m-1}^{(j)}(A)\bm r_m^{(j-1)}+\sum_{k=j-l}^{j-1}\alpha_{j,k}(A)\bm z^{(k)},\]
\BT{where $q_{m-1}^{(j)}(A)$ is a polynomial of degree $m-1$, and $\alpha_{j,k}(A)$'s are polynomials of degree $k$. These polynomials are determined by minimizing the residual 2-norm of $\bm x_m^{(j)}$. For the related algorithms, see, e.g., \BookCite{13}, \BookCite{103}, \BookCite{104}, \BookCite{106}.}{其中 $q_{m-1}^{(j)}(A)$ 为 $m-1$ 次多项式, $\alpha_{j,k}(A)$ 为 $k$ 次多项式. 这些多项式通过极小化 $\bm x_m^{(j)}$ 的残差 2-范数确定. 相关算法例如见 \BookCite{13}, \BookCite{103}, \BookCite{104}, \BookCite{106}.}
% block: gmres-e-augmentation page:127 type:ordinary-paragraph
\BT{The approximate solution of the GMRES-E method after the $j$th cycle is denoted by $\bm x_m^{(j)}$, and $\bm p_k^{(j)}$ is an approximate eigenvector (Harmonic Ritz vectors) that corresponds to the $k$th smallest approximate eigenvalue (Harmonic Ritz value) in magnitude. Then, for a given $p$ the GMRES-E method finds an approximate solution in the following way:}{将 GMRES-E 方法在第 $j$ 个循环后的近似解记为 $\bm x_m^{(j)}$, $\bm p_k^{(j)}$ 为对应于模第 $k$ 小的近似特征值 (调和 Ritz 值) 的近似特征向量 (调和 Ritz 向量). 则对给定的 $p$, GMRES-E 方法按下列方式寻找近似解:}
\[\bm x_m^{(j)}=\bm x_m^{(j-1)}+q_{m-1}^{(j)}(A)\bm r_m^{(j-1)}+\sum_{k=1}^{p}\beta_{j,k}(A)\bm p_k^{(j-1)},\]
\BT{where $q_{m-1}^{(j)}(A)$ is a polynomial of degree $m-1$, and $\beta_{j,k}(A)$'s are polynomials of degree $k$. These polynomials are determined by minimizing the residual 2-norm of $\bm x_m^{(j)}$. For the related algorithms, see, e.g., \BookCite{11}, \BookCite{135}. By combining the LGMRES method and the GMRES-E method, the LGMRES-E method using a switching controller is proposed in \BookCite{31}. Roughly speaking, the switching rule is to choose the LGMRES method while good convergence appears and to choose the GMRES-E method when slow convergence is found.}{其中 $q_{m-1}^{(j)}(A)$ 为 $m-1$ 次多项式, $\beta_{j,k}(A)$ 为 $k$ 次多项式. 这些多项式通过极小化 $\bm x_m^{(j)}$ 的残差 2-范数确定. 相关算法例如见 \BookCite{11}, \BookCite{135}. 文献 \BookCite{31} 将 LGMRES 方法和 GMRES-E 方法结合, 提出了使用切换控制器的 LGMRES-E 方法. 粗略地说, 切换规则是在收敛良好时选择 LGMRES 方法, 在发现收敛缓慢时选择 GMRES-E 方法.}
\BT{The augmentation is not only for the GMRES method. In fact, frameworks of augmentation and deflation for Krylov subspace methods including the BiCG method, the BiCR method, and the other important Krylov space methods are developed in \BookCite{74}, \BookCite{89}, \BookCite{90}.}{增广并不只用于 GMRES 方法. 事实上, 文献 \BookCite{74}, \BookCite{89}, \BookCite{90} 建立了包括 BiCG 方法, BiCR 方法及其他重要 Krylov 空间方法在内的 Krylov 子空间方法的增广与消去框架.}
% block: gmres-numerical-example pages:127,128 type:ordinary-paragraph
\BT{Finally, we provide the result of a numerical example for the GMRES method and the GMRES($m$) method. Consider nonsymmetric linear systems \eqref{eq:three-dimensional-difference-system} with parameters $N=20$, $a_1=a_2=a_3=1$, $b_1=b_2=b_3=1$, and $c=1$. The result is shown in Fig.\ \ref{fig:gmres-restarts}, where the stopping criterion used for the GMRES method and the GMRES($m$) method is $\|\bm r_n\|/\|\bm b\|\le10^{-10}$, and the horizontal axis is the number of matrix--vector multiplications.}{最后给出 GMRES 方法和 GMRES($m$) 方法的一个数值例子的结果. 考虑非对称线性方程组 \eqref{eq:three-dimensional-difference-system}, 参数为 $N=20$, $a_1=a_2=a_3=1$, $b_1=b_2=b_3=1$, 以及 $c=1$. 结果见 Figure\ \ref{fig:gmres-restarts}, 其中 GMRES 方法和 GMRES($m$) 方法所用的停止准则为 $\|\bm r_n\|/\|\bm b\|\le10^{-10}$, 横轴为矩阵向量乘法的次数.}
\BT{We see from Fig.\ \ref{fig:gmres-restarts} that the residual 2-norms of the GMRES method, the GMRES(10) method, and the GMRES(20) method decrease monotonically, which are theoretically guaranteed. In terms of the number of matrix--vector multiplications, the GMRES method performs better than the restarted GMRES method (GMRES(10), GMRES(20)). As for the GMRES method, the required number of matrix--vector multiplications is the smallest of all, since it finds the best approximate solution over the Krylov subspace. On the other hand, as mentioned before, the computational cost of the GMRES method grows linearly as the number of iterations increases. This means that matrix--vector multiplications may not become the most time-consuming parts when the required number of iterations is very large. In terms of CPU time, the GMRES method required 0.39 sec., GMRES(10) 0.31 sec., GMRES(20) 0.26 sec.}{由 Figure\ \ref{fig:gmres-restarts} 可见, GMRES 方法, GMRES(10) 方法和 GMRES(20) 方法的残差 2-范数单调下降, 这有理论保证. 从矩阵向量乘法次数看, GMRES 方法优于重启 GMRES 方法 (GMRES(10), GMRES(20)). 对 GMRES 方法而言, 由于它在 Krylov 子空间上寻找最佳近似解, 所需的矩阵向量乘法次数最少. 另一方面, 如前所述, GMRES 方法的计算开销随迭代次数线性增加. 这意味着, 当所需迭代次数很大时, 矩阵向量乘法可能不再是最耗时的部分. 从 CPU 时间看, GMRES 方法需要 0.39 秒, GMRES(10) 需要 0.31 秒, GMRES(20) 需要 0.26 秒.}\footnote{\BT{Codes: GNU Octave version 6.2.0, OS: macOS Moterey version 12.2.1, CPU: Apple M1 pro.}{代码: GNU Octave 6.2.0 版, 操作系统: macOS Moterey 12.2.1 版, CPU: Apple M1 pro.}}
\BT{Therefore, with respect to CPU time, the GMRES(20) method was the best of all, which balanced the increasing computational costs and the required number of iterations.}{因此, 就 CPU 时间而言, GMRES(20) 方法在所有方法中最好, 它平衡了增加的计算开销与所需的迭代次数.}
% block: gmres-restart-figure page:128 type:figure
\begin{figure}[H]\centering\includegraphics[width=.87\linewidth]{figures/gmres-restarts.png}
\caption{\BT{Convergence histories of GMRES, GMRES(10), and GMRES(20)}{GMRES, GMRES(10) 和 GMRES(20) 的收敛历程}}\label{fig:gmres-restarts}\end{figure}
% block: gcr-heading pages:128,129 type:subsection
\subsubsection{\BT{The Generalized Conjugate Residual (GCR) Method}{广义共轭残差 (GCR) 方法}}\label{subsec:gcr-method}
\BT{The Generalized Conjugate Residual (GCR) method \BookCite{54} is another Krylov solver based on the Arnoldi process. In exact precision arithmetic, the GCR method produces the same approximate solution as the GMRES method. The algorithm is obtained by the following Arnoldi process with the weight $A^{\mathrm H}A$ and $\bm r_0=\bm b-A\bm x_0$:}{广义共轭残差 (GCR) 方法 \BookCite{54} 是另一种基于 Arnoldi 过程的 Krylov 求解器. 在精确运算中, GCR 方法产生与 GMRES 方法相同的近似解. 该算法由下列以 $A^{\mathrm H}A$ 为权且 $\bm r_0=\bm b-A\bm x_0$ 的 Arnoldi 过程得到:}
\BT{set}{置} $\bm v_1=\dfrac{\bm r_0}{\|\bm r_0\|_{A^{\mathrm H}A}}$,\par
\BT{for $n=1,2,\ldots$, do:}{对 $n=1,2,\ldots$, 执行:}
\[h_{i,n}=(\bm v_i,A\bm v_n)_{A^{\mathrm H}A},\quad i=1,2,\ldots,n,\]
\begin{equation}\label{eq:gcr-weighted-arnoldi-remainder}\widetilde{\bm v}_{n+1}=A\bm v_n-\sum_{i=1}^nh_{i,n}\bm v_i,\end{equation}
\[h_{n+1,n}=\|\widetilde{\bm v}_{n+1}\|_{A^{\mathrm H}A},\]
\begin{equation}\label{eq:gcr-weighted-arnoldi-normalization}\bm v_{n+1}=\frac{\widetilde{\bm v}_{n+1}}{h_{n+1,n}}.\end{equation}
\BT{end}{结束}
\BT{The process generates $A^{\mathrm H}A$-orthonormal vectors, i.e.,}{该过程生成 $A^{\mathrm H}A$-标准正交向量, 即,}
\[(\bm v_i,A^{\mathrm H}A\bm v_j)=(A\bm v_i,A\bm v_j)=\delta_{i,j}.\]
% block: gcr-residual-expansion page:129 type:ordinary-paragraph
\BT{Let $V_n$ be a matrix whose columns are $\bm v_1,\bm v_2,\ldots,\bm v_n$. Then, the GCR method generates}{设 $V_n$ 为以 $\bm v_1,\bm v_2,\ldots,\bm v_n$ 为列的矩阵. 则 GCR 方法生成}
\[\bm x_n=\bm x_0+V_n\bm y_n\]
\BT{such that the residual 2-norm $\|\bm r_n\|$ is minimized. The corresponding residual vector can be written as}{使残差 2-范数 $\|\bm r_n\|$ 达到极小值. 相应的残差向量可写为}
\begin{equation}\label{eq:gcr-residual-expansion}\bm r_n=\bm r_0-AV_n\bm y_n=\bm r_0-\sum_{i=1}^ny_iA\bm v_i.\end{equation}
\BT{From the above relation, we have}{由上述关系有}
\begin{equation}\label{eq:gcr-residual-basis-update}\bm r_n=\bm r_{n-1}-y_nA\bm v_n.\end{equation}
\BT{Unknown parameters $y_i$ ($i=1,2,\ldots,n$) are determined by minimizing the residual 2-norm, i.e., from \eqref{eq:gcr-residual-expansion} this can be achieved by}{未知参数 $y_i$ ($i=1,2,\ldots,n$) 通过极小化残差 2-范数确定, 即由 \eqref{eq:gcr-residual-expansion}, 这可通过下式实现}
\[y_i=(A\bm v_i,\bm r_n),\quad\text{for }i=1,2,\ldots,n.\]
\BT{Since $A\bm v_i$ is an orthonormal vector, we have}{由于 $A\bm v_i$ 为标准正交向量, 有}
\[y_i=(A\bm v_i,\bm r_0),\quad\text{for }i=1,2,\ldots,n.\]
\BT{Here, we define the vector $\bm p_n:=-y_n\widetilde{\bm v}_{n+1}$. Then, it follows from \eqref{eq:gcr-weighted-arnoldi-remainder} that}{这里定义向量 $\bm p_n:=-y_n\widetilde{\bm v}_{n+1}$. 则由 \eqref{eq:gcr-weighted-arnoldi-remainder} 得}
% block: gcr-direction-projection page:130 type:ordinary-paragraph
\begin{equation}\label{eq:gcr-direction-from-arnoldi}\bm p_n=-y_nA\bm v_n-\sum_{i=1}^n(\bm v_i,-y_nA\bm v_n)_{A^{\mathrm H}A}\bm v_i.\end{equation}
\BT{From \eqref{eq:gcr-residual-basis-update} we have $-y_nA\bm v_n=\bm r_n-\bm r_{n-1}$. Substituting this relation into \eqref{eq:gcr-direction-from-arnoldi} yields}{由 \eqref{eq:gcr-residual-basis-update} 有 $-y_nA\bm v_n=\bm r_n-\bm r_{n-1}$. 将此关系代入 \eqref{eq:gcr-direction-from-arnoldi}, 得}
\begin{equation}\label{eq:gcr-direction-residual-difference}\bm p_n=\bm r_n-\bm r_{n-1}-\sum_{i=1}^n(\bm v_i,\bm r_n-\bm r_{n-1})_{A^{\mathrm H}A}\bm v_i\end{equation}
\[=\bm r_n-\sum_{i=1}^n(\bm v_i,\bm r_n)_{A^{\mathrm H}A}\bm v_i-\left(\bm r_{n-1}-\sum_{i=1}^n(\bm v_i,\bm r_{n-1})_{A^{\mathrm H}A}\bm v_i\right).\]
\BT{Since $\bm r_{n-1}$ lies in $\mathcal K_n(A,\bm r_0)$, it can be expanded as $\bm r_{n-1}=\sum_{i=1}^nc_i\bm v_i$. Then,}{由于 $\bm r_{n-1}$ 位于 $\mathcal K_n(A,\bm r_0)$ 中, 可展开为 $\bm r_{n-1}=\sum_{i=1}^nc_i\bm v_i$. 于是,}
\[(\bm v_j,\bm r_{n-1})_{A^{\mathrm H}A}=\sum_{i=1}^nc_i(\bm v_j,\bm v_i)=c_i\]
\BT{because $(\bm v_j,\bm v_i)=\delta_{j,i}$. Thus we have}{因为 $(\bm v_j,\bm v_i)=\delta_{j,i}$. 因而有}
\begin{equation}\label{eq:gcr-previous-residual-projection}\bm r_{n-1}=\sum_{i=1}^n(\bm v_i,\bm r_{n-1})_{A^{\mathrm H}A}\bm v_i.\end{equation}
\BT{From \eqref{eq:gcr-direction-residual-difference} and \eqref{eq:gcr-previous-residual-projection}, the vector $\bm p_n$ can be written as}{由 \eqref{eq:gcr-direction-residual-difference} 和 \eqref{eq:gcr-previous-residual-projection}, 向量 $\bm p_n$ 可写为}
\begin{equation}\label{eq:gcr-direction-weighted-projection}\bm p_n=\bm r_n-\sum_{i=1}^n(\bm v_i,\bm r_n)_{A^{\mathrm H}A}\bm v_i.\end{equation}
% block: gcr-normalized-direction page:130 type:ordinary-paragraph
\BT{On the other hand, $\bm p_n$ and $\bm v_{n+1}$ are related by}{另一方面, $\bm p_n$ 和 $\bm v_{n+1}$ 的关系为}
\[\bm p_n=-\frac{y_n}{|y_n|}\|A\bm p_n\|\bm v_{n+1},\]
\BT{because it follows from the definition of $\bm p_n$ and \eqref{eq:gcr-weighted-arnoldi-normalization} that}{因为由 $\bm p_n$ 的定义及 \eqref{eq:gcr-weighted-arnoldi-normalization} 得}
\[\begin{aligned}\bm p_n&=-y_n\widetilde{\bm v}_{n+1}\\&=-y_n\|\widetilde{\bm v}_{n+1}\|_{A^{\mathrm H}A}\bm v_{n+1}\\&=-y_n\sqrt{(A\widetilde{\bm v}_{n+1},A\widetilde{\bm v}_{n+1})}\,\bm v_{n+1}\\&=-y_n\sqrt{\left(-\frac1{y_n}A\bm p_n,-\frac1{y_n}A\bm p_n\right)}\,\bm v_{n+1}\\&=-\frac{y_n}{|y_n|}\sqrt{(A\bm p_n,A\bm p_n)}\,\bm v_{n+1}.\end{aligned}\]
% block: gcr-coefficients page:131 type:ordinary-paragraph
\BT{Hence, using $\bm v_n=-(|y_{n-1}|/y_{n-1})\times\bm p_{n-1}/\sqrt{(A\bm p_{n-1},A\bm p_{n-1})}$, we have the following relations:}{因此, 使用 $\bm v_n=-(|y_{n-1}|/y_{n-1})\times\bm p_{n-1}/\sqrt{(A\bm p_{n-1},A\bm p_{n-1})}$, 有下列关系:}
\[\begin{aligned}y_nA\bm v_n&=(A\bm v_n,\bm r_0)A\bm v_n=\frac{(A\bm p_{n-1},\bm r_0)}{(A\bm p_{n-1},A\bm p_{n-1})}A\bm p_{n-1},\\(\bm v_i,\bm r_n)_{A^{\mathrm H}A}\bm v_i&=\frac{(A\bm p_{i-1},A\bm r_n)}{(A\bm p_{i-1},A\bm p_{i-1})}\bm p_{i-1}.\end{aligned}\]
\BT{Therefore, substituting the above results into \eqref{eq:gcr-residual-basis-update} and \eqref{eq:gcr-direction-weighted-projection}, we have}{因此, 将上述结果代入 \eqref{eq:gcr-residual-basis-update} 和 \eqref{eq:gcr-direction-weighted-projection}, 得}
\[\begin{aligned}\bm r_n&=\bm r_{n-1}-\frac{(A\bm p_{n-1},\bm r_0)}{(A\bm p_{n-1},A\bm p_{n-1})}A\bm p_{n-1},\\\bm p_n&=\bm r_n-\sum_{i=1}^n\frac{(A\bm p_{i-1},A\bm r_n)}{(A\bm p_{i-1},A\bm p_{i-1})}\bm p_{i-1}.\end{aligned}\]
\BT{Now, from the definition of $\bm p_n$, we have $\bm p_{n-1}=c\bm v_n$ for a constant $c$. Thus $\bm p_{n-1}$ is $A^{\mathrm H}A$-orthogonal to $\bm v_1,\ldots,\bm v_{n-1}$. From \eqref{eq:gcr-residual-expansion}, it follows that}{现在, 由 $\bm p_n$ 的定义, 对某个常数 $c$ 有 $\bm p_{n-1}=c\bm v_n$. 因而 $\bm p_{n-1}$ 与 $\bm v_1,\ldots,\bm v_{n-1}$ 关于 $A^{\mathrm H}A$ 正交. 由 \eqref{eq:gcr-residual-expansion} 得}
\[(A\bm p_{n-1},\bm r_{n-1})=(A\bm p_{n-1},\bm r_0)-\sum_{i=1}^{n-1}y_i(A\bm p_{n-1},A\bm v_i)=(A\bm p_{n-1},\bm r_0).\]
\BT{Now, let us define $\alpha_{n-1}$ and $\beta_{n-1,i}$ as}{现在定义 $\alpha_{n-1}$ 和 $\beta_{n-1,i}$ 为}
\begin{equation}\label{eq:gcr-alpha}\alpha_{n-1}=\frac{(A\bm p_{n-1},\bm r_{n-1})}{(A\bm p_{n-1},A\bm p_{n-1})},\end{equation}
\begin{equation}\label{eq:gcr-beta}\beta_{n-1,i}=-\frac{(A\bm p_{i-1},A\bm r_n)}{(A\bm p_{i-1},A\bm p_{i-1})}.\end{equation}
\BT{Then, we have}{于是有}
\begin{equation}\label{eq:gcr-residual-recurrence}\bm r_n=\bm r_{n-1}-\alpha_{n-1}A\bm p_{n-1},\end{equation}
\begin{equation}\label{eq:gcr-direction-recurrence}\bm p_n=\bm r_n+\sum_{i=1}^n\beta_{n-1,i}\bm p_{i-1}.\end{equation}
\BT{From the recurrence relation of the residual vector, we obtain}{由残差向量的递推关系得到}
\begin{equation}\label{eq:gcr-solution-recurrence}\bm x_n=\bm x_{n-1}+\alpha_{n-1}\bm p_{n-1}.\end{equation}
\BT{To reduce the number of matrix--vector multiplications, use the following recurrences:}{为减少矩阵向量乘法的次数, 使用下列递推式:}
% block: gcr-matvec-recurrence page:132 type:ordinary-paragraph
\begin{equation}\label{eq:gcr-matvec-recurrence}A\bm p_n=A\bm r_n+\sum_{i=1}^n\beta_{n-1,i}A\bm p_{i-1}.\end{equation}
\BT{From \eqref{eq:gcr-alpha}--\eqref{eq:gcr-matvec-recurrence}, the algorithm of the GCR method is obtained in Algorithm\ \ref{alg:gcr}.}{由 \eqref{eq:gcr-alpha}--\eqref{eq:gcr-matvec-recurrence}, 得到 Algorithm\ \ref{alg:gcr} 中的 GCR 算法.}
% block: gcr-algorithm page:132 type:ordinary-paragraph
\begin{algorithm}[H]\caption{\BT{The GCR method}{GCR 方法}}\label{alg:gcr}\begin{algorithmic}[1]
\Statex\textbf{Input:} $\bm x_0\in\mathbb C^N,\ \beta_{-1}=0,\ \bm r_0=\bm b-A\bm x_0,\ \bm p_0=\bm r_0$
\Statex\textbf{Output:} $\bm x_n$
\For{$n=0,1,2,\ldots$, \BT{until convergence}{直至收敛}}
\State $\alpha_n=\frac{(A\bm p_n,\bm r_n)}{(A\bm p_n,A\bm p_n)}$
\State $\bm x_{n+1}=\bm x_n+\alpha_n\bm p_n$
\State $\bm r_{n+1}=\bm r_n-\alpha_nA\bm p_n$
\State $\beta_{n,i}=-\frac{(A\bm p_{i-1},A\bm r_{n+1})}{(A\bm p_{i-1},A\bm p_{i-1})},\quad1\le i\le n+1$
\State $\bm p_{n+1}=\bm r_{n+1}+\sum_{i=1}^{n+1}\beta_{n,i}\bm p_{i-1}$
\State $(A\bm p_{n+1}=A\bm r_{n+1}+\sum_{i=1}^{n+1}\beta_{n,i}A\bm p_{i-1})$
\EndFor\end{algorithmic}\end{algorithm}
% block: gcr-properties-proposition page:132 type:proposition
\BT{The properties of the GCR method are given next.}{下面给出 GCR 方法的性质.}
\par\Needspace{3\baselineskip}
\begin{Proposition}\label{prop:gcr-orthogonality}
\BT{The GCR iterates hold the following properties:}{GCR 迭代具有下列性质:}
\begin{enumerate}[label=(G\arabic*),ref=G\arabic*]
\item\label{part:gcr-direction-orthogonality} $(A\bm p_i,A\bm p_j)=0,\quad i\ne j$,
\item $(\bm r_i,A\bm p_j)=0,\quad i>j$,
\item $(\bm r_i,A\bm p_i)=(\bm r_i,A\bm r_i)$,
\item $(\bm r_i,A\bm r_j)=0,\quad i>j$,
\item $(\bm r_i,A\bm p_i)=(\bm r_0,A\bm p_i),\quad i\ge j$.
\end{enumerate}\end{Proposition}
\BT{For the proof, see \BookCite{54}, Theorem 3.1.}{证明见 \BookCite{54}, Theorem 3.1.}
% block: gcr-restarted-and-orthomin page:132 type:ordinary-paragraph
\BT{From the above properties, we see that the previous derivation is based on the property (\ref{part:gcr-direction-orthogonality}). In exact precision arithmetic, the GCR method, as well as the GMRES method, converges within at most $N$ iterations. However, the GCR method has the same shortcoming as the GMRES method in that the computational work and the required memory increase with the number of iterations. Hence, a restarted version of GCR which is referred to as the GCR($k$) method was proposed in \BookCite{54}. The algorithm is described in Algorithm\ \ref{alg:gcr-restarted}.}{由上述性质可见, 前面的推导基于性质 (\ref{part:gcr-direction-orthogonality}). 在精确运算中, GCR 方法和 GMRES 方法一样, 至多在 $N$ 次迭代内收敛. 然而, GCR 方法与 GMRES 方法具有相同的缺点, 即计算工作量和所需内存随迭代次数增加. 因此 \BookCite{54} 提出了称为 GCR($k$) 方法的重启 GCR 版本. 其算法见 Algorithm\ \ref{alg:gcr-restarted}.}
\BT{Another alternative is to keep only $k$-direction vectors:}{另一种选择是只保留 $k$ 个方向向量:}
\[\bm p_{n+1}=\bm r_{n+1}+\sum_{i=1}^{\min\{k,n+1\}}\beta_{n,i}\bm p_{i-1}.\]
\BT{This method is known as the Orthomin method proposed by Vinsome \BookCite{201}. Later the Orthomin method was referred to as the Orthomin($k$) method in \BookCite{54}. The algorithm of the Orthomin($k$) method is described in Algorithm\ \ref{alg:orthomin}.}{这一方法称为 Vinsome \BookCite{201} 提出的 Orthomin 方法. 后来 \BookCite{54} 将 Orthomin 方法称为 \linebreak[4]Orthomin($k$) 方法. Orthomin($k$) 方法的算法见 Algorithm\ \ref{alg:orthomin}.}
% block: gcr-restarted-algorithm page:133 type:ordinary-paragraph
\begin{algorithm}[H]\caption{\BT{The GCR($k$) method}{GCR($k$) 方法}}\label{alg:gcr-restarted}\begin{algorithmic}[1]
\Statex\textbf{Input:} $\bm x_0\in\mathbb C^N,\ \bm r_0=\bm b-A\bm x_0,\ \bm p_0=\bm r_0$
\Statex\textbf{Output:} $\bm x_{k+1}$
\For{$n=0,1,\ldots,k$ \BT{until convergence}{直至收敛}}
\State $\alpha_n=\frac{(A\bm p_n,\bm r_n)}{(A\bm p_n,A\bm p_n)}$
\State $\bm x_{n+1}=\bm x_n+\alpha_n\bm p_n$
\State $\bm r_{n+1}=\bm r_n-\alpha_nA\bm p_n$
\State $\beta_{n,i}=-\frac{(A\bm p_{i-1},A\bm r_{n+1})}{(A\bm p_{i-1},A\bm p_{i-1})},\quad1\le i\le n+1$
\State $\bm p_{n+1}=\bm r_{n+1}+\sum_{i=1}^{n+1}\beta_{n,i}\bm p_{i-1}$
\State $(A\bm p_{n+1}=A\bm r_{n+1}+\sum_{i=1}^{n+1}\beta_{n,i}A\bm p_{i-1})$
\EndFor\State $\bm x_0=\bm x_{k+1}$\State\BT{repeat}{重复}
\end{algorithmic}\end{algorithm}
% block: orthomin-algorithm page:133 type:ordinary-paragraph
\begin{algorithm}[H]\caption{\BT{The Orthomin($k$) method}{Orthomin($k$) 方法}}\label{alg:orthomin}\begin{algorithmic}[1]
\Statex\textbf{Input:} $\bm x_0\in\mathbb C^N,\ \bm r_0=\bm b-A\bm x_0,\ \bm p_0=\bm r_0$
\Statex\textbf{Output:} $\bm x_n$
\For{$n=0,1,\ldots$, \BT{until $\|\bm r_n\|\le\epsilon\|\bm b\|$}{直至 $\|\bm r_n\|\le\epsilon\|\bm b\|$}}
\State $\alpha_n=\frac{(A\bm p_n,\bm r_n)}{(A\bm p_n,A\bm p_n)}$
\State $\bm x_{n+1}=\bm x_n+\alpha_n\bm p_n$
\State $\bm r_{n+1}=\bm r_n-\alpha_nA\bm p_n$
\State $\beta_{n,i}=-\frac{(A\bm p_{i-1},A\bm r_{n+1})}{(A\bm p_{i-1},A\bm p_{i-1})},\quad n-k+2\le i\le n+1$
\State $\bm p_{n+1}=\bm r_{n+1}+\sum_{i=1}^{\min\{k,n+1\}}\beta_{n,i}\bm p_{i-1}$
\State $(A\bm p_{n+1}=A\bm r_{n+1}+\sum_{i=1}^{n+1}\beta_{n,i}A\bm p_{i-1})$
\EndFor\end{algorithmic}\end{algorithm}
% block: gcr-convergence-proposition pages:133,134 type:proposition
\BT{The convergence analysis of the GCR method, the GCR($k$) method, and the Orthomin($k$) method is given in Proposition\ \ref{prop:gcr-restarted-convergence}.}{GCR 方法, GCR($k$) 方法和 Orthomin($k$) 方法的收敛性分析见 Proposition\ \ref{prop:gcr-restarted-convergence}.}
\par\Needspace{3\baselineskip}
\begin{Proposition}\label{prop:gcr-restarted-convergence}
\BT{Let $M$ be the symmetric part of $A\in\mathbb R^{N\times N}$ and let $R$ be the skew-symmetric part of $A$, i.e., $M=(A+A^{\top})/2$, $R=(A-A)^{\top}/2$. If $M$ is positive definite, then $\{\bm r_n\}$ that are the sequence of residuals generated by the GCR method, the GCR($k$) method, or the Orthomin($k$) method satisfy}{设 $M$ 为 $A\in\mathbb R^{N\times N}$ 的对称部分, $R$ 为 $A$ 的斜对称部分, 即 $M=(A+A^{\top})/2$, $R=(A-A)^{\top}/2$. 若 $M$ 正定, 则 GCR 方法, GCR($k$) 方法或 Orthomin($k$) 方法生成的残差序列 $\{\bm r_n\}$ 满足}
\[\|\bm r_n\|\le\left(1-\frac{\lambda_{\min}(M)^2}{\lambda_{\max}(A^{\top}A)}\right)^{n/2}\|\bm r_0\|\]
\BT{and}{以及}
\[\|\bm r_n\|\le\left(1-\frac{\lambda_{\min}(M)^2}{\lambda_{\min}(M)\lambda_{\max}(M)+\rho(R)^2}\right)^{n/2}\|\bm r_0\|,\]
\BT{where $\lambda_{\min}(M)$ and $\lambda_{\max}(M)$ are the minimum and maximum eigenvalues of the symmetric positive definite matrix $M$, and $\rho(R)$ is the spectral radius of $R$.}{其中 $\lambda_{\min}(M)$ 和 $\lambda_{\max}(M)$ 为对称正定矩阵 $M$ 的最小和最大特征值, $\rho(R)$ 为 $R$ 的谱半径.}
\end{Proposition}
\BT{See \eqref{eq:stationary-spectral-radius} for the spectral radius. For the proof, see \BookCite{54}, Theorem 4.4.}{谱半径见 \eqref{eq:stationary-spectral-radius}. 证明见 \BookCite{54}, Theorem 4.4.}
\BT{We can see from Proposition\ \ref{prop:gcr-restarted-convergence} that if $A$ is close to $I$, then the three methods converge fast since $\lambda_{\min}(M)\approx\lambda_{\max}(M)\approx1$ and $R\approx O$.}{由 Proposition\ \ref{prop:gcr-restarted-convergence} 可见, 若 $A$ 接近 $I$, 则这三种方法收敛很快, 因为 $\lambda_{\min}(M)\approx\lambda_{\max}(M)\approx1$ 且 $R\approx O$.}
% block: fom-heading page:134 type:subsection
\subsubsection{\BT{The Full Orthogonalization Method (FOM)}{完全正交化方法 (FOM)}}\label{subsec:fom}
\BT{The FOM \BookCite{149} is a generalization of the CG method. As well as the CG method, the FOM finds approximate solutions such that}{FOM \BookCite{149} 是 CG 方法的推广. 与 CG 方法一样, FOM 寻找满足下列条件的近似解}
\[\begin{aligned}\bm x_n&=\bm x_0+\bm z_n,\quad\bm z_n\in\mathcal K_n(A,\bm r_0),\\\bm r_n&\perp\mathcal K_n(A,\bm r_0).\end{aligned}\]
\BT{The FOM uses the Arnoldi method in Sect.\ \ref{subsec:arnoldi-process} to produce the orthonormal basis vectors of $\mathcal K_n(A,\bm r_0)$. Thus the residual vector $\bm r_n$ has the same form as \eqref{eq:gmres-projected-residual}, i.e.,}{FOM 使用 Section\ \ref{subsec:arnoldi-process} 中的 Arnoldi 方法生成 $\mathcal K_n(A,\bm r_0)$ 的标准正交基向量. 因而残差向量 $\bm r_n$ 与 \eqref{eq:gmres-projected-residual} 具有相同的形式, 即,}
\begin{equation}\label{eq:fom-projected-residual}\bm r_n=V_{n+1}(\beta\bm e_1-H_{n+1,n}\bm y_n).\end{equation}
\BT{The difference between the GMRES method and the FOM is that the GMRES method finds $\bm y_n$ such that the 2-norm of the residual vector $\|\bm r_n\|$ is minimized, and the FOM finds $\bm y_n$ such that the residual vector is orthogonal to the $n$ dimensional Krylov subspace $\bm r_n\perp\mathcal K_n(A,\bm r_0)$. Let $\bm v_i$ be the $i$th column of $V_{n+1}$. Then}{GMRES 方法与 FOM 的区别在于, GMRES 方法寻找使残差向量的 2-范数 $\|\bm r_n\|$ 达到极小值的 $\bm y_n$, 而 FOM 寻找使残差向量与 $n$ 维 Krylov 子空间正交的 $\bm y_n$, 即 $\bm r_n\perp\mathcal K_n(A,\bm r_0)$. 设 $\bm v_i$ 为 $V_{n+1}$ 的第 $i$ 列. 则}
\[\begin{aligned}V_n^{\mathrm H}\bm r_n=\bm0&\Leftrightarrow\bm v_i^{\mathrm H}\bm r_n=0\quad\text{for }i=1,2,\ldots,n\\&\Leftrightarrow\bm r_n^{\mathrm H}(c_1\bm v_1+c_2\bm v_2+\cdots+c_n\bm v_n)=0\quad\text{for all }c_i\in\mathbb C\\&\Leftrightarrow\bm r_n\perp\mathcal K_n(A,\bm r_0).\end{aligned}\]
% block: fom-projected-solve page:134 type:ordinary-paragraph
\BT{Note that the statement $\bm v_i^{\mathrm H}\bm r_n=0\Leftarrow\bm r_n^{\mathrm H}(c_1\bm v_1+c_2\bm v_2+\cdots+c_n\bm v_n)=0$ can be shown by setting $c_i=1$ and $c_j=0$ for $i\ne j$. From this, finding $\bm y_n$ such that $\bm r_n\perp\mathcal K_n(A,\bm r_0)$ is equivalent to determining $\bm y_n$ such that $V_n^{\mathrm H}\bm r_n=\bm0$. Then, from \eqref{eq:fom-projected-residual} it follows that}{注意, 令 $c_i=1$, 且对 $i\ne j$ 令 $c_j=0$, 就可说明断言 $\bm v_i^{\mathrm H}\bm r_n=0\Leftarrow\bm r_n^{\mathrm H}(c_1\bm v_1+c_2\bm v_2+\cdots+c_n\bm v_n)=0$. 由此, 寻找使 $\bm r_n\perp\mathcal K_n(A,\bm r_0)$ 的 $\bm y_n$, 等价于确定使 $V_n^{\mathrm H}\bm r_n=\bm0$ 的 $\bm y_n$. 于是由 \eqref{eq:fom-projected-residual} 得}
\[\bm0=V_n^{\mathrm H}\bm r_n=V_n^{\mathrm H}V_{n+1}(\beta\bm e_1-H_{n+1,n}\bm y_n)=\beta\bm e_1-H_{n,n}\bm y_n.\]
\BT{Thus $\bm y_n$ can be obtained by solving the following linear systems:}{因此 $\bm y_n$ 可通过求解下列线性方程组得到:}
\[H_{n,n}\bm y_n=\beta\bm e_1.\]
\BT{Since $H_{n,n}$ is a Hessenberg matrix, the linear systems can be solved by the LU decomposition with the computational cost of $O(n^2)$. After the advent of the GMRES method, the FOM became less attractive because, unlike the GMRES method, the FOM may suffer from breakdown due to using the LU decomposition without pivoting. On the other hand, the FOM is still attractive to use for solving shifted linear systems in Chap.\ \ref{chap:shifted-linear-systems}, when considering the restart.}{由于 $H_{n,n}$ 为 Hessenberg 矩阵, 可通过 LU 分解以 $O(n^2)$ 的计算开销求解该线性方程组. GMRES 方法出现后, FOM 的吸引力有所降低, 因为与 GMRES 方法不同, FOM 可能因使用不选主元的 LU 分解而发生中断. 另一方面, 在考虑重启时, 用 FOM 求解 Chapter\ \ref{chap:shifted-linear-systems} 中的移位线性方程组仍有吸引力.}
% block: product-type-heading page:135 type:subsection
\subsubsection{\BT{Product-Type Krylov Subspace Methods}{乘积型 Krylov 子空间方法}}\label{subsec:product-type-krylov}
\BT{As seen in Sect.\ \ref{subsec:bicg-method}, the BiCG method requires $A^{\mathrm H}$ for solving $A\bm x=\bm b$. P. Sonneveld found that squaring the BiCG polynomials leads to an algorithm without using the information of $A^{\mathrm H}$, and the resulting algorithm is known as the CGS method \BookCite{173}. Though the convergence depends on linear systems, the CGS method is much faster than the BiCG method when the BiCG method shows smooth convergence behavior. However, the CGS method sometimes has much irregular convergence and may suffer from loss of accuracy of the approximate solution or may diverge. H. A. van der Vorst observed that squaring the BiCG polynomials can be replaced by the multiplication of degree-one minimum residual polynomials and the BiCG polynomial, leading to the BiCGSTAB method \BookCite{195}, which tends to have smoother convergence behavior than the CGS method. M. H. Gutknecht proposed the multiplication of degree-two minimum residual polynomials and the BiCG polynomials, which is known as the BiCGSTAB2 method \BookCite{87}. G. L. G. Sleijpen and D. R. Fokkema proposed the multiplication of degree-$\ell$ minimal residual polynomials and the BiCG polynomials, which is known as the BiCGSTAB($\ell$) method \BookCite{162}. S.-L. Zhang constructed a framework that includes the CGS method, the BiCGSTAB method, and the BiCGSTAB2 method. In the framework, the GPBiCG method was proposed in \BookCite{212}. These methods are referred to as product-type methods (based on the Bi-CG method) first named by S.-L. Zhang.}{如 Section\ \ref{subsec:bicg-method} 所示, BiCG 方法求解 $A\bm x=\bm b$ 时需要 $A^{\mathrm H}$. P. Sonneveld 发现, 将 BiCG 多项式平方可得到一种不使用 $A^{\mathrm H}$ 信息的算法, 所得算法称为 CGS 方法 \BookCite{173}. 虽然收敛性取决于线性方程组, 但当 BiCG 方法表现出平滑的收敛行为时, CGS 方法比 BiCG 方法快得多. 然而, CGS 方法有时表现出很不规则的收敛, 可能发生近似解精度损失或发散. H. A. van der Vorst 观察到, BiCG 多项式的平方可替换为一次最小残差多项式与 BiCG 多项式的乘积, 由此得到 BiCGSTAB 方法 \BookCite{195}, 它往往比 CGS 方法具有更平滑的收敛行为. M. H. Gutknecht 提出二次最小残差多项式与 BiCG 多项式相乘, 称为 BiCGSTAB2 方法 \BookCite{87}. G. L. G. Sleijpen 和 D. R. Fokkema 提出 $\ell$ 次最小残差多项式与 BiCG 多项式相乘, 称为 BiCGSTAB($\ell$) 方法 \BookCite{162}. S.-L. Zhang 构造了包含 CGS 方法, BiCGSTAB 方法和 BiCGSTAB2 方法的框架. 在这一框架中, 文献 \BookCite{212} 提出了 GPBiCG 方法. 这些方法称为 (基于 Bi-CG 方法的) 乘积型方法, 这一名称首先由 S.-L. Zhang 提出.}
\BT{In this section, Zhang's framework is explained, based on Zhang's book in \BookCite{72}.}{本节根据 Zhang 的著作 \BookCite{72} 介绍 Zhang 的框架.}
% block: product-bicg-polynomials page:135 type:ordinary-paragraph
\BT{As seen in \eqref{eq:cg-residual-direction-polynomials}, the residual vector of the CG method can be written as the multiplication of Lanczos polynomials and the initial residual vector $\bm r_0$. Similarly, the residual vector of the BiCG method can also be written as follows:}{如 \eqref{eq:cg-residual-direction-polynomials} 所示, CG 方法的残差向量可写为 Lanczos 多项式与初始残差向量 $\bm r_0$ 的乘积. 类似地, BiCG 方法的残差向量也可写为}
\[\bm r_n^{\mathrm{BiCG}}=R_n(A)\bm r_0,\]
\BT{where}{其中}
\begin{equation}\label{eq:product-bicg-polynomials-initial}R_0(\lambda)=1,\quad P_0(\lambda)=1,\end{equation}
\begin{equation}\label{eq:product-bicg-residual-polynomial}R_n(\lambda)=R_{n-1}(\lambda)-\alpha_{n-1}\lambda P_{n-1}(\lambda),\end{equation}
\begin{equation}\label{eq:product-bicg-direction-polynomial}P_n(\lambda)=R_n(\lambda)+\beta_{n-1}P_{n-1}(\lambda)\quad n=1,2,\ldots\end{equation}
\BT{The differences between the (matrix) polynomial representation of the CG method and that of the BiCG method are $\alpha_i$ and $\beta_i$ of the Lanczos polynomials. The $n$th residual vector of product-type methods based on the BiCG method is defined by the product of $n$ degree polynomials and the $n$th residual vector of the BiCG method as follows:}{CG 方法与 BiCG 方法的 (矩阵) 多项式表示的区别在于 Lanczos 多项式中的 $\alpha_i$ 和 $\beta_i$. 基于 BiCG 方法的乘积型方法的第 $n$ 个残差向量, 定义为 $n$ 次多项式与 BiCG 方法的第 $n$ 个残差向量的乘积, 如下所示:}
\begin{equation}\label{eq:product-residual-polynomial}\bm r_n=H_n(A)\bm r_n^{\mathrm{BiCG}}=H_n(A)R_n(A)\bm r_0.\end{equation}
% block: product-polynomial-restructuring page:136 type:subsection
\paragraph{\BT{Restructuring of Residual Polynomials}{残差多项式的重构}}\label{subsec:product-polynomial-restructuring}\mbox{}\par
\BT{S.-L. Zhang proposed $H_n$ in \eqref{eq:product-residual-polynomial} as the following three-term recurrence relations:}{S.-L. Zhang 提出将 \eqref{eq:product-residual-polynomial} 中的 $H_n$ 定为下列三项递推关系:}
\begin{equation}\label{eq:product-h-initial}H_0(\lambda):=1,\end{equation}
\begin{equation}\label{eq:product-h-first}H_1(\lambda):=(1-\zeta_0\lambda)H_0(\lambda),\end{equation}
\begin{equation}\label{eq:product-h-three-term}H_n(\lambda):=(1+\eta_{n-1}-\zeta_{n-1}\lambda)H_{n-1}(\lambda)-\eta_{n-1}H_{n-2}(\lambda),\quad n=2,3,\ldots\end{equation}
\BT{We see that $H_n$ is similar to the Lanczos polynomials \eqref{eq:cg-polynomial-initial-zero}--\eqref{eq:cg-polynomial-three-term}.}{可见 $H_n$ 与 Lanczos 多项式 \eqref{eq:cg-polynomial-initial-zero}--\eqref{eq:cg-polynomial-three-term} 相似.}
\BT{In the following, $H_n$ is rewritten as the form of coupled two-term recurrence relation. We introduce $G_{n-1}$ as}{下面将 $H_n$ 改写为耦合的两项递推关系的形式. 引入 $G_{n-1}$ 为}
\[G_{n-1}(\lambda):=\frac{H_{n-1}(\lambda)-H_n(\lambda)}{\lambda},\]
\BT{and we rewrite \eqref{eq:product-h-three-term} as}{并将 \eqref{eq:product-h-three-term} 改写为}
\[-(H_{n-1}(\lambda)-H_n(\lambda))=-\zeta_{n-1}\lambda H_{n-1}(\lambda)-\eta_{n-1}(H_{n-2}(\lambda)-H_{n-1}(\lambda)).\]
\BT{Then $H_n(\lambda)$ in \eqref{eq:product-h-initial}--\eqref{eq:product-h-three-term} can be rewritten as}{于是 \eqref{eq:product-h-initial}--\eqref{eq:product-h-three-term} 中的 $H_n(\lambda)$ 可改写为}
\begin{equation}\label{eq:product-coupled-initial}H_0(\lambda)=1,\quad G_0(\lambda)=\zeta_0,\end{equation}
\begin{equation}\label{eq:product-coupled-h}H_n(\lambda)=H_{n-1}(\lambda)-\lambda G_{n-1}(\lambda),\end{equation}
\begin{equation}\label{eq:product-coupled-g}G_n(\lambda)=\zeta_nH_n(\lambda)+\eta_nG_{n-1}(\lambda),\quad n=1,2,\ldots\end{equation}
% block: product-iterate-polynomial-recurrences page:136 type:subsection
\paragraph{\BT{Recurrence Formulas for the Iterates of the Product-Type Methods}{乘积型方法迭代量的递推公式}}\label{subsec:product-iterate-recurrences}\mbox{}\par
\BT{Since recurrence relations of $R_n,P_n,H_n,G_n$ are already given in \eqref{eq:product-bicg-polynomials-initial}--\eqref{eq:product-bicg-direction-polynomial} and in \eqref{eq:product-coupled-initial}--\eqref{eq:product-coupled-g}, the residual vector $\bm r_{n+1}=H_{n+1}(A)R_{n+1}(A)\bm r_0$ can be computed by the previous residual vector $\bm r_n=H_n(A)R_n(A)\bm r_0$ using the recurrence relations. In fact, $H_{n+1}R_{n+1}$ can be expanded as follows:}{由于 \eqref{eq:product-bicg-polynomials-initial}--\eqref{eq:product-bicg-direction-polynomial} 及 \eqref{eq:product-coupled-initial}--\eqref{eq:product-coupled-g} 已给出 $R_n,P_n,H_n,G_n$ 的递推关系, 可以使用这些递推关系, 由前一个残差向量 $\bm r_n=H_n(A)R_n(A)\bm r_0$ 计算残差向量 $\bm r_{n+1}=H_{n+1}(A)R_{n+1}(A)\bm r_0$. 事实上, $H_{n+1}R_{n+1}$ 可作如下展开:}
\begin{equation}\label{eq:product-polynomial-residual-three}H_{n+1}R_{n+1}=H_nR_{n+1}-\eta_n\lambda G_{n-1}R_{n+1}-\zeta_n\lambda H_nR_{n+1}\end{equation}
\begin{equation}\label{eq:product-polynomial-residual-two}=H_nR_n-\alpha_n\lambda H_nP_n-\lambda G_nR_{n+1},\end{equation}
\begin{equation}\label{eq:product-polynomial-t}H_nR_{n+1}=H_nR_n-\alpha_n\lambda H_nP_n,\end{equation}
\begin{equation}\label{eq:product-polynomial-y}\lambda G_nR_{n+2}=H_nR_{n+1}-H_{n+1}R_{n+1}-\alpha_{n+1}\lambda H_nP_{n+1}+\alpha_{n+1}\lambda H_{n+1}P_{n+1},\end{equation}
\begin{equation}\label{eq:product-polynomial-p}H_{n+1}P_{n+1}=H_{n+1}R_{n+1}+\beta_nH_nP_n-\beta_n\lambda G_nP_n,\end{equation}
\begin{equation}\label{eq:product-polynomial-w}\lambda H_nP_{n+1}=\lambda H_nR_{n+1}+\beta_n\lambda H_nP_n,\end{equation}
\begin{equation}\label{eq:product-polynomial-u}\lambda G_nP_n=\zeta_n\lambda H_nP_n+\eta_n(H_{n-1}R_n-H_nR_n+\beta_{n-1}\lambda G_{n-1}P_{n-1}),\end{equation}
\begin{equation}\label{eq:product-polynomial-z}G_nR_{n+1}=\zeta_nH_nR_n+\eta_nG_{n-1}R_n-\alpha_n\lambda G_nP_n.\end{equation}
% block: product-auxiliary-vector-recurrences page:137 type:ordinary-paragraph
\BT{Let us define new auxiliary vectors as}{定义新的辅助向量为}
\[\begin{aligned}\bm t_n&:=H_n(A)\bm r_{n+1}^{\mathrm{BiCG}},&\bm y_n&:=AG_{n-1}(A)\bm r_{n+1}^{\mathrm{BiCG}},\\\bm p_n&:=H_n(A)\bm p_n^{\mathrm{BiCG}},&\bm w_n&:=AH_n(A)\bm p_{n+1}^{\mathrm{BiCG}},\\\bm u_n&:=AG_n(A)\bm p_n^{\mathrm{BiCG}},&\bm z_n&:=G_n(A)\bm r_{n+1}^{\mathrm{BiCG}}.\end{aligned}\]
\BT{Then, from \eqref{eq:product-polynomial-residual-three}--\eqref{eq:product-polynomial-z} we have the following recurrence formulas:}{则由 \eqref{eq:product-polynomial-residual-three}--\eqref{eq:product-polynomial-z} 有下列递推公式:}
\begin{equation}\label{eq:product-residual-three}\bm r_{n+1}=\bm t_n-\eta_n\bm y_n-\zeta_nA\bm t_n\end{equation}
\begin{equation}\label{eq:product-residual-two}=\bm r_n-\alpha_nA\bm p_n-A\bm z_n,\end{equation}
\begin{equation}\label{eq:product-t}\bm t_n=\bm r_n-\alpha_nA\bm p_n,\end{equation}
\begin{equation}\label{eq:product-y}\bm y_{n+1}=\bm t_n-\bm r_{n+1}-\alpha_{n+1}\bm w_n+\alpha_{n+1}A\bm p_{n+1},\end{equation}
\begin{equation}\label{eq:product-p}\bm p_{n+1}=\bm r_{n+1}+\beta_n(\bm p_n-\bm u_n),\end{equation}
\begin{equation}\label{eq:product-w}\bm w_n=A\bm t_n+\beta_nA\bm p_n,\end{equation}
\begin{equation}\label{eq:product-u}\bm u_n=\zeta_nA\bm p_n+\eta_n(\bm t_{n-1}-\bm r_n+\beta_{n-1}\bm u_{n-1}),\end{equation}
\begin{equation}\label{eq:product-z}\bm z_n=\zeta_n\bm r_n+\eta_n\bm z_{n-1}-\alpha_n\bm u_n.\end{equation}
\BT{Two recurrence relations, \eqref{eq:product-residual-three} and \eqref{eq:product-residual-two}, are described for the residual vector and the approximate solution. The first recurrence relation \eqref{eq:product-residual-three} will be used to determine the two parameters $\zeta_n$ and $\eta_n$, and the second recurrence relation \eqref{eq:product-residual-two} can be used to obtain the approximate solution as follows:}{针对残差向量和近似解, 给出了两个递推关系 \eqref{eq:product-residual-three} 和 \eqref{eq:product-residual-two}. 第一个递推关系 \eqref{eq:product-residual-three} 将用于确定两个参数 $\zeta_n$ 和 $\eta_n$, 第二个递推关系 \eqref{eq:product-residual-two} 可用于按下式得到近似解:}
\[\bm x_{n+1}=\bm x_n+\alpha_n\bm p_n+\bm z_n.\]
% block: product-alpha-beta page:137 type:subsection
\paragraph{\BT{Computations for $\alpha_n$ and $\beta_n$}{$\alpha_n$ 和 $\beta_n$ 的计算}}\label{subsec:product-alpha-beta}\mbox{}\par
\BT{Since the coefficient of the highest-order term of $H_n$ is $(-1)^n\prod_{i=0}^{n-1}\zeta_i$, we obtain}{由于 $H_n$ 最高次项的系数为 $(-1)^n\prod_{i=0}^{n-1}\zeta_i$, 得到}
\[\begin{aligned}(\bm r_0^{*},\bm r_n)&=(\bar H_n(A^{\mathrm H})\bm r_0^{*},\bm r_n^{\mathrm{BiCG}})=\left((-1)^n\prod_{i=0}^{n-1}\zeta_i\right)((A^{\mathrm H})^n\bm r_0^{*},\bm r_n^{\mathrm{BiCG}}),\\(\bm r_0^{*},A\bm p_n)&=(\bar H_n(A^{\mathrm H})\bm r_0^{*},A\bm p_n^{\mathrm{BiCG}})=\left((-1)^n\prod_{i=0}^{n-1}\zeta_i\right)((A^{\mathrm H})^n\bm r_0^{*},A\bm p_n^{\mathrm{BiCG}}).\end{aligned}\]
\BT{Thus, from \eqref{eq:bicg-alpha-shadow-power} and \eqref{eq:bicg-beta-shadow-power} it follows that}{因此由 \eqref{eq:bicg-alpha-shadow-power} 和 \eqref{eq:bicg-beta-shadow-power} 得}
\[\alpha_n=\frac{(\bm r_0^{*},\bm r_n)}{(\bm r_0^{*},A\bm p_n)},\qquad\beta_n=\frac{\alpha_n}{\zeta_n}\times\frac{(\bm r_0^{*},\bm r_{n+1})}{(\bm r_0^{*},\bm r_n)}.\]
% block: product-implementation page:138 type:subsection
\paragraph{\BT{Implementation Details}{实现细节}}\label{subsec:product-implementation}\mbox{}\par
\BT{Krylov subspace methods, such as the CGS method, the BiCGSTAB method, the BiCGSTAB2 method, and the GPBiCG method, can be derived from Zhang's framework by determining the two parameters $\zeta_n$ and $\eta_n$. The relations among these Krylov subspace methods and the parameters are listed in Table\ \ref{tab:product-parameters}.}{通过确定两个参数 $\zeta_n$ 和 $\eta_n$, 可从 Zhang 的框架导出 CGS 方法, BiCGSTAB 方法, BiCGSTAB2 方法和 GPBiCG 方法等 Krylov 子空间方法. 这些 Krylov 子空间方法与参数的关系列于 Table\ \ref{tab:product-parameters}.}
% block: product-parameters-table page:138 type:table
\begin{table}[H]\centering
\caption{\BT{Choices of $\zeta_n$ and $\eta_n$ for the product-type methods}{乘积型方法中 $\zeta_n$ 和 $\eta_n$ 的选择}}\label{tab:product-parameters}
\begin{tabular}{l|l}\hline
CGS&$\zeta_n=\alpha_n,\quad\eta_n=\dfrac{\beta_{n-1}}{\alpha_{n-1}}\alpha_n$\\\hline
BiCGSTAB&$\zeta_n=\operatorname*{arg\,min}_{\zeta_n\in\mathbb C}\|\bm r_{n+1}\|,\quad\eta_n=0$\\\hline
GPBiCG&$\{\zeta_n,\eta_n\}=\operatorname*{arg\,min}_{\zeta_n,\eta_n\in\mathbb C}\|\bm r_{n+1}\|$\\\hline
BiCGSTAB2&\BT{BiCGSTAB parameters at even iterations}{偶数次迭代采用 BiCGSTAB 参数}\\
&\BT{GPBiCG parameters at odd iterations}{奇数次迭代采用 GPBiCG 参数}\\\hline
\end{tabular}\end{table}
% block: cgs-parameter-choice page:138 type:subsection
\paragraph{\BT{The Choice for the CGS Method}{CGS 方法的选择}}\label{subsec:cgs-parameter-choice}\mbox{}\par
\BT{Setting $\zeta_n=\alpha_n$ and $\eta_n=(\beta_{n-1}/\alpha_{n-1})\alpha_n$ in recurrence relations \eqref{eq:product-coupled-h} and \eqref{eq:product-coupled-g} yields the CGS method, i.e., the polynomials $H_n$ and $G_n$ become the Lanczos polynomials:}{在递推关系 \eqref{eq:product-coupled-h} 和 \eqref{eq:product-coupled-g} 中置 $\zeta_n=\alpha_n$ 且 $\eta_n=(\beta_{n-1}/\alpha_{n-1})\alpha_n$, 就得到 CGS 方法, 即多项式 $H_n$ 和 $G_n$ 变为 Lanczos 多项式:}
\[H_n(\lambda)=R_n(\lambda),\qquad G_n(\lambda)=\alpha_nP_n(\lambda).\]
\BT{This fact leads to the relation $\bm p_n-\bm u_n=\bm z_n/\alpha_n$, since from the definitions of auxiliary vectors we have}{这一事实得到关系 $\bm p_n-\bm u_n=\bm z_n/\alpha_n$, 因为由辅助向量的定义有}
\[\begin{aligned}\bm p_n-\bm u_n&=H_n(A)P_n(A)\bm r_0-AG_n(A)P_n(A)\bm r_0\\&=R_n(A)P_n(A)\bm r_0-\alpha_nAP_n(A)P_n(A)\bm r_0\\&=(R_n(A)-\alpha_nAP_n(A))P_n(A)\bm r_0\\&=R_{n+1}(A)P_n(A)\bm r_0\\&=P_n(A)R_{n+1}(A)\bm r_0\\&=P_n(A)\bm r_{n+1}^{\mathrm{BiCG}}\\&=\frac1{\alpha_n}\times G_n(A)\bm r_{n+1}^{\mathrm{BiCG}}\\&=\frac1{\alpha_n}\times\bm z_n.\end{aligned}\]
% block: cgs-auxiliary-rewrite page:139 type:ordinary-paragraph
\BT{Notice that $\bm t_{n-1}-\bm r_n=A\bm z_{n-1}$ from \eqref{eq:product-residual-two} and \eqref{eq:product-t}. We use the recurrence formula \eqref{eq:product-residual-two} to update $\bm r_{n+1}$, and then the auxiliary vectors $\bm t_n,\bm y_n$, and $\bm w_n$ are not required in the recurrence formulas \eqref{eq:product-residual-three}--\eqref{eq:product-z}. Now we introduce the following auxiliary vectors:}{注意, 由 \eqref{eq:product-residual-two} 和 \eqref{eq:product-t} 有 $\bm t_{n-1}-\bm r_n=A\bm z_{n-1}$. 使用递推公式 \eqref{eq:product-residual-two} 更新 $\bm r_{n+1}$ 后, 在递推公式 \eqref{eq:product-residual-three}--\eqref{eq:product-z} 中就不再需要辅助向量 $\bm t_n,\bm y_n$ 和 $\bm w_n$. 现在引入下列辅助向量:}
\[\widetilde{\bm u}_n:=A^{-1}\bm u_n/\alpha_n,\qquad\widetilde{\bm z}_n:=\bm z_n/\alpha_n,\]
\BT{instead of $\bm u_n$ and $\bm z_n$. By using the relation $(\bm r_0^{*},A\bm p_n)=(\bm r_0^{*},\bm u_n/\alpha_n)$, the CGS method can be obtained and the algorithm is described in Algorithm\ \ref{alg:cgs}. Here $\widetilde{\bm u}_n$ and $\widetilde{\bm z}_n$ were rewritten as $\bm u_n$ and $\bm z_n$.}{以代替 $\bm u_n$ 和 $\bm z_n$. 使用关系 $(\bm r_0^{*},A\bm p_n)=(\bm r_0^{*},\bm u_n/\alpha_n)$ 即可得到 CGS 方法, 其算法见 Algorithm\ \ref{alg:cgs}. 这里已将 $\widetilde{\bm u}_n$ 和 $\widetilde{\bm z}_n$ 改写为 $\bm u_n$ 和 $\bm z_n$.}
% block: cgs-algorithm page:139 type:ordinary-paragraph
\begin{algorithm}[H]\caption{\BT{The CGS method}{CGS 方法}}\label{alg:cgs}\begin{algorithmic}[1]
\Statex\textbf{Input:} $\bm x_0\in\mathbb C^N,\ \beta_{-1}=0,\ \bm u_{-1}=\bm z_{-1}=\bm0,\ \bm r_0=\bm b-A\bm x_0$
\Statex\textbf{Input:} \BT{Choose $\bm r_0^{*}\in\mathbb C^N$, e.g., $\bm r_0^{*}=\bm r_0$}{选择 $\bm r_0^{*}\in\mathbb C^N$, 例如 $\bm r_0^{*}=\bm r_0$}
\Statex\textbf{Output:} $\bm x_n$
\For{$n=0,1,\ldots$}
\State $\bm p_n=\bm r_n+\beta_{n-1}\bm z_{n-1}$
\State $\bm u_n=\bm p_n+\beta_{n-1}(\bm z_{n-1}+\beta_{n-1}\bm u_{n-1})$
\State $\alpha_n=\frac{(\bm r_0^{*},\bm r_n)}{(\bm r_0^{*},A\bm u_n)}$
\State $\bm z_n=\bm p_n-\alpha_nA\bm u_n$
\State $\bm x_{n+1}=\bm x_n+\alpha_n(\bm p_n+\bm z_n)$
\State $\bm r_{n+1}=\bm r_n-\alpha_nA(\bm p_n+\bm z_n)$
\State $\beta_n=\frac{(\bm r_0^{*},\bm r_{n+1})}{(\bm r_0^{*},\bm r_n)}$
\EndFor\end{algorithmic}\end{algorithm}
\BT{The residual vector of the CGS method is written as $\bm r_{n+1}=R_{n+1}(A)R_{n+1}(A)\bm r_0$. The CGS method is short for the (Bi)CG squared method, which comes from the term $R_{n+1}(A)R_{n+1}(A)$. From this, the CGS method is expected to converge as twice as fast as the BiCG method. On the other hand, when the BiCG method shows irregular convergence behavior as is often the case, the CGS method gives much irregular convergence behavior, leading to much less accurate approximate solutions or no convergence.}{CGS 方法的残差向量写为 $\bm r_{n+1}=R_{n+1}(A)R_{n+1}(A)\bm r_0$. CGS 方法是 (Bi)CG 平方方法的缩写, 名称来自项 $R_{n+1}(A)R_{n+1}(A)$. 由此, 可以预期 CGS 方法的收敛速度为 BiCG 方法的两倍. 另一方面, 当 BiCG 方法像通常那样表现出不规则的收敛行为时, CGS 方法的收敛行为会很不规则, 导致近似解的精度低得多, 或不收敛.}
% block: bicgstab-parameter-choice pages:139,140 type:subsection
\paragraph{\BT{The Choice for the BiCGSTAB Method}{BiCGSTAB 方法的选择}}\label{subsec:bicgstab-parameter-choice}\mbox{}\par
\BT{If we choose $\eta_n=0$ and $\zeta_n$ such that the 2-norm of the residual vector $\|\bm r_{n+1}\|=\|\bm t_n-\zeta_nA\bm t_n\|$ is minimized, then the BiCGSTAB method is obtained. To achieve the minimization, we choose $\zeta_n$ such that the following property holds:}{若选择 $\eta_n=0$, 并选择 $\zeta_n$ 使残差向量的 2-范数 $\|\bm r_{n+1}\|=\|\bm t_n-\zeta_nA\bm t_n\|$ 达到极小值, 就得到 BiCGSTAB 方法. 为实现这一极小化, 选择 $\zeta_n$ 使下列性质成立:}
\[\bm r_{n+1}\perp A\bm t_n,\]
\BT{or equivalently it follows from \eqref{eq:scalar-residual-minimizer} that}{或者等价地, 由 \eqref{eq:scalar-residual-minimizer} 得}
\[\zeta_n=\frac{\bm t_n^{\mathrm H}A\bm t_n}{(A\bm t_n)^{\mathrm H}(A\bm t_n)}.\]
\BT{Since $\eta_n=0$, we obtain}{由于 $\eta_n=0$, 得到}
\[\bm u_n=\zeta_nA\bm p_n,\qquad\bm z_n=\zeta_n\bm t_n.\]
\BT{The algorithm of the BiCGSTAB method is described in Algorithm\ \ref{alg:bicgstab}.}{BiCGSTAB 方法的算法见 Algorithm\ \ref{alg:bicgstab}.}
\BT{The BiCGSTAB method is short for the BiCG stabilized method. Below is an explanation of the meaning of the stabilization. Since $\eta_n=0$, the recurrence relations in \eqref{eq:product-coupled-h} and \eqref{eq:product-coupled-g} reduce to}{BiCGSTAB 方法是 BiCG 稳定化方法的缩写. 下面说明稳定化的含义. 由于 $\eta_n=0$, \eqref{eq:product-coupled-h} 和 \eqref{eq:product-coupled-g} 中的递推关系化为}
\[H_n(\lambda)=H_{n-1}(\lambda)-\zeta_{n-1}\lambda H_{n-1}(\lambda)=(1-\zeta_{n-1}\lambda)H_{n-1}(\lambda).\]
\BT{This means that the residual vector of the BiCGSTAB method is given by}{这意味着 BiCGSTAB 方法的残差向量为}
\[\bm r_{n+1}=H_{n+1}(A)R_{n+1}(A)\bm r_0=(I-\zeta_nA)\bm r_n^{\mathrm{BiCG}}.\]
\BT{This implies that since $\zeta_n$ is chosen such that $\|\bm r_{n+1}\|$ is minimized, the residual vector of the BiCG method is expected to be stabilized, i.e., the BiCGSTAB method is expected to have smoother convergence behavior than the BiCG method.}{这表明, 由于选择 $\zeta_n$ 使 $\|\bm r_{n+1}\|$ 达到极小值, 可以预期 BiCG 方法的残差向量得到稳定化, 即可以预期 BiCGSTAB 方法比 BiCG 方法具有更平滑的收敛行为.}
% block: bicgstab-algorithm page:140 type:ordinary-paragraph
\begin{algorithm}[H]\caption{\BT{The BiCGSTAB method}{BiCGSTAB 方法}}\label{alg:bicgstab}\begin{algorithmic}[1]
\Statex\textbf{Input:} $\bm x_0\in\mathbb C^N,\ \beta_{-1}=0,\ \bm p_{-1}=\bm0,\ \bm r_0=\bm b-A\bm x_0$
\Statex\textbf{Input:} \BT{Choose $\bm r_0^{*}\in\mathbb C^N$, e.g., $\bm r_0^{*}=\bm r_0$}{选择 $\bm r_0^{*}\in\mathbb C^N$, 例如 $\bm r_0^{*}=\bm r_0$}
\Statex\textbf{Output:} $\bm x_n$
\For{$n=0,1,\ldots$}
\State $\bm p_n=\bm r_n+\beta_{n-1}(\bm p_{n-1}-\zeta_{n-1}A\bm p_{n-1})$
\State $\alpha_n=\frac{(\bm r_0^{*},\bm r_n)}{(\bm r_0^{*},A\bm p_n)}$
\State $\bm t_n=\bm r_n-\alpha_nA\bm p_n$
\State $\zeta_n=\frac{(A\bm t_n,\bm t_n)}{(A\bm t_n,A\bm t_n)}$
\State $\bm x_{n+1}=\bm x_n+\alpha_n\bm p_n+\zeta_n\bm t_n$
\State $\bm r_{n+1}=\bm t_n-\zeta_nA\bm t_n$
\State $\beta_n=\frac{\alpha_n}{\zeta_n}\times\frac{(\bm r_0^{*},\bm r_{n+1})}{(\bm r_0^{*},\bm r_n)}$
\EndFor\end{algorithmic}\end{algorithm}
% block: gpbicg-parameter-choice pages:140,141 type:subsection
\paragraph{\BT{The Choice for the GPBiCG Method}{GPBiCG 方法的选择}}\label{subsec:gpbicg-parameter-choice}\mbox{}\par
\BT{For the GPBiCG method, two parameters $\eta_n$ and $\zeta_n$ are chosen by minimizing the 2-norm of the residual vector $\|\bm r_{n+1}\|=\|\bm t_n-\eta_n\bm y_n-\zeta_nA\bm t_n\|$. The following orthogonal condition enables us to achieve the minimization:}{对于 GPBiCG 方法, 通过极小化残差向量的 2-范数 $\|\bm r_{n+1}\|=\|\bm t_n-\eta_n\bm y_n-\zeta_nA\bm t_n\|$ 来选择两个参数 $\eta_n$ 和 $\zeta_n$. 下列正交条件可实现这一极小化:}
\[\bm r_{n+1}\perp A\bm t_n,\qquad\bm r_{n+1}\perp\bm y_n,\]
\BT{or equivalently it follows from \eqref{eq:least-squares-decomposition} and \eqref{eq:least-squares-minimizer} with setting $\bm d=\bm t_n$, $M=[\bm y_n,A\bm t_n]$, $\bm x=[\eta_n,\zeta_n]^{\top}$ that we have}{或者等价地, 在 \eqref{eq:least-squares-decomposition} 和 \eqref{eq:least-squares-minimizer} 中置 $\bm d=\bm t_n$, $M=[\bm y_n,A\bm t_n]$, $\bm x=[\eta_n,\zeta_n]^{\top}$, 得到}
\[\begin{aligned}\zeta_n&=\frac{(\bm y_n,\bm y_n)(A\bm t_n,\bm t_n)-(\bm y_n,\bm t_n)(A\bm t_n,\bm y_n)}{(A\bm t_n,A\bm t_n)(\bm y_n,\bm y_n)-(\bm y_n,A\bm t_n)(A\bm t_n,\bm y_n)},\\\eta_n&=\frac{(A\bm t_n,A\bm t_n)(\bm y_n,\bm t_n)-(\bm y_n,A\bm t_n)(A\bm t_n,\bm t_n)}{(A\bm t_n,A\bm t_n)(\bm y_n,\bm y_n)-(\bm y_n,A\bm t_n)(A\bm t_n,\bm y_n)}.\end{aligned}\]
\BT{The algorithm of the GPBiCG method is described in Algorithm\ \ref{alg:gpbicg}. Note the GPBiCG method uses two parameters $\zeta_n$ and $\eta_n$ for the minimization of the residual 2-norm, and the BiCGSTAB method uses one parameter $\zeta_n$ for the minimization of the residual 2-norm and $\eta_n=0$. From this, the GPBiCG method is expected to give smoother convergence behavior than the BiCGSTAB method.}{GPBiCG 方法的算法见 Algorithm\ \ref{alg:gpbicg}. 注意, GPBiCG 方法使用两个参数 $\zeta_n$ 和 $\eta_n$ 来极小化残差 2-范数, 而 BiCGSTAB 方法使用一个参数 $\zeta_n$ 来极小化残差 2-范数, 且 $\eta_n=0$. 由此, 可以预期 GPBiCG 方法比 BiCGSTAB 方法表现出更平滑的收敛行为.}
% block: gpbicg-algorithm page:141 type:ordinary-paragraph
\begin{algorithm}[H]\caption{\BT{The GPBiCG method}{GPBiCG 方法}}\label{alg:gpbicg}\begin{algorithmic}[1]
\Statex\textbf{Input:} $\bm x_0\in\mathbb C^N,\ \beta_{-1}=0,\ \bm p_{-1}=\bm t_{-1}=\bm u_{-1}=\bm w_{-1}=\bm z_{-1}=\bm0,\ \bm r_0=\bm b-A\bm x_0$
\Statex\textbf{Input:} \BT{Choose $\bm r_0^{*}\in\mathbb C^N$, e.g., $\bm r_0^{*}=\bm r_0$}{选择 $\bm r_0^{*}\in\mathbb C^N$, 例如 $\bm r_0^{*}=\bm r_0$}
\Statex\textbf{Output:} $\bm x_n$
\For{$n=0,1,\ldots$}
\State $\bm p_n=\bm r_n+\beta_{n-1}(\bm p_{n-1}-\bm u_{n-1})$
\State $\alpha_n=\frac{(\bm r_0^{*},\bm r_n)}{(\bm r_0^{*},A\bm p_n)}$
\State $\bm y_n=\bm t_{n-1}-\bm r_n-\alpha_n\bm w_{n-1}+\alpha_nA\bm p_n$
\State $\bm t_n=\bm r_n-\alpha_nA\bm p_n$
\State $\zeta_n=\frac{(\bm y_n,\bm y_n)(A\bm t_n,\bm t_n)-(\bm y_n,\bm t_n)(A\bm t_n,\bm y_n)}{(A\bm t_n,A\bm t_n)(\bm y_n,\bm y_n)-(\bm y_n,A\bm t_n)(A\bm t_n,\bm y_n)}$
\State $\eta_n=\frac{(A\bm t_n,A\bm t_n)(\bm y_n,\bm t_n)-(\bm y_n,A\bm t_n)(A\bm t_n,\bm t_n)}{(A\bm t_n,A\bm t_n)(\bm y_n,\bm y_n)-(\bm y_n,A\bm t_n)(A\bm t_n,\bm y_n)}$
\If{$n=0$}
\State $\zeta_n=\frac{(A\bm t_n,\bm t_n)}{(A\bm t_n,A\bm t_n)},\quad\eta_n=0$
\EndIf
\State $\bm u_n=\zeta_nA\bm p_n+\eta_n(\bm t_{n-1}-\bm r_n+\beta_{n-1}\bm u_{n-1})$
\State $\bm z_n=\zeta_n\bm r_n+\eta_n\bm z_{n-1}-\alpha_n\bm u_n$
\State $\bm x_{n+1}=\bm x_n+\alpha_n\bm p_n+\bm z_n$
\State $\bm r_{n+1}=\bm t_n-\eta_n\bm y_n-\zeta_nA\bm t_n$
\State $\beta_n=\frac{\alpha_n}{\zeta_n}\times\frac{(\bm r_0^{*},\bm r_{n+1})}{(\bm r_0^{*},\bm r_n)}$
\State $\bm w_n=A\bm t_n+\beta_nA\bm p_n$
\EndFor\end{algorithmic}\end{algorithm}
\BT{From Zhang's framework, the GPBiCG($m,\ell$) method was proposed by Fujino \BookCite{71}. A unified generalization of the GPBiCG method and the BiCGSTAB($\ell$) method was proposed by K. Aihara \BookCite{4}, which is referred to as the GPBiCGstab($L$) method.}{Fujino \BookCite{71} 从 Zhang 的框架提出 GPBiCG($m,\ell$) 方法. K. Aihara \BookCite{4} 提出了 GPBiCG 方法和 BiCGSTAB($\ell$) 方法的统一推广, 称为 GPBiCGstab($L$) 方法.}
% block: product-numerical-examples page:142 type:subsection
\paragraph{\BT{Numerical Examples}{数值例子}}\label{subsec:product-numerical-examples}\mbox{}\par
\BT{In this subsection, some illustrative numerical examples are shown. We consider nonsymmetric linear systems \eqref{eq:three-dimensional-difference-system} with parameters $N=20$, $a_1=a_2=a_3=1$, $b_1=b_2=b_3=1$, and two different parameters $c=300$ and $c=900$. The stopping criterion used for all the Krylov subspace methods is $\|\bm r_n\|/\|\bm b\|\le10^{-10}$.}{本小节给出一些说明性的数值例子. 考虑非对称线性方程组 \eqref{eq:three-dimensional-difference-system}, 参数为 $N=20$, $a_1=a_2=a_3=1$, $b_1=b_2=b_3=1$, 以及两个不同参数 $c=300$ 和 $c=900$. 所有 Krylov 子空间方法均采用停止准则 $\|\bm r_n\|/\|\bm b\|\le10^{-10}$.}
\BT{The convergence histories of the CGS method, the BiCGSTAB method, and the GPBiCG method are shown in Fig.\ \ref{fig:product-methods-moderate} for $c=300$. The horizontal axis is not the number of iterations but the number of matrix--vector multiplications. (The CGS method, the BiCGSTAB method, and the GPBiCG method require two matrix--vector multiplications per iteration.)}{当 $c=300$ 时, CGS 方法, BiCGSTAB 方法和 GPBiCG 方法的收敛历程见 Figure\ \ref{fig:product-methods-moderate}. 横轴不是迭代次数, 而是矩阵向量乘法的次数. (CGS 方法, BiCGSTAB 方法和 GPBiCG 方法每次迭代均需两次矩阵向量乘法.)}
\BT{From Fig.\ \ref{fig:product-methods-moderate}, the GPBiCG method converges faster than BiCGSTAB method, and the CGS method is the best in terms of the number of matrix--vector multiplications, or equivalently the number of iterations. On the other hand, the CGS method shows irregular convergence behavior, which may lead to a loss of accuracy. The GPBiCG method and the BiCGSTAB method show relatively smooth convergence behavior, which may be caused by local minimization of the residual 2-norms. In terms of accuracy, the $\log_{10}$ of the true relative residual 2-norms for the CGS method, the BiCGSTAB method, and the GPBiCG method are $-9.79$, $-10.50$, $-10.17$ respectively. From which, all the methods produced sufficiently accurate approximate solutions.}{由 Figure\ \ref{fig:product-methods-moderate}, GPBiCG 方法比 BiCGSTAB 方法收敛更快, 而从矩阵向量乘法次数, 或等价的迭代次数看, CGS 方法最好. 另一方面, CGS 方法表现出不规则的收敛行为, 这可能导致精度损失. GPBiCG 方法和 BiCGSTAB 方法表现出相对平滑的收敛行为, 这可能由残差 2-范数的局部极小化所致. 从精度看, CGS 方法, BiCGSTAB 方法和 GPBiCG 方法的真实相对残差 2-范数的 $\log_{10}$ 分别为 $-9.79$, $-10.50$, $-10.17$. 由此, 所有方法都产生了足够精确的近似解.}
% block: product-residual-gap pages:142,143 type:ordinary-paragraph
\BT{Next, the convergence histories for $c=900$ are shown in Fig.\ \ref{fig:product-methods-difficult}. We see from Fig.\ \ref{fig:product-methods-difficult} that the CGS method shows much irregular convergence behavior and the slowest convergence, and the GPBiCG method is the best in terms of the number of matrix--vector multiplications. In terms of accuracy, the $\log_{10}$ of the true relative residual 2-norms for the CGS method, the BiCGSTAB method, and the GPBiCG method are $-6.37$, $-10.03$, $-9.79$ respectively. This means that the CGS method suffers from a loss of accuracy of about four digits. This loss of accuracy is known as a residual gap. Therefore, it is strongly recommended to check the true relative residual 2-norm, i.e., $\|\bm b-A\bm x_n\|/\|\bm b\|$ (not $\|\bm r_n\|/\|\bm b\|$!), after satisfying the stopping criterion.}{接着, $c=900$ 时的收敛历程见 Figure\ \ref{fig:product-methods-difficult}. 由 Figure\ \ref{fig:product-methods-difficult} 可见, CGS 方法表现出很不规则的收敛行为, 收敛也最慢; 从矩阵向量乘法次数看, GPBiCG 方法最好. 从精度看, CGS 方法, BiCGSTAB 方法和 GPBiCG 方法的真实相对残差 2-范数的 $\log_{10}$ 分别为 $-6.37$, $-10.03$, $-9.79$. 这意味着 CGS 方法损失了约四位精度. 这种精度损失称为残差间隙. 因此, 强烈建议在满足停止准则后检查真实相对残差 2-范数, 即 $\|\bm b-A\bm x_n\|/\|\bm b\|$ (不是 $\|\bm r_n\|/\|\bm b\|$!).}
% block: product-moderate-figure page:142 type:figure
\begin{figure}[H]\centering\includegraphics[width=.87\linewidth]{figures/product-methods-moderate.png}
\caption{\BT{Convergence histories of CGS, BiCGSTAB, and GPBiCG for $N=20,c=300$}{当 $N=20,c=300$ 时 CGS, BiCGSTAB 和 GPBiCG 的收敛历程}}\label{fig:product-methods-moderate}\end{figure}
% block: product-difficult-figure page:143 type:figure
\begin{figure}[H]\centering\includegraphics[width=.87\linewidth]{figures/product-methods-difficult.png}
\caption{\BT{Convergence histories of CGS, BiCGSTAB, and GPBiCG for $N=20,c=900$}{当 $N=20,c=900$ 时 CGS, BiCGSTAB 和 GPBiCG 的收敛历程}}\label{fig:product-methods-difficult}\end{figure}
% block: product-related-methods pages:143,144 type:subsection
\paragraph{\BT{Other Related Methods}{其他相关方法}}\label{subsec:product-related-methods}\mbox{}\par
\BT{The QMR method in Sect.\ \ref{subsec:qmr-method} requires the information of $A^{\mathrm H}$, as well as the BiCG method. Similar to the idea of the CGS method, R. W. Freund proposed a squared variant of the QMR method referred to as the TFQMR (transpose-free QMR) method \BookCite{63}, which tends to give smooth convergence behavior. Based on the BiCGSTAB method and the idea of the QMR method, the QMRCGSTAB method was proposed in \BookCite{34}.}{Section\ \ref{subsec:qmr-method} 中的 QMR 方法与 BiCG 方法一样需要 $A^{\mathrm H}$ 的信息. 与 CGS 方法的思想类似, R. W. Freund 提出了 QMR 方法的平方变型, 称为 TFQMR (无转置 QMR) 方法 \BookCite{63}, 它往往表现出平滑的收敛行为. 文献 \BookCite{34} 基于 BiCGSTAB 方法和 QMR 方法的思想提出 QMRCGSTAB 方法.}
\BT{The BiCR method in Sect.\ \ref{subsec:bicr-method} may give a smoother convergence behavior than the BiCG method when the coefficient matrix $A$ is close to Hermitian, since the BiCR method with the choice $\bm r_0^{*}=\bm r_0$ becomes the CR method when $A$ is Hermitian and the residual 2-norm of the CR method monotonically decreases. Thus it is natural to replace the product-type Krylov subspace methods based on the BiCG method with those based on the BiCR method. In fact, the product-type methods based on the BiCR method can be constructed, and the resulting algorithms (CRS, BiCRSTAB, GPBiCR) are derived in \BookCite{165}. K. Abe and G. L. G. Sleijpen independently proposed product-type methods based on the BiCR method, see \BookCite{2}.}{当系数矩阵 $A$ 接近 Hermitian 矩阵时, Section\ \ref{subsec:bicr-method} 中的 BiCR 方法可能比 BiCG 方法表现出更平滑的收敛行为, 因为当 $A$ 为 Hermitian 矩阵时, 选择 $\bm r_0^{*}=\bm r_0$ 的 BiCR 方法变为 CR 方法, 而 CR 方法的残差 2-范数单调下降. 因此, 自然可以将基于 BiCG 方法的乘积型 Krylov 子空间方法替换为基于 BiCR 方法的乘积型方法. 事实上, 可构造基于 BiCR 方法的乘积型方法, 文献 \BookCite{165} 导出了所得算法 (CRS, BiCRSTAB, GPBiCR). K. Abe 和 G. L. G. Sleijpen 独立提出了基于 BiCR 方法的乘积型方法, 见 \BookCite{2}.}
\BT{Y.-F. Jing et al. proposed an approach similar to the BiCR method called the BiCOR method \BookCite{114}, and the product-type variants were proposed in \BookCite{32}.}{Y.-F. Jing 等提出了类似于 BiCR 方法的一种途径, 称为 BiCOR 方法 \BookCite{114}, 其乘积型变型见 \BookCite{32}.}
% block: idr-heading page:144 type:subsection
\subsubsection{\BT{Induced Dimension Reduction (IDR($s$)) Method}{诱导降维 (IDR($s$)) 方法}}\label{subsec:idr-method}
\BT{The Induced Dimension Reduction (IDR) method was developed by Wesseling and Sonneveld around 1979 \BookCite{205}. The IDR method is mathematically equivalent to the BiCGSTAB method, and the BiCGSTAB method can be regarded as a stabilized variant of the IDR method. It is written in the acknowledgements in \BookCite{195} that Sonneveld suggested that van der Vorst reconsiders Sonneveld's IDR.}{诱导降维 (IDR) 方法由 Wesseling 和 Sonneveld 于 1979 年前后建立 \BookCite{205}. IDR 方法在数学上等价于 BiCGSTAB 方法, 而 BiCGSTAB 方法可视为 IDR 方法的稳定化变型. 文献 \BookCite{195} 的致谢中写道, Sonneveld 建议 van der Vorst 重新考虑 Sonneveld 的 IDR.}
\BT{Though the IDR method was not given further attention, about 30 years later, Sonneveld and van Gijzen proposed an innovative algorithm called the IDR($s$) method \BookCite{174}. In this section, the theory and the derivation of the IDR($s$) method are described. The derivation in this section looks to be redundant, but it makes it easier to understand the principle of the IDR($s$) method.}{虽然 IDR 方法未受到进一步关注, 但大约 30 年后, Sonneveld 和 van Gijzen 提出了一种称为 IDR($s$) 方法的创新算法 \BookCite{174}. 本节介绍 IDR($s$) 方法的理论与推导. 本节推导看似冗长, 但使 IDR($s$) 方法的原理更易理解.}
% block: idr-theorem page:144 type:theorem
\par\Needspace{14\baselineskip}
\begin{Theorem}[\BT{The IDR theorem}{IDR 定理} \BookCite{174}, \BookCite{205}, p.550]\label{thm:idr-dimension-reduction}
\BT{Let $A$ be any matrix in $\mathbb C^{N\times N}$, let $\bm v_0$ be any nonzero vector in $\mathbb C^N$, and let $\mathcal G_0$ be the complete Krylov space, i.e., $\mathcal G_0:=\mathcal K_N(A,\bm v_0)$. Let $\mathcal S$ denote any (proper) subspace of $\mathbb C^N$ such that $\mathcal S$ and $\mathcal G_0$ do not share a nontrivial invariant subspace,}{设 $A$ 为 $\mathbb C^{N\times N}$ 中的任意矩阵, $\bm v_0$ 为 $\mathbb C^N$ 中的任意非零向量, $\mathcal G_0$ 为完全 Krylov 空间, 即 $\mathcal G_0:=\mathcal K_N(A,\bm v_0)$. 设 $\mathcal S$ 表示 $\mathbb C^N$ 的任意 (真) 子空间, 使 $\mathcal S$ 和 $\mathcal G_0$ 不具有共同的非平凡不变子空间,}\footnote{\BT{This means that $\mathcal S\cap\mathcal G_0$ does not contain any eigenvector of $A$, and the trivial invariant subspace is $\{\bm0\}$.}{这意味着 $\mathcal S\cap\mathcal G_0$ 不包含 $A$ 的任何特征向量, 平凡不变子空间为 $\{\bm0\}$.}}
\BT{and define the sequences $\mathcal G_j$, $j=1,2,\ldots$ as}{并定义序列 $\mathcal G_j$, $j=1,2,\ldots$ 为}
\begin{equation}\label{eq:idr-subspace-recurrence}\mathcal G_j=(I-\omega_jA)(\mathcal G_{j-1}\cap\mathcal S),\end{equation}
\BT{where the $\omega_j$'s are nonzero scalars. Then:}{其中 $\omega_j$ 为非零标量. 则:}
\begin{enumerate}[label=(\arabic*)]
\item $\mathcal G_j\subset\mathcal G_{j-1}$ \BT{for all $j>0$.}{对所有 $j>0$ 成立.}
\item $\mathcal G_j=\{\bm0\}$ \BT{for some $j\le N$.}{对某个 $j\le N$ 成立.}
\end{enumerate}\end{Theorem}
% block: idr-theorem-proof-nesting pages:144,145 type:proof
\begin{Proof}
\BT{The proof is based on \BookCite{174}, Theorem 2.1.}{本证明基于 \BookCite{174}, Theorem 2.1.}
\BT{The first statement is shown by induction. First, we show $\mathcal G_1\subset\mathcal G_0$. From $\mathcal G_0=\mathcal K_N(A,\bm v_0)$, it follows that $(I-\omega_1A)\mathcal G_0\subset\mathcal G_0$, because for a given $\bm v\in(I-\omega_1A)\mathcal G_0$ we have}{第一个断言用归纳法证明. 首先证明 $\mathcal G_1\subset\mathcal G_0$. 由 $\mathcal G_0=\mathcal K_N(A,\bm v_0)$ 得 $(I-\omega_1A)\mathcal G_0\subset\mathcal G_0$, 因为对给定的 $\bm v\in(I-\omega_1A)\mathcal G_0$, 有}
\[\begin{aligned}\bm v&=(I-\omega_1A)(c_0\bm v_0+c_1A\bm v_0+\cdots+c_{N-1}A^{N-1}\bm v_0)\\&=c_0(I-\omega_1A)\bm v_0+c_1(I-\omega_1A)A\bm v_0+\cdots+c_{N-1}(I-\omega_1A)A^{N-1}\bm v_0\\&=d_0\bm v_0+d_1A\bm v_0+\cdots+d_{N-1}A^{N-1}\bm v_0+d_NA^N\bm v_0\\&=\widetilde d_0\bm v_0+\widetilde d_1A\bm v_0+\cdots+\widetilde d_{N-1}A^{N-1}\bm v_0\\&\in\mathcal G_0.\end{aligned}\]
\BT{Thus we have $\mathcal G_1\subset\mathcal G_0$ because}{因此有 $\mathcal G_1\subset\mathcal G_0$, 因为}
\[\mathcal G_1=(I-\omega_1A)(\mathcal G_0\cap\mathcal S)\subset(I-\omega_1A)\mathcal G_0\subset\mathcal G_0.\]
\BT{Next we show $\mathcal G_{j+1}\subset\mathcal G_j$ (i.e., $\bm x\in\mathcal G_{j+1}\Rightarrow\bm x\in\mathcal G_j$) for all $j>0$ using the assumption $\mathcal G_j\subset\mathcal G_{j-1}$. Let $\bm x\in\mathcal G_{j+1}$. Then $\bm x$ can be written as $\bm x=(I-\omega_{j+1}A)\bm y$ for some $\bm y\in\mathcal G_j\cap\mathcal S$. From the assumption $\mathcal G_j\subset\mathcal G_{j-1}$, it follows that $\bm y\in\mathcal G_j\cap\mathcal S\subset\mathcal G_{j-1}\cap\mathcal S$, and thus $(I-\omega_jA)\bm y\in\mathcal G_j$. Since $\bm y\in\mathcal G_j$ and $\bm y-\omega_jA\bm y\in\mathcal G_j$, we have $A\bm y\in\mathcal G_j$, and thus $(I-\omega_{j+1}A)\bm y=\bm x\in\mathcal G_j$.}{接着, 在假设 $\mathcal G_j\subset\mathcal G_{j-1}$ 下, 对所有 $j>0$ 证明 $\mathcal G_{j+1}\subset\mathcal G_j$ (即 $\bm x\in\mathcal G_{j+1}\Rightarrow\bm x\in\mathcal G_j$). 设 $\bm x\in\mathcal G_{j+1}$. 则对某个 $\bm y\in\mathcal G_j\cap\mathcal S$, $\bm x$ 可写为 $\bm x=(I-\omega_{j+1}A)\bm y$. 由假设 $\mathcal G_j\subset\mathcal G_{j-1}$, 得 $\bm y\in\mathcal G_j\cap\mathcal S\subset\mathcal G_{j-1}\cap\mathcal S$, 因而 $(I-\omega_jA)\bm y\in\mathcal G_j$. 由于 $\bm y\in\mathcal G_j$ 且 $\bm y-\omega_jA\bm y\in\mathcal G_j$, 有 $A\bm y\in\mathcal G_j$, 从而 $(I-\omega_{j+1}A)\bm y=\bm x\in\mathcal G_j$.}
% block: idr-theorem-proof-dimension page:145 type:proof
\BT{The second statement is shown next. The property $\mathcal G_{j+1}\subset\mathcal G_j$ means that $\dim(\mathcal G_{j+1})<\dim(\mathcal G_j)$ or $\dim(\mathcal G_{j+1})=\dim(\mathcal G_j)$.}{下面证明第二个断言. 性质 $\mathcal G_{j+1}\subset\mathcal G_j$ 意味着 $\dim(\mathcal G_{j+1})<\dim(\mathcal G_j)$ 或 $\dim(\mathcal G_{j+1})=\dim(\mathcal G_j)$.}
\BT{If $\mathcal G_j\cap\mathcal S\ne\mathcal G_j$, then $\dim(\mathcal G_j\cap\mathcal S)<\dim(\mathcal G_j)$ and thus we have $\dim(\mathcal G_{j+1})<\dim(\mathcal G_j)$ because $\dim(\mathcal G_{j+1})=\dim((I-\omega_jA)(\mathcal G_j\cap\mathcal S))\le\dim(\mathcal G_j\cap\mathcal S)<\dim(\mathcal G_j)$.}{若 $\mathcal G_j\cap\mathcal S\ne\mathcal G_j$, 则 $\dim(\mathcal G_j\cap\mathcal S)<\dim(\mathcal G_j)$, 从而有 $\dim(\mathcal G_{j+1})<\dim(\mathcal G_j)$, 因为 $\dim(\mathcal G_{j+1})=\dim((I-\omega_jA)(\mathcal G_j\cap\mathcal S))\le\dim(\mathcal G_j\cap\mathcal S)<\dim(\mathcal G_j)$.}
\BT{On the other hand, if $\mathcal G_j\cap\mathcal S=\mathcal G_j$, then $\mathcal G_j\subset\mathcal S$, and we have $\mathcal G_{j+1}=(I-\omega_{j+1}A)(\mathcal G_j\cap\mathcal S)=(I-\omega_{j+1}A)\mathcal G_j\subset\mathcal G_j$. (The last inclusion follows from the first statement $\mathcal G_{j+1}\subset\mathcal G_j$.) This implies $A\mathcal G_j\subset\mathcal G_j$, which means that $\mathcal G_j$ is an invariant subspace of $A$. Recalling $\mathcal G_j\subset\mathcal S$ and $\mathcal G_j\subset\mathcal G_0$ leads to $\mathcal G_j\subset\mathcal G_0\cap\mathcal S$. From the assumption that $\mathcal G_0\cap\mathcal S$ do not share a nontrivial invariant subspace of $A$, the subspace $\mathcal G_j$ is the trivial invariant subspace of $A$, i.e., $\mathcal G_j=\{\bm0\}$.}{另一方面, 若 $\mathcal G_j\cap\mathcal S=\mathcal G_j$, 则 $\mathcal G_j\subset\mathcal S$, 且有 $\mathcal G_{j+1}=(I-\omega_{j+1}A)(\mathcal G_j\cap\mathcal S)=(I-\omega_{j+1}A)\mathcal G_j\subset\mathcal G_j$. (最后一个包含关系来自第一个断言 $\mathcal G_{j+1}\subset\mathcal G_j$.) 这表明 $A\mathcal G_j\subset\mathcal G_j$, 即 $\mathcal G_j$ 为 $A$ 的不变子空间. 回顾 $\mathcal G_j\subset\mathcal S$ 和 $\mathcal G_j\subset\mathcal G_0$, 得 $\mathcal G_j\subset\mathcal G_0\cap\mathcal S$. 由 $\mathcal G_0\cap\mathcal S$ 不具有共同的 $A$ 的非平凡不变子空间这一假设, 子空间 $\mathcal G_j$ 为 $A$ 的平凡不变子空间, 即 $\mathcal G_j=\{\bm0\}$.}
\BT{Therefore either $\dim(\mathcal G_{j+1})<\dim(\mathcal G_j)$ or $\dim(\mathcal G_j)=0$ occurs, which indicates that $\dim(\mathcal G_k)=0$, i.e., $\mathcal G_k=\{\bm0\}$ for some $k\le N$.}{因此, 或者 $\dim(\mathcal G_{j+1})<\dim(\mathcal G_j)$, 或者 $\dim(\mathcal G_j)=0$, 这表明对某个 $k\le N$ 有 $\dim(\mathcal G_k)=0$, 即 $\mathcal G_k=\{\bm0\}$.}
\par\nopagebreak[4]
\end{Proof}
% block: idr-dimension-corollary page:145 type:corollary
\par\Needspace{3\baselineskip}
\begin{Corollary}\label{cor:idr-dimension-step}
\BT{Let $\mathcal S$ be a subspace of $\mathbb C^N$ with $\dim(\mathcal S)=N-s$ for a given $s\ge1$. If $\mathcal G_{i-1}+\mathcal S=\mathbb C^N$ and $(I-\omega_jA)$ is nonsingular, then}{对给定的 $s\ge1$, 设 $\mathcal S$ 为 $\mathbb C^N$ 的子空间且 $\dim(\mathcal S)=N-s$. 若 $\mathcal G_{i-1}+\mathcal S=\mathbb C^N$ 且 $(I-\omega_jA)$ 非奇异, 则}
\begin{equation}\label{eq:idr-dimension-step}\dim(\mathcal G_j)=\dim(\mathcal G_{j-1})-s.\end{equation}
\end{Corollary}
% block: idr-dimension-corollary-proof page:145 type:proof
\begin{Proof}
\BT{Let $W_1,W_2$ be subspaces of $\mathbb C^N$ (or a general vector space $V$). Then it is well known in linear algebra that $\dim(W_1)+\dim(W_2)=\dim(W_1+W_2)+\dim(W_1\cap W_2)$. Let $W_1=\mathcal G_{i-1}$, $W_2=\mathcal S$. Then, from the assumption $\mathcal G_{i-1}+\mathcal S=\mathbb C^N$, we have $\dim(\mathcal G_{i-1}+\mathcal S)=N$. Thus $\dim(\mathcal G_i)=\dim((I-\omega_jA)(\mathcal G_{i-1}\cap\mathcal S))=\dim(\mathcal G_{i-1}\cap\mathcal S)=\dim(\mathcal G_{i-1})+\dim(\mathcal S)-N=\dim(\mathcal G_{i-1})-s$.}{设 $W_1,W_2$ 为 $\mathbb C^N$ (或一般向量空间 $V$) 的子空间. 线性代数中众所周知, $\dim(W_1)+\dim(W_2)=\dim(W_1+W_2)+\dim(W_1\cap W_2)$. 令 $W_1=\mathcal G_{i-1}$, $W_2=\mathcal S$. 则由假设 $\mathcal G_{i-1}+\mathcal S=\mathbb C^N$, 有 $\dim(\mathcal G_{i-1}+\mathcal S)=N$. 因而 $\dim(\mathcal G_i)=\dim((I-\omega_jA)(\mathcal G_{i-1}\cap\mathcal S))=\dim(\mathcal G_{i-1}\cap\mathcal S)=\dim(\mathcal G_{i-1})+\dim(\mathcal S)-N=\dim(\mathcal G_{i-1})-s$.}
\par\nopagebreak[4]
\end{Proof}
% block: idr-dimension-remark page:145 type:remark
\par\Needspace{3\baselineskip}
\begin{Remark}\label{rem:idr-dimension-count}
\BT{From Corollary\ \ref{cor:idr-dimension-step}, if $\mathcal G_{i-1}+\mathcal S=\mathbb C^N$ for $i=1,2,\ldots,n$, then $\dim(\mathcal G_n)=\dim(\mathcal G_0)-ns$. Thus, if $N/s$ is an integer and $\dim(\mathcal G_0)=N$, we have $\dim(\mathcal G_{N/s})=0$, i.e., $\mathcal G_{N/s}=\{\bm0\}$. This remark will be cited in Remark\ \ref{rem:idr-matvec-bound} for the maximum required number of matrix--vector multiplications of the IDR($s$) method.}{由 Corollary\ \ref{cor:idr-dimension-step}, 若对 $i=1,2,\ldots,n$ 有 $\mathcal G_{i-1}+\mathcal S=\mathbb C^N$, 则 $\dim(\mathcal G_n)=\dim(\mathcal G_0)-ns$. 因而, 若 $N/s$ 为整数且 $\dim(\mathcal G_0)=N$, 则有 $\dim(\mathcal G_{N/s})=0$, 即 $\mathcal G_{N/s}=\{\bm0\}$. Remark\ \ref{rem:idr-matvec-bound} 将引用本注以给出 IDR($s$) 方法所需矩阵向量乘法次数的最大值.}
\end{Remark}
% block: idr-derivation-heading pages:145,146 type:subsection
\paragraph{\BT{The Derivation of the IDR($s$) Method}{IDR($s$) 方法的推导}}\label{subsec:idr-derivation}\mbox{}\par
\BT{The IDR($s$) method with a given number $s\ge1$ generates residual vectors $\bm r_0,\bm r_1,\ldots$ such that:}{对于给定数 $s\ge1$, IDR($s$) 方法生成满足下列条件的残差向量 $\bm r_0,\bm r_1,\ldots$:}
\begin{enumerate}
\item $\bm r_{j(s+1)},\bm r_{j(s+1)+1},\ldots,\bm r_{j(s+1)+s}\in\mathcal G_j,\quad j=0,1,2,\ldots$;
\item\BT{The subspace $\mathcal S$ in \eqref{eq:idr-subspace-recurrence} is chosen so that $\dim\mathcal S=N-s$.}{选择 \eqref{eq:idr-subspace-recurrence} 中的子空间 $\mathcal S$, 使 $\dim\mathcal S=N-s$.}
\end{enumerate}
\BT{For simplicity, we consider the case $s=2$, i.e., the IDR(2) method. In this case, the IDR(2) method produces the following residual vectors:}{为简单起见, 考虑 $s=2$ 的情形, 即 IDR(2) 方法. 在此情形下, IDR(2) 方法产生下列残差向量:}
\[\bm r_0,\bm r_1,\bm r_2\in\mathcal G_0,\quad\bm r_3,\bm r_4,\bm r_5\in\mathcal G_1,\quad\bm r_6,\bm r_7,\bm r_8\in\mathcal G_2,\ldots\]
% block: idr-two-shadow-space page:146 type:ordinary-paragraph
\BT{Here we explain how to produce $\bm r_3,\bm r_4,\bm r_5\in\mathcal G_1$ from the given residual vectors $\bm r_0,\bm r_1,\bm r_2\in\mathcal G_0$. Since $s=2$, it follows from Remark\ \ref{rem:idr-dimension-count} that we need to choose $\mathcal S$ such that $\dim(\mathcal S)=N-2$. The standard choice is}{这里说明如何从给定残差向量 $\bm r_0,\bm r_1,\bm r_2\in\mathcal G_0$ 产生 $\bm r_3,\bm r_4,\bm r_5\in\mathcal G_1$. 由于 $s=2$, 由 Remark\ \ref{rem:idr-dimension-count} 可知, 需选择 $\mathcal S$ 使 $\dim(\mathcal S)=N-2$. 标准选择为}
\begin{equation}\label{eq:idr-shadow-complement}\mathcal S=\operatorname{span}\{\bm s_1,\bm s_2\}^{\perp}\end{equation}
\BT{with arbitrary linearly independent vectors $\bm s_1$ and $\bm s_2$, where the symbol $\perp$ is the orthogonal complement.}{其中 $\bm s_1$ 和 $\bm s_2$ 为任意线性无关向量, 符号 $\perp$ 表示正交补.}
\BT{Now, let $S=[\bm s_1,\bm s_2]\in\mathbb C^{N\times2}$ with linearly independent vectors $\bm s_1$ and $\bm s_2$. Then for a vector $\bm v\in\mathbb C^N$, the following fact holds true:}{现在令 $S=[\bm s_1,\bm s_2]\in\mathbb C^{N\times2}$, 其中向量 $\bm s_1$ 和 $\bm s_2$ 线性无关. 则对于向量 $\bm v\in\mathbb C^N$, 下列事实成立:}
\begin{equation}\label{eq:idr-shadow-orthogonality}\bm v\perp\mathcal S\Leftrightarrow S^{\mathrm H}\bm v=\bm0.\end{equation}
\BT{In what follows, we use the following definitions:}{下面使用下列定义:}
\begin{equation}\label{eq:idr-residual-differences}\Delta\bm r_k:=\bm r_{k+1}-\bm r_k,\quad dR_k:=[\Delta\bm r_k,\Delta\bm r_{k+1}],\quad\bm c_k=(S^{\mathrm H}dR_k)^{-1}S^{\mathrm H}\bm r_{k+2}.\end{equation}
% block: idr-first-subspace-residuals pages:146,147 type:ordinary-paragraph
\BT{The following vector $\bm v_0$ lies in $\mathcal G_0\cap\mathcal S$:}{下列向量 $\bm v_0$ 位于 $\mathcal G_0\cap\mathcal S$ 中:}
\begin{equation}\label{eq:idr-first-projected-vector}\bm v_0=\bm r_2-dR_0\bm c_0\in\mathcal G_0\cap\mathcal S.\end{equation}
\BT{In fact, $\bm v_0\in\mathcal G_0$ because $\bm v_0\in\operatorname{span}\{\bm r_2,\Delta\bm r_0,\Delta\bm r_1\}$ with $\bm r_2,\Delta\bm r_0,\Delta\bm r_1\in\mathcal G_0$ from the assumption $\bm r_0,\bm r_1,\bm r_2\in\mathcal G_0$, and $\bm v_0\in\mathcal S$ because it is easy to see $S^{\mathrm H}\bm v_0=\bm0$ and the relation \eqref{eq:idr-shadow-orthogonality}.}{事实上, $\bm v_0\in\mathcal G_0$, 因为 $\bm v_0\in\operatorname{span}\{\bm r_2,\Delta\bm r_0,\Delta\bm r_1\}$, 而由假设 $\bm r_0,\bm r_1,\bm r_2\in\mathcal G_0$ 有 $\bm r_2,\Delta\bm r_0,\Delta\bm r_1\in\mathcal G_0$; 又有 $\bm v_0\in\mathcal S$, 因为容易看出 $S^{\mathrm H}\bm v_0=\bm0$, 并使用关系 \eqref{eq:idr-shadow-orthogonality}.}
\BT{From \eqref{eq:idr-subspace-recurrence}, we can produce $\bm r_3\in\mathcal G_1$ by}{由 \eqref{eq:idr-subspace-recurrence}, 可按下式产生 $\bm r_3\in\mathcal G_1$}
\begin{equation}\label{eq:idr-first-reduced-residual}\bm r_3=(I-\omega_0A)\bm v_0\in\mathcal G_1,\end{equation}
\BT{where $\omega_0$ is usually chosen so that $\|\bm r_3\|$ is minimized, i.e., $\omega_0=(A\bm v_0,\bm v_0)/(A\bm v_0,A\bm v_0)$.}{其中通常选择 $\omega_0$ 使 $\|\bm r_3\|$ 达到极小值, 即 $\omega_0=(A\bm v_0,\bm v_0)/(A\bm v_0,A\bm v_0)$.}
\BT{In order to produce $\bm r_4\in\mathcal G_1$, we need a vector $\bm v_1$ such that $\bm v_1\in\mathcal G_0\cap\mathcal S$. To this end, $\bm v_1$ is constructed by using $\bm r_3$ and $dR_1$ as follows:}{为产生 $\bm r_4\in\mathcal G_1$, 需要一个使 $\bm v_1\in\mathcal G_0\cap\mathcal S$ 的向量 $\bm v_1$. 为此, 利用 $\bm r_3$ 和 $dR_1$ 按下式构造 $\bm v_1$:}
\[\bm v_1=\bm r_3-dR_1\bm c_1\in\mathcal G_0\cap\mathcal S.\]
\BT{From a similar discussion to that above, we see $\bm v_1\in\mathcal G_0$ and $\bm v_1\in\mathcal S$. Thus, we have}{由与上面类似的讨论, 可见 $\bm v_1\in\mathcal G_0$ 且 $\bm v_1\in\mathcal S$. 因此有}
\[\bm r_4=(I-\omega_0A)\bm v_1\in\mathcal G_1.\]
\BT{Similarly,}{类似地,}
\[\bm v_2=\bm r_4-dR_2\bm c_2\in\mathcal G_0\cap\mathcal S,\quad\bm r_5=(I-\omega_0A)\bm v_2\in\mathcal G_1.\]
% block: idr-second-subspace-residuals page:147 type:ordinary-paragraph
\BT{Now, we obtained $\bm r_3,\bm r_4,\bm r_5\in\mathcal G_1$ from the given residual vectors $\bm r_0,\bm r_1,\bm r_2\in\mathcal G_0$. Similarly, the following algorithm produces $\bm r_6,\bm r_7,\bm r_8\in\mathcal G_2$ from $\bm r_3,\bm r_4,\bm r_5\in\mathcal G_1$:}{现在, 已从给定残差向量 $\bm r_0,\bm r_1,\bm r_2\in\mathcal G_0$ 得到 $\bm r_3,\bm r_4,\bm r_5\in\mathcal G_1$. 类似地, 下列算法从 $\bm r_3,\bm r_4,\bm r_5\in\mathcal G_1$ 产生 $\bm r_6,\bm r_7,\bm r_8\in\mathcal G_2$:}
\[\begin{aligned}\bm v_3&=\bm r_5-dR_3\bm c_3\in\mathcal G_1\cap\mathcal S,&\bm r_6&=(I-\omega_1A)\bm v_3\in\mathcal G_2,\\\bm v_4&=\bm r_6-dR_4\bm c_4\in\mathcal G_1\cap\mathcal S,&\bm r_7&=(I-\omega_1A)\bm v_4\in\mathcal G_2,\\\bm v_5&=\bm r_7-dR_5\bm c_5\in\mathcal G_1\cap\mathcal S,&\bm r_8&=(I-\omega_1A)\bm v_5\in\mathcal G_2,\end{aligned}\]
\BT{where $\omega_1$ is usually chosen so that $\|\bm r_6\|$ is minimized, i.e., $\omega_1=(A\bm v_3,\bm v_3)/(A\bm v_3,A\bm v_3)$.}{其中通常选择 $\omega_1$ 使 $\|\bm r_6\|$ 达到极小值, 即 $\omega_1=(A\bm v_3,\bm v_3)/(A\bm v_3,A\bm v_3)$.}
% block: idr-solution-extraction page:147 type:ordinary-paragraph
\BT{Now, we describe how to extract approximate solutions $\bm x_3,\bm x_4,\bm x_5$ from $\bm x_0,\bm x_1,\bm x_2$. In what follows, we use the following definition:}{现在说明如何从 $\bm x_0,\bm x_1,\bm x_2$ 提取近似解 $\bm x_3,\bm x_4,\bm x_5$. 下面使用下列定义:}
\[\Delta\bm x_k:=\bm x_{k+1}-\bm x_k,\qquad dX_k:=[\Delta\bm x_k,\Delta\bm x_{k+1}].\]
\BT{Then from \eqref{eq:idr-first-projected-vector} and \eqref{eq:idr-first-reduced-residual} $\bm x_3$ is obtained as follows:}{于是由 \eqref{eq:idr-first-projected-vector} 和 \eqref{eq:idr-first-reduced-residual}, 按下式得到 $\bm x_3$:}
\[\begin{aligned}\bm r_3&=(I-\omega_0A)\bm v_0=\bm r_2-dR_0\bm c_0-\omega_0A\bm v_0\\&\Leftrightarrow\bm b-A\bm x_3=\bm b-A\bm x_2-dR_0\bm c_0-\omega_0A\bm v_0\\&\Leftrightarrow\bm x_3=\bm x_2+A^{-1}dR_0\bm c_0+\omega_0\bm v_0\\&\Leftrightarrow\bm x_3=\bm x_2-dX_0\bm c_0+\omega_0\bm v_0.\end{aligned}\]
\BT{Similarly,}{类似地,}
\[\begin{aligned}\bm x_4&=\bm x_3-dX_1\bm c_1+\omega_0\bm v_1,\\\bm x_5&=\bm x_4-dX_2\bm c_2+\omega_0\bm v_2.\end{aligned}\]
\BT{$\bm x_6,\bm x_7,\bm x_8$ can be obtained as follows:}{$\bm x_6,\bm x_7,\bm x_8$ 可按下式得到:}
\[\begin{aligned}\bm x_6&=\bm x_5-dX_3\bm c_3+\omega_1\bm v_3,\\\bm x_7&=\bm x_7-dX_4\bm c_4+\omega_1\bm v_4,\\\bm x_8&=\bm x_8-dX_5\bm c_5+\omega_1\bm v_5.\end{aligned}\]
\BT{Finally, we describe how to determine initial residual vectors $\bm r_0,\bm r_1,\bm r_2\in\mathcal G_0$ and $\bm x_0,\bm x_1,\bm x_2$. These vectors are given as follows: we compute}{最后说明如何确定初始残差向量 $\bm r_0,\bm r_1,\bm r_2\in\mathcal G_0$ 及 $\bm x_0,\bm x_1,\bm x_2$. 按如下方式给出这些向量: 计算}
\[\begin{aligned}\bm r_0&=\bm b-A\bm x_0,\\\bm r_1&=(I-c_0A)\bm r_0,&\bm x_1&=\bm x_0+c_0\bm r_0,\\\bm r_2&=(I-c_1A)\bm r_1,&\bm x_2&=\bm x_1+c_1\bm r_1,\end{aligned}\]
\BT{where $\bm x_0$ is an arbitrary vector, and $c_0$ and $c_1$ are chosen so that $\|\bm r_1\|$ and $\|\bm r_2\|$ are minimized.}{其中 $\bm x_0$ 为任意向量, $c_0$ 和 $c_1$ 选择为使 $\|\bm r_1\|$ 和 $\|\bm r_2\|$ 达到极小值.}
% block: idr-general-algorithm-introduction page:148 type:ordinary-paragraph
\BT{From the above discussion, the IDR($s$) method with $s=2$ can be easily generalized to the IDR($s$) method with a general number $s\ge1$, and the algorithm is described in Algorithm\ \ref{alg:idr}. The symbol ``$\Delta X(:,i)$'' means that the $i$th column of matrix $\Delta X$.}{由上述讨论, $s=2$ 的 IDR($s$) 方法容易推广到一般数 $s\ge1$ 的 IDR($s$) 方法, 其算法见 Algorithm\ \ref{alg:idr}. 符号 ``$\Delta X(:,i)$'' 表示矩阵 $\Delta X$ 的第 $i$ 列.}
% block: idr-algorithm page:148 type:ordinary-paragraph
\begin{algorithm}[H]\caption{\BT{The IDR($s$) method}{IDR($s$) 方法}}\label{alg:idr}\begin{algorithmic}[1]
\Statex\textbf{Input:} $\bm x_0\in\mathbb C^N,\ P\in\mathbb C^{N\times s},\ s\ge1$
\Statex\textbf{Output:} $\bm x_i$
\State $\bm r_0=\bm b-A\bm x_0$
\Statex (\BT{The other initial residuals $\bm r_1,\ldots,\bm r_s\in\mathcal G_0$}{其他初始残差 $\bm r_1,\ldots,\bm r_s\in\mathcal G_0$})
\For{$i=1$ \BT{to}{至} $s$}
\State $\bm v=A\bm r_{i-1},\quad\omega=\frac{\bm v^{\mathrm H}\bm r_{i-1}}{\bm v^{\mathrm H}\bm v}$
\State $\Delta X(:,i)=\omega\bm r_{i-1},\quad\Delta R(:,i)=-\omega\bm v$
\State $\bm x_i=\bm x_{i-1}+\Delta X(:,i);\quad\bm r_i=\bm r_{i-1}+\Delta R(:,i)$
\EndFor
\State $j=1,\quad i=s$
\State $M=P^{\mathrm H}\Delta R,\quad\bm h=P^{\mathrm H}\bm r_i$
\While{$\frac{\|\bm r_i\|}{\|\bm b\|}>\epsilon$}
\For{$k=0$ \BT{to}{至} $s$}
\State\BT{Solve $\bm c$ from $M\bm c=\bm h$}{由 $M\bm c=\bm h$ 解出 $\bm c$}
\State $\bm q=-\Delta R\bm c$
\State $\bm v=\bm r_i+\bm q$
\If{$k=0$}
\State $\bm t=A\bm v,\quad\omega=\frac{\bm t^{\mathrm H}\bm v}{\bm t^{\mathrm H}\bm t}$
\State $\Delta R(:,j)=\bm q-\omega\bm t$
\State $\Delta X(:,j)=-\Delta X\bm c+\omega\bm v$
\Else
\State $\Delta X(:,j)=-\Delta X\bm c+\omega\bm v$
\State $\Delta R(:,j)=-A\Delta X(:,j)$
\EndIf
\Statex (\BT{Update approximate solutions $\bm x_i$}{更新近似解 $\bm x_i$})
\State $\bm r_{i+1}=\bm r_i+\Delta R(:,j)$
\State $\bm x_{i+1}=\bm x_i+\Delta X(:,j)$
\State $\Delta\bm m=P^{\mathrm H}\Delta R(:,j)$
\State $M(:,j)=\Delta\bm m$
\State $\bm h=\bm h+\Delta\bm m$
\State $i=i+1,\quad j=j+1$
\State $j=(j-1)\%s+1$ (\BT{$\%$: modulo operation, i.e.\ $a\%n=r$, where $a=mn+r$.}{$\%$: 取模运算, 即 $a\%n=r$, 其中 $a=mn+r$.})
\EndFor\EndWhile\end{algorithmic}\end{algorithm}
% block: idr-matvec-remark page:148 type:remark
\par\Needspace{3\baselineskip}
\begin{Remark}\label{rem:idr-matvec-bound}
\BT{In order to produce residual vectors in $\mathcal G_j$ from residual vectors in $\mathcal G_{j-1}$, the IDR($s$) method requires $s+1$ matrix--vector multiplications. From Remark\ \ref{rem:idr-dimension-count}, the residual vector in $\mathcal G_{N/s}$ is zero. Thus, in exact precision arithmetic, the number of matrix--vector multiplications for the IDR($s$) method is at most $(s+1)\times N/s=N+N/s$ to obtain the exact solution. Note that the number of matrix--vector multiplications for the product-type Krylov subspace methods in Sect.\ \ref{subsec:product-type-krylov} is at most $2N$.}{为从 $\mathcal G_{j-1}$ 中的残差向量产生 $\mathcal G_j$ 中的残差向量, IDR($s$) 方法需要 $s+1$ 次矩阵向量乘法. 由 Remark\ \ref{rem:idr-dimension-count}, $\mathcal G_{N/s}$ 中的残差向量为零. 因此, 在精确运算中, IDR($s$) 方法得到精确解所需的矩阵向量乘法次数至多为 $(s+1)\times N/s=N+N/s$. 注意, Section\ \ref{subsec:product-type-krylov} 中乘积型 Krylov 子空间方法的矩阵向量乘法次数至多为 $2N$.}
\end{Remark}
% block: idr-bicgstab-remark page:149 type:remark
\BT{From the following remark, the IDR($s$) method can be regarded as an extension of the BiCGSTAB method.}{由下面的注, IDR($s$) 方法可视为 BiCGSTAB 方法的推广.}
\par\Needspace{3\baselineskip}
\begin{Remark}\label{rem:idr-bicgstab-equivalence}
\BT{The IDR(1) method and the BiCGSTAB method give the same residual vectors at the even steps.}{IDR(1) 方法和 BiCGSTAB 方法在偶数步给出相同的残差向量.}
\end{Remark}
% block: idr-sonneveld-remark page:149 type:remark
\par\Needspace{3\baselineskip}
\begin{Remark}\label{rem:idr-sonneveld-subspace}
\BT{In terms of orthogonality, the IDR($s$) method produces residual vectors that belong to the following so-called Sonneveld subspace:}{从正交性看, IDR($s$) 方法产生的残差向量属于下列所谓的 Sonneveld 子空间:}
\begin{equation}\label{eq:idr-sonneveld-subspace}\{\Omega_j(A)\bm v:\bm v\perp\mathcal B_j(A^{\mathrm H},S)\}.\end{equation}
\BT{Here $\Omega_j(A):=\prod_{i=1}^j(I-\omega_iA)$, and $\mathcal B_j(A^{\mathrm H},S)$ is a subspace of $\mathbb C^N$ generated by using $A$ and $S:=[\bm s_1,\ldots,\bm s_s]$, which is defined by}{这里 $\Omega_j(A):=\prod_{i=1}^j(I-\omega_iA)$, $\mathcal B_j(A^{\mathrm H},S)$ 为使用 $A$ 和 $S:=[\bm s_1,\ldots,\bm s_s]$ 生成的 $\mathbb C^N$ 的子空间, 定义为}
\[\mathcal B_j(A^{\mathrm H},S):=\mathcal K_j(A^{\mathrm H},\bm s_1)+\cdots+\mathcal K_j(A^{\mathrm H},\bm s_s).\]
\BT{Here the symbol ``$+$'' means the sum of subspaces.}{这里符号 ``$+$'' 表示子空间之和.}\footnote{\BT{Given two subspaces $W_1$ and $W_2$, the sum of $W_1$ and $W_2$ is defined by $W_1+W_2:=\{\bm w_1+\bm w_2:\bm w_1\in W_1,\bm w_2\in W_2\}$.}{给定两个子空间 $W_1$ 和 $W_2$, $W_1$ 和 $W_2$ 的和定义为 $W_1+W_2:=\{\bm w_1+\bm w_2:\bm w_1\in W_1,\bm w_2\in W_2\}$.}}
\BT{$\mathcal B_j(A^{\mathrm H},R_0^{*})$ will be redefined in Definition\ \ref{def:block-krylov-subspace} for block Krylov subspace methods. For the details of the view in \eqref{eq:idr-sonneveld-subspace}, see \BookCite{161}, \BookCite{163}.}{$\mathcal B_j(A^{\mathrm H},R_0^{*})$ 将在块 Krylov 子空间方法的 Definition\ \ref{def:block-krylov-subspace} 中重新定义. 关于 \eqref{eq:idr-sonneveld-subspace} 这一观点的详情, 见 \BookCite{161}, \BookCite{163}.}
\end{Remark}
% block: idr-sonneveld-first-level page:149 type:ordinary-paragraph
\BT{We give a rough explanation of Remark\ \ref{rem:idr-sonneveld-subspace} using the IDR($s$) method with $s=2$ for the two cases $j=1,2$ in \eqref{eq:idr-sonneveld-subspace}. First we consider the case $j=1$. From \eqref{eq:idr-subspace-recurrence} $\mathcal G_1$ is generated by}{使用 $s=2$ 的 IDR($s$) 方法, 对 \eqref{eq:idr-sonneveld-subspace} 中 $j=1,2$ 两种情形, 粗略说明 Remark\ \ref{rem:idr-sonneveld-subspace}. 首先考虑 $j=1$ 的情形. 由 \eqref{eq:idr-subspace-recurrence}, $\mathcal G_1$ 由下式生成}
\[\mathcal G_1=(I-\omega_1A)(\mathcal G_0\cap\mathcal S),\]
\BT{where $\mathcal S=\operatorname{span}\{\bm s_1,\bm s_2\}^{\perp}$, see also \eqref{eq:idr-shadow-complement}. Then, $\bm r\in\mathcal G_1$ can be written as}{其中 $\mathcal S=\operatorname{span}\{\bm s_1,\bm s_2\}^{\perp}$, 另见 \eqref{eq:idr-shadow-complement}. 则 $\bm r\in\mathcal G_1$ 可写为}
\[\bm r=(I-\omega_1A)\bm v,\quad\bm v\in\mathcal G_0,\quad\bm v\perp\operatorname{span}\{\bm s_1,\bm s_2\}.\]
\BT{Note that $\bm v\in\mathcal S$ means $\bm v\perp\operatorname{span}\{\bm s_1,\bm s_2\}$. Since $\mathcal B_1(A^{\mathrm H},S)=\mathcal K_1(A^{\mathrm H},\bm s_1)+\mathcal K_1(A^{\mathrm H},\bm s_2)=\operatorname{span}\{\bm s_1,\bm s_2\}$, we have}{注意, $\bm v\in\mathcal S$ 意味着 $\bm v\perp\operatorname{span}\{\bm s_1,\bm s_2\}$. 由于 $\mathcal B_1(A^{\mathrm H},S)=\mathcal K_1(A^{\mathrm H},\bm s_1)+\mathcal K_1(A^{\mathrm H},\bm s_2)=\operatorname{span}\{\bm s_1,\bm s_2\}$, 有}
\[\bm r\in\{(I-\omega_1A)\bm v:\bm v\perp\mathcal B_1(A^{\mathrm H},S)\}=\{\Omega_1(A)\bm v:\bm v\perp\mathcal B_1(A^{\mathrm H},S)\},\]
\BT{which corresponds to the case $j=1$ in \eqref{eq:idr-sonneveld-subspace}.}{这对应于 \eqref{eq:idr-sonneveld-subspace} 中 $j=1$ 的情形.}
% block: idr-sonneveld-second-level pages:149,150 type:ordinary-paragraph
\BT{Next, we consider the case $j=2$. $\mathcal G_2$ is generated by}{接着考虑 $j=2$ 的情形. $\mathcal G_2$ 由下式生成}
\[\mathcal G_2=(I-\omega_2A)(\mathcal G_1\cap\mathcal S).\]
\BT{Then a residual vector $\widetilde{\bm r}\in\mathcal G_2$ can be written as}{则残差向量 $\widetilde{\bm r}\in\mathcal G_2$ 可写为}
\begin{equation}\label{eq:idr-second-level-residual}\widetilde{\bm r}=(I-\omega_2A)\widetilde{\bm v},\quad\widetilde{\bm v}\in\mathcal G_1,\quad\widetilde{\bm v}\perp\operatorname{span}\{\bm s_1,\bm s_2\}.\end{equation}
\BT{Since $\widetilde{\bm v}\in\mathcal G_1=(I-\omega_1A)(\mathcal G_0\cap\mathcal S)$, the vector $\widetilde{\bm v}$ can be written as}{由于 $\widetilde{\bm v}\in\mathcal G_1=(I-\omega_1A)(\mathcal G_0\cap\mathcal S)$, 向量 $\widetilde{\bm v}$ 可写为}
\begin{equation}\label{eq:idr-first-level-preimage}\widetilde{\bm v}=(I-\omega_1A)\bm v,\quad\bm v\in\mathcal G_0,\quad\bm v\perp\operatorname{span}\{\bm s_1,\bm s_2\}.\end{equation}
\BT{From \eqref{eq:idr-second-level-residual} and \eqref{eq:idr-first-level-preimage}, we have}{由 \eqref{eq:idr-second-level-residual} 和 \eqref{eq:idr-first-level-preimage}, 有}
\begin{equation}\label{eq:idr-two-factor-residual}\widetilde{\bm r}=(I-\omega_2A)(I-\omega_1A)\bm v,\end{equation}
\BT{where}{其中}
\[\widetilde{\bm v}=(I-\omega_1A)\bm v\perp\operatorname{span}\{\bm s_1,\bm s_2\},\quad\bm v\perp\operatorname{span}\{\bm s_1,\bm s_2\}.\]
\BT{Note that}{注意}
\[\begin{aligned}(I-\omega_1A)\bm v\perp\operatorname{span}\{\bm s_1,\bm s_2\}&\Leftrightarrow((I-\omega_1A)\bm v)^{\mathrm H}\bm s_i=0\quad(i=1,2)\\&\Leftrightarrow\bm v^{\mathrm H}(I-\omega_1A)^{\mathrm H}\bm s_i=0\quad(i=1,2)\\&\Leftrightarrow\bm v\perp\operatorname{span}\{(I-\omega_1A)^{\mathrm H}\bm s_1,(I-\omega_1A)^{\mathrm H}\bm s_2\}.\end{aligned}\]
\BT{$\bm v\perp\operatorname{span}\{(I-\omega_1A)^{\mathrm H}\bm s_1,(I-\omega_1A)^{\mathrm H}\bm s_2\}$ and $\bm v\perp\operatorname{span}\{\bm s_1,\bm s_2\}$ imply}{$\bm v\perp\operatorname{span}\{(I-\omega_1A)^{\mathrm H}\bm s_1,(I-\omega_1A)^{\mathrm H}\bm s_2\}$ 及 $\bm v\perp\operatorname{span}\{\bm s_1,\bm s_2\}$ 表明}
\[\bm v\perp\bm s_1,\quad\bm v\perp\bm s_2,\quad\bm v\perp A^{\mathrm H}\bm s_1,\quad\bm v\perp A^{\mathrm H}\bm s_1.\]
\BT{Thus,}{因此,}
\[\begin{aligned}\bm v\perp\operatorname{span}\{\bm s_1,A^{\mathrm H}\bm s_1,\bm s_2,A^{\mathrm H}\bm s_2\}.\\=\operatorname{span}\{\bm s_1,A^{\mathrm H}\bm s_1\}+\operatorname{span}\{\bm s_2,A^{\mathrm H}\bm s_2\}\\=\mathcal K_2(A^{\mathrm H},\bm s_1)+\mathcal K_2(A^{\mathrm H},\bm s_2)\\=\mathcal B_2(A^{\mathrm H},S).\end{aligned}\]
\BT{Then, from \eqref{eq:idr-two-factor-residual}, together with the above result, we have}{于是, 结合 \eqref{eq:idr-two-factor-residual} 与上述结果, 有}
\begin{equation}\label{eq:idr-second-level-block-orthogonality}\widetilde{\bm r}=(I-\omega_2A)(I-\omega_1A)\bm v,\quad\bm v\perp\mathcal B_2(A^{\mathrm H},S),\end{equation}
\BT{from which we obtain}{由此得到}
\begin{equation}\label{eq:idr-second-sonneveld-membership}\widetilde{\bm r}\in\{\Omega_2(A)\bm v:\bm v\perp\mathcal B_2(A^{\mathrm H},S)\}.\end{equation}
\BT{This corresponds to the case $j=2$ in \eqref{eq:idr-sonneveld-subspace}.}{这对应于 \eqref{eq:idr-sonneveld-subspace} 中 $j=2$ 的情形.}
% block: idr-related-literature pages:150,151 type:ordinary-paragraph
\BT{In what follows, some references related to the IDR($s$) method are described. An elegant explanation of the IDR($s$) method is found in \BookCite{88}, where a visualization of the IDR theorem in \BookCite{88}, Fig.3.1 is useful for intuitively understanding the IDR theorem (Theorem\ \ref{thm:idr-dimension-reduction}). A more stable and accurate variant of the IDR($s$) method is proposed by the same authors in \BookCite{199}, and the IDR($s$) method with the quasi-minimal residual strategy is proposed in \BookCite{47}. Variants of the IDR($s$) method with partial orthonormalization is proposed in \BookCite{209}. Extensions of the IDR($s$) methods using the idea of the BiCGSTAB($\ell$) method are known as the IDRstab method \BookCite{164} and the GBiCGSTAB($s,L$) method \BookCite{185}. Variants of the IDRstab method are found in \BookCite{5}--\BookCite{7}.}{下面介绍与 IDR($s$) 方法有关的一些文献. 文献 \BookCite{88} 对 IDR($s$) 方法给出了优美的解释, 其中 \BookCite{88}, Figure 3.1 对 IDR 定理的可视化有助于直观理解 IDR 定理 (Theorem\ \ref{thm:idr-dimension-reduction}). 同一批作者在 \BookCite{199} 中提出了更稳定且更精确的 IDR($s$) 变型, 文献 \BookCite{47} 提出了采用准最小残差策略的 IDR($s$) 方法. 文献 \BookCite{209} 提出了采用部分标准正交化的 IDR($s$) 变型. 使用 BiCGSTAB($\ell$) 方法思想的 IDR($s$) 推广称为 IDRstab 方法 \BookCite{164} 和 GBiCGSTAB($s,L$) 方法 \BookCite{185}. IDRstab 方法的变型见 \BookCite{5}--\BookCite{7}.}
% block: idr-numerical-example page:151 type:subsection
\paragraph{\BT{Numerical Example}{数值例子}}\label{subsec:idr-numerical-example}\mbox{}\par
\BT{In this subsection, typical convergence behavior of the IDR($s$) method is shown by one numerical experiment. We consider nonsymmetric linear systems \eqref{eq:three-dimensional-difference-system} with parameters $N=20$, $a_1=a_2=a_3=1$, $b_1=b_2=b_3=1$, and $c=300$. The stopping criterion is $\|\bm r_n\|/\|\bm b\|\le10^{-10}$.}{本小节通过一个数值实验展示 IDR($s$) 方法的典型收敛行为. 考虑非对称线性方程组 \eqref{eq:three-dimensional-difference-system}, 参数为 $N=20$, $a_1=a_2=a_3=1$, $b_1=b_2=b_3=1$, 以及 $c=300$. 停止准则为 $\|\bm r_n\|/\|\bm b\|\le10^{-10}$.}
\BT{The convergence histories of the IDR($s$) method ($s=1,2,4$) and the BiCGSTAB method are shown in Fig.\ \ref{fig:idr-comparison}. The horizontal axis is the number of matrix--vector multiplications.}{IDR($s$) 方法 ($s=1,2,4$) 和 BiCGSTAB 方法的收敛历程见 Figure\ \ref{fig:idr-comparison}. 横轴为矩阵向量乘法的次数.}
\BT{From Fig.\ \ref{fig:idr-comparison}, we see that as the number of $s$ in the IDR($s$) method increases, the required number of matrix--vector multiplications decreases. The IDR(4) method and the IDR(2) method converge faster than the BiCGSTAB method. In terms of accuracy, the true relative residual 2-norms for the BiCGSTAB method, the IDR(1) method, the IDR(2) method, and the IDR(4) method are $-10.50$, $-10.04$, $-10.2$, $-8.76$.}{由 Figure\ \ref{fig:idr-comparison} 可见, 随着 IDR($s$) 方法中的数 $s$ 增大, 所需矩阵向量乘法次数减少. IDR(4) 方法和 IDR(2) 方法比 BiCGSTAB 方法收敛更快. 从精度看, BiCGSTAB 方法, IDR(1) 方法, IDR(2) 方法和 IDR(4) 方法的真实相对残差 2-范数分别为 $-10.50$, $-10.04$, $-10.2$, $-8.76$.}
% block: idr-comparison-figure page:151 type:figure
\begin{figure}[H]\centering\includegraphics[width=.87\linewidth]{figures/idr-comparison.png}
\caption{\BT{Convergence histories of BiCGSTAB, IDR(1), IDR(2), and IDR(4)}{BiCGSTAB, IDR(1), IDR(2) 和 IDR(4) 的收敛历程}}\label{fig:idr-comparison}\end{figure}
% block: block-idr-heading page:152 type:subsection
\subsubsection{\BT{Block Induced Dimension Reduction (Block IDR($s$)) Method}{块诱导降维 (块 IDR($s$)) 方法}}\label{subsec:block-idr-method}
\BT{The block IDR($s$) method \BookCite{46} is an extension of the IDR($s$) method to solve the block linear systems of the form}{块 IDR($s$) 方法 \BookCite{46} 是 IDR($s$) 方法的推广, 用以求解形如下式的块线性方程组}
\begin{equation}\label{eq:block-linear-system}AX=B,\end{equation}
\BT{where $A\in\mathbb C^{N\times N}$, $B\in\mathbb C^{N\times m}$, and $m$ is usually much less than $N$.}{其中 $A\in\mathbb C^{N\times N}$, $B\in\mathbb C^{N\times m}$, 且 $m$ 通常远小于 $N$.}
\BT{Prior to the theory of the block IDR($s$) method, the framework of block Krylov subspace methods is briefly described after the following definitions:}{在介绍块 IDR($s$) 方法的理论之前, 先在下列定义之后简要描述块 Krylov 子空间方法的框架:}
% block: block-krylov-definition page:152 type:definition
\par\Needspace{3\baselineskip}
\begin{Definition}{\BT{Block Krylov subspace}{块 Krylov 子空间}}\label{def:block-krylov-subspace}
\BT{Let $V\in\mathbb C^{N\times m}$. Then $\mathcal K_n(A,V)\subset\mathbb C^{N\times m}$ defined by}{设 $V\in\mathbb C^{N\times m}$. 则由下式定义的 $\mathcal K_n(A,V)\subset\mathbb C^{N\times m}$}
\[\begin{aligned}\mathcal K_n(A,V)&:=\{VC_0+AVC_1+\cdots+A^{n-1}VC_{n-1}:C_0,C_1,\ldots,C_{n-1}\in\mathbb C^{m\times m}\}\\&(=\operatorname{block\ span}\{V,AV,\ldots,A^{n-1}V\})\end{aligned}\]
\BT{is called the block Krylov subspace with respect to $A$ and $V$.}{称为关于 $A$ 和 $V$ 的块 Krylov 子空间.}
\end{Definition}
% block: block-krylov-sum-definition page:152 type:definition
\par\Needspace{3\baselineskip}
\begin{Definition}{\BT{Sum of Krylov subspaces}{Krylov 子空间之和}}\label{def:block-krylov-sum}
\BT{Let $V=[\bm v_1,\bm v_2,\ldots,\bm v_m]\in\mathbb C^{N\times m}$. Then $\mathcal B_n(A,V)\subset\mathbb C^N$ is defined by}{设 $V=[\bm v_1,\bm v_2,\ldots,\bm v_m]\in\mathbb C^{N\times m}$. 则 $\mathcal B_n(A,V)\subset\mathbb C^N$ 定义为}
\[\begin{aligned}\mathcal B_n(A,V)&:=\mathcal K_n(A,\bm v_1)+\mathcal K_n(A,\bm v_2)+\cdots+\mathcal K_n(A,\bm v_m)\\&=\{V\bm c_0+AV\bm c_1+\cdots+A^{n-1}V\bm c_{n-1}:\bm c_0,\bm c_1,\ldots,\bm c_{n-1}\in\mathbb C^m\}.\end{aligned}\]
\end{Definition}
% block: block-krylov-grade-definition page:152 type:definition
\par\Needspace{3\baselineskip}
\begin{Definition}{\BT{Grade of block Krylov subspace}{块 Krylov 子空间的阶}}\label{def:block-krylov-grade}
\BT{The smallest number $n$ such that}{满足下式的最小数 $n$}
\[\dim(\mathcal B_n(A,V))=\dim(\mathcal B_{n+1}(A,V))\]
\BT{is called the block grade of $A$ with respect to $V$, which is denoted by $\nu(V,A)$.}{称为 $A$ 关于 $V$ 的块阶, 记为 $\nu(V,A)$.}
\end{Definition}
% block: block-grade-reduction-remark page:152 type:remark
\par\Needspace{3\baselineskip}
\begin{Remark}\label{rem:block-grade-scalar-reduction}
\BT{Definitions\ \ref{def:block-krylov-subspace} and \ref{def:block-krylov-grade} are extensions of Definitions\ \ref{def:krylov-subspace} and \ref{def:krylov-grade}. In fact, if $m=1$, these definitions are identical to Definitions\ \ref{def:krylov-subspace} and \ref{def:krylov-grade}.}{Definition\ \ref{def:block-krylov-subspace} 和 Definition\ \ref{def:block-krylov-grade} 是 Definition\ \ref{def:krylov-subspace} 和 Definition\ \ref{def:krylov-grade} 的推广. 事实上, 若 $m=1$, 这些定义就与 Definition\ \ref{def:krylov-subspace} 和 Definition\ \ref{def:krylov-grade} 相同.}
\end{Remark}
% block: block-krylov-framework pages:152,153 type:ordinary-paragraph
\BT{We now describe the framework of the block Krylov subspace methods. Let $X_0\in\mathbb C^{N\times s}$ be an initial guess, and let $R_0:=B-AX_0$ be the corresponding initial residual matrix. Then, block Krylov subspace methods find the $n$th approximate solution $X_n$ as follows:}{现在描述块 Krylov 子空间方法的框架. 设 $X_0\in\mathbb C^{N\times s}$ 为初始猜测, $R_0:=B-AX_0$ 为相应的初始残差矩阵. 则块 Krylov 子空间方法按下式寻找第 $n$ 个近似解 $X_n$:}
\[X_n=X_0+Z_n,\quad Z_n\in\mathcal K_n(A,R_0).\]
\BT{The corresponding residual matrix is given by}{相应的残差矩阵为}
\[R_n=B-AX_n=R_0-AZ_n\in\mathcal K_{n+1}(A,R_0).\]
\BT{Determining $Z_n$ yields various block Krylov subspace methods, e.g.,}{确定 $Z_n$ 就得到各种块 Krylov 子空间方法, 例如,}
\begin{itemize}
\item\BT{The block CG method:}{块 CG 方法:} $R_n\perp\mathcal B_n(A,R_0)$.
\item\BT{The block COCG method:}{块 COCG 方法:} $R_n\perp\overline{\mathcal B_n(A,R_0)}$.
\item\BT{The block BiCG method:}{块 BiCG 方法:} $R_n\perp\mathcal B_n(A^{\mathrm H},R_0^{*})$.
\end{itemize}
\BT{Here, $R_0^{*}$ is an initial shadow residual matrix,}{这里 $R_0^{*}$ 为初始影子残差矩阵,}\footnote{\BT{In practice, $R_0^{*}$ is usually set to $R_0^{*}=R_0$ or a random matrix.}{实际中通常取 $R_0^{*}=R_0$, 或将 $R_0^{*}$ 设为随机矩阵.}}
\BT{and the symbol $R_n\perp\mathcal B_n(A,V)$ means that all the column vectors of $R_n$ are orthogonal to $\mathcal B_n(A,V)$.}{符号 $R_n\perp\mathcal B_n(A,V)$ 表示 $R_n$ 的所有列向量都与 $\mathcal B_n(A,V)$ 正交.}
% block: block-krylov-history page:153 type:ordinary-paragraph
\BT{Here, we give a brief history of the block Krylov subspace methods. In 1980, O’Leary proposed the block BiCG method (and the block CG method) in \BookCite{142}. In 1990, Vital proposed the block GMRES method in \BookCite{202}. Block Krylov subspace methods based on the QMR method were proposed in 1997 \BookCite{65}, \BookCite{157}. In 2003, Guennouni et al. proposed the block BiCGSTAB method \BookCite{86}. In 2011, Du et al. proposed the block IDR($s$) method \BookCite{46}. For other block Krylov subspace methods, see \BookCite{153}, \BookCite{181}, \BookCite{182}, \BookCite{210}.}{这里简要回顾块 Krylov 子空间方法的历史. 1980 年, O’Leary 在 \BookCite{142} 中提出块 BiCG 方法 (及块 CG 方法). 1990 年, Vital 在 \BookCite{202} 中提出块 GMRES 方法. 1997 年提出了基于 QMR 方法的块 Krylov 子空间方法 \BookCite{65}, \BookCite{157}. 2003 年, Guennouni 等提出块 BiCGSTAB 方法 \BookCite{86}. 2011 年, Du 等提出块 IDR($s$) 方法 \BookCite{46}. 其他块 Krylov 子空间方法见 \BookCite{153}, \BookCite{181}, \BookCite{182}, \BookCite{210}.}
\BT{Now we turn our attention to deriving the block IDR($s$) method. The block IDR($s$) method is based on the following theorem (see \BookCite{46}), which is a slight modification of Theorem\ \ref{thm:idr-dimension-reduction}.}{现在转向块 IDR($s$) 方法的推导. 块 IDR($s$) 方法基于下述定理 (见 \BookCite{46}), 它是 Theorem\ \ref{thm:idr-dimension-reduction} 的轻微修改.}
% block: block-idr-theorem page:153 type:theorem
\par\Needspace{14\baselineskip}
\begin{Theorem}\label{thm:block-idr-dimension-reduction}
\BT{Let $A$ be any matrix in $\mathbb C^{N\times N}$, let $\bm v_0$ be any nonzero vector in $\mathbb C^N$, and let $\mathcal G_0=\mathcal B_{\nu(R_0,A)}(A,R_0)$. Let $\mathcal S$ denote any (proper) subspace of $\mathbb C^N$ such that $\mathcal S$ and $\mathcal G_0$ do not share a nontrivial invariant subspace,}{设 $A$ 为 $\mathbb C^{N\times N}$ 中的任意矩阵, $\bm v_0$ 为 $\mathbb C^N$ 中的任意非零向量, 并设 $\mathcal G_0=\mathcal B_{\nu(R_0,A)}(A,R_0)$. 设 $\mathcal S$ 表示 $\mathbb C^N$ 的任意 (真) 子空间, 使 $\mathcal S$ 和 $\mathcal G_0$ 不具有共同的非平凡不变子空间,}\footnote{\BT{This means that $\mathcal S\cap\mathcal G_0$ does not contain any eigenvector of $A$, and the trivial invariant subspace is $\{\bm0\}$.}{这意味着 $\mathcal S\cap\mathcal G_0$ 不包含 $A$ 的任何特征向量, 平凡不变子空间为 $\{\bm0\}$.}}
\BT{and define the sequences $\mathcal G_j$, $j=1,2,\ldots$ as}{并定义序列 $\mathcal G_j$, $j=1,2,\ldots$ 为}
\begin{equation}\label{eq:block-idr-subspace-recurrence}\mathcal G_j=(I-\omega_jA)(\mathcal G_{j-1}\cap\mathcal S),\end{equation}
\BT{where the $\omega_j$'s are nonzero scalars. Then:}{其中 $\omega_j$ 为非零标量. 则:}
\begin{enumerate}[label=(\arabic*)]
\item $\mathcal G_j\subset\mathcal G_{j-1}$ \BT{for all $j>0$.}{对所有 $j>0$ 成立.}
\item $\mathcal G_j=\{\bm0\}$ \BT{for some $j\le N$.}{对某个 $j\le N$ 成立.}
\end{enumerate}\end{Theorem}
% block: block-idr-theorem-proof pages:153,154 type:proof
\begin{Proof}
\BT{The only difference between Theorems\ \ref{thm:idr-dimension-reduction} and \ref{thm:block-idr-dimension-reduction} is the definition of $\mathcal G_0$, i.e., $\mathcal G_0=\mathcal K_N(A,\bm v_0)$ or $\mathcal G_0=\mathcal B_{\nu(R_0,A)}(A,R_0)$. In the proof of Theorem\ \ref{thm:idr-dimension-reduction}-(1), $\mathcal G_0=\mathcal K_N(A,\bm v_0)$ is used only for proving $\mathcal G_1\subset\mathcal G_0$. Thus, all we have to do is to prove $\mathcal G_1\subset\mathcal G_0$. From $\mathcal G_0=\mathcal B_{\nu(R_0,A)}(A,R_0)$, we have $(I-\omega_1A)\mathcal G_0\subset\mathcal G_0$, because for a given $\bm v\in(I-\omega_1A)\mathcal G_0=(I-\omega_1A)\mathcal B_{\nu(R_0,A)}(A,R_0)$ we have}{Theorem\ \ref{thm:idr-dimension-reduction} 和 Theorem\ \ref{thm:block-idr-dimension-reduction} 的唯一区别在于 $\mathcal G_0$ 的定义, 即 $\mathcal G_0=\mathcal K_N(A,\bm v_0)$ 或 $\mathcal G_0=\mathcal B_{\nu(R_0,A)}(A,R_0)$. 在 Theorem\ \ref{thm:idr-dimension-reduction}-(1) 的证明中, $\mathcal G_0=\mathcal K_N(A,\bm v_0)$ 仅用于证明 $\mathcal G_1\subset\mathcal G_0$. 因此只需证明 $\mathcal G_1\subset\mathcal G_0$. 由 $\mathcal G_0=\mathcal B_{\nu(R_0,A)}(A,R_0)$ 有 $(I-\omega_1A)\mathcal G_0\subset\mathcal G_0$, 因为对于给定的 $\bm v\in(I-\omega_1A)\mathcal G_0=(I-\omega_1A)\mathcal B_{\nu(R_0,A)}(A,R_0)$, 有}
\[\begin{aligned}\bm v&=(I-\omega_1A)(R_0\bm c_0+AR_0\bm c_1+\cdots+A^{\nu(R_0,A)-1}R_0\bm c_{\nu(R_0,A)-1})\\&=(I-\omega_1A)R_0\bm c_0+(I-\omega_1A)AR_0\bm c_1+\cdots+(I-\omega_1A)A^{\nu(R_0,A)-1}R_0\bm c_{\nu(R_0,A)-1}\\&=R_0\bm d_0+AR_0\bm d_1+\cdots+A^{\nu(R_0,A)-1}R_0\bm d_{\nu(R_0,A)-1}+A^{\nu(R_0,A)}R_0\bm d_{\nu(R_0,A)}\\&=R_0\widetilde{\bm d}_0+AR_0\widetilde{\bm d}_1+\cdots+A^{\nu(R_0,A)-1}R_0\widetilde{\bm d}_{\nu(R_0,A)-1}\\&\in\mathcal G_0.\end{aligned}\]
\BT{Thus $\mathcal G_1=(I-\omega_1A)(\mathcal G_0\cap\mathcal S)\subset(I-\omega_1A)\mathcal G_0\subset\mathcal G_0$.}{因此 $\mathcal G_1=(I-\omega_1A)(\mathcal G_0\cap\mathcal S)\subset(I-\omega_1A)\mathcal G_0\subset\mathcal G_0$.}
\BT{Theorem\ \ref{thm:block-idr-dimension-reduction}-(2) is shown by following the proof of Theorem\ \ref{thm:idr-dimension-reduction}-(2) and using $\dim\mathcal G_0\le N$.}{按照 Theorem\ \ref{thm:idr-dimension-reduction}-(2) 的证明并使用 $\dim\mathcal G_0\le N$, 即可证明 Theorem\ \ref{thm:block-idr-dimension-reduction}-(2).}
\par\nopagebreak[4]
\end{Proof}
% block: block-idr-dimension-corollary page:154 type:corollary
\BT{Similar to Corollary\ \ref{cor:idr-dimension-step}, we have the following fact:}{与 Corollary\ \ref{cor:idr-dimension-step} 类似, 有下列事实:}
\par\Needspace{3\baselineskip}
\begin{Corollary}\label{cor:block-idr-dimension-step}
\BT{Let $\mathcal S$ be a subspace of $\mathbb C^N$ with $\dim(\mathcal S)=N-sm$ for a given $s,m\ge1$. If $\mathcal G_{i-1}+\mathcal S=\mathbb C^N$ and $(I-\omega_jA)$ is nonsingular, then}{对于给定的 $s,m\ge1$, 设 $\mathcal S$ 为 $\mathbb C^N$ 的子空间, 且 $\dim(\mathcal S)=N-sm$. 若 $\mathcal G_{i-1}+\mathcal S=\mathbb C^N$ 且 $(I-\omega_jA)$ 非奇异, 则}
\[\dim(\mathcal G_j)=\dim(\mathcal G_{j-1})-sm.\]
\end{Corollary}
\BT{The difference between Corollaries\ \ref{cor:idr-dimension-step}, \ref{cor:block-idr-dimension-step} is the parameter $m$. The parameter $m$ corresponds to the size of the number of right-hand sides of $AX=B$, where $B$ is an $N$-by-$m$ matrix.}{Corollary\ \ref{cor:idr-dimension-step} 与 Corollary\ \ref{cor:block-idr-dimension-step} 的区别在于参数 $m$. 参数 $m$ 对应于 $AX=B$ 的右端项个数, 其中 $B$ 为 $N$ 行 $m$ 列的矩阵.}
% block: block-idr-subspace-choice page:154 type:ordinary-paragraph
\BT{In what follows, the block IDR($s$) method is explained. The block IDR($s$) method generates residual matrices $R_k=B-AX_k$ for $k=0,1,\ldots$ such that:}{下面解释块 IDR($s$) 方法. 块 IDR($s$) 方法对 $k=0,1,\ldots$ 生成残差矩阵 $R_k=B-AX_k$, 使得:}
\begin{enumerate}
\item\BT{all the column vectors of $R_{j(s+1)},R_{j(s+1)+1},\ldots,R_{j(s+1)+s}$ belong to $\mathcal G_j$ in \eqref{eq:block-idr-subspace-recurrence} for $j=0,1,2,\ldots$;}{对 $j=0,1,2,\ldots$, $R_{j(s+1)},R_{j(s+1)+1},\ldots,R_{j(s+1)+s}$ 的所有列向量属于 \eqref{eq:block-idr-subspace-recurrence} 中的 $\mathcal G_j$;}
\item\BT{the subspace $\mathcal S$ in \eqref{eq:block-idr-subspace-recurrence} is chosen so that $\dim\mathcal S=N-sm$.}{选择 \eqref{eq:block-idr-subspace-recurrence} 中的子空间 $\mathcal S$, 使 $\dim\mathcal S=N-sm$.}
\end{enumerate}
\BT{Let $S=[\bm s_1,\bm s_2,\ldots,\bm s_{sm}]\in\mathbb C^{N\times sm}$. Then, a standard choice of the subspace $\mathcal S$ is as follows:}{设 $S=[\bm s_1,\bm s_2,\ldots,\bm s_{sm}]\in\mathbb C^{N\times sm}$. 则子空间 $\mathcal S$ 的标准选择如下:}
\[\mathcal S=\operatorname{span}\{S\}^{\perp}=\operatorname{span}\{\bm s_1,\bm s_2,\ldots,\bm s_{sm}\}^{\perp},\]
\BT{where all the column vectors of $S\in\mathbb C^{N\times sm}$ are linearly independent, and thus it is easy to see that $\dim\mathcal S=N-sm$.}{其中 $S\in\mathbb C^{N\times sm}$ 的所有列向量线性无关, 因此容易看出 $\dim\mathcal S=N-sm$.}
% block: block-idr-two-residuals page:154 type:ordinary-paragraph
\BT{We now briefly explain the block IDR($s$) method with $s=2$. For $j=0,1,2$, the block IDR(2) method produces the following residual matrices:}{现在简要说明 $s=2$ 的块 IDR($s$) 方法. 对 $j=0,1,2$, 块 IDR(2) 方法产生下列残差矩阵:}
\[R_0,R_1,R_2,\quad R_3,R_4,R_5,\quad R_6,R_7,R_8,\]
\BT{where all the column vectors $R_0,R_1,R_2$ belong to $\mathcal G_0$, all the column vectors $R_3,R_4,R_5$ belong to $\mathcal G_1$, and all the column vectors $R_3,R_4,R_5$ belong to $\mathcal G_2$.}{其中 $R_0,R_1,R_2$ 的所有列向量属于 $\mathcal G_0$, $R_3,R_4,R_5$ 的所有列向量属于 $\mathcal G_1$, $R_3,R_4,R_5$ 的所有列向量属于 $\mathcal G_2$.}
\BT{In what follows, we describe how to obtain $R_3,R_4,R_5$ from $R_0,R_1,R_2$ under the assumption that all the column vectors of $R_0,R_1,R_2$ belong to $\mathcal G_0$. Similar to the IDR($s$) method, the following definitions are useful for the block IDR($s$) method:}{下面在 $R_0,R_1,R_2$ 的所有列向量属于 $\mathcal G_0$ 的假设下, 说明如何从 $R_0,R_1,R_2$ 得到 $R_3,R_4,R_5$. 与 IDR($s$) 方法类似, 下列定义对块 IDR($s$) 方法有用:}
\begin{equation}\label{eq:block-idr-residual-differences}\Delta R_k:=R_{k+1}-R_k,\quad dR_k^{\blacksquare}:=[\Delta R_k,\Delta R_{k+1}],\quad C_k=(S^{\mathrm H}dR_k^{\blacksquare})^{-1}S^{\mathrm H}R_{k+2}.\end{equation}
\BT{All the column vectors of the following matrix $V_0$ lie in $\mathcal G_0\cap\mathcal S$:}{下列矩阵 $V_0$ 的所有列向量位于 $\mathcal G_0\cap\mathcal S$ 中:}
\[V_0=R_2-dR_0^{\blacksquare}C_0,\quad\operatorname{Im}(V_0)\subset\mathcal G_0\cap\mathcal S\]
% block: block-idr-first-reduced-residual page:155 type:ordinary-paragraph
\BT{because all the column vectors of $R_2\in\mathbb C^{N\times m}$ and $dR_0\in\mathbb C^{N\times2m}$ belong to $\mathcal G_0$ from the assumption that all the column vectors of $R_0,R_1,R_2$ belong to $\mathcal G_0$, and all the column vectors of $V_0$ also belong to $\mathcal S$ from the relation $S^{\mathrm H}V_0=O$, where $O$ is the $2m\times m$ zero matrix. Here $\operatorname{Im}(V_0)$ is the image of $V_0$, i.e., the subspace spanned by all the column vectors of $V_0$.}{因为由 $R_0,R_1,R_2$ 的所有列向量属于 $\mathcal G_0$ 这一假设, $R_2\in\mathbb C^{N\times m}$ 和 $dR_0\in\mathbb C^{N\times2m}$ 的所有列向量属于 $\mathcal G_0$; 又由关系 $S^{\mathrm H}V_0=O$, $V_0$ 的所有列向量也属于 $\mathcal S$, 其中 $O$ 为 $2m\times m$ 零矩阵. 这里 $\operatorname{Im}(V_0)$ 为 $V_0$ 的像, 即 $V_0$ 的所有列向量张成的子空间.}
\BT{From \eqref{eq:block-idr-subspace-recurrence}, the residual matrix $R_3$ whose column vectors belong to $\mathcal G_1$ can be produced by}{由 \eqref{eq:block-idr-subspace-recurrence}, 列向量属于 $\mathcal G_1$ 的残差矩阵 $R_3$ 可按下式产生}
\begin{equation}\label{eq:block-idr-first-reduced-residual}R_3=(I-\omega_0A)V_0,\quad\operatorname{Im}(R_3)\subset\mathcal G_1,\end{equation}
\BT{where $\omega_0$ is usually chosen so that $\|R_3\|_{\mathrm F}$ is minimized,}{其中通常选择 $\omega_0$ 使 $\|R_3\|_{\mathrm F}$ 达到极小值,}\footnote{\BT{For $R=(r_{i,j})\in\mathbb C^{N\times m}$, the symbol $\|R\|_{\mathrm F}:=(\sum_{i=1}^N\sum_{j=1}^m\bar r_{i,j}r_{i,j})^{1/2}$ is called the Frobenius norm of $R$, and for a square matrix $A=(a_{i,j})\in\mathbb C^{N\times N}$, the symbol $\operatorname{Tr}(A):=\sum_{i=1}^Na_{i,i}$ is called the trace of $A$.}{对 $R=(r_{i,j})\in\mathbb C^{N\times m}$, 符号 $\|R\|_{\mathrm F}:=(\sum_{i=1}^N\sum_{j=1}^m\bar r_{i,j}r_{i,j})^{1/2}$ 称为 $R$ 的 Frobenius 范数; 对方阵 $A=(a_{i,j})\in\mathbb C^{N\times N}$, 符号 $\operatorname{Tr}(A):=\sum_{i=1}^Na_{i,i}$ 称为 $A$ 的迹.}}
\BT{i.e.,}{即,}
\[\omega_0=\frac{\operatorname{Tr}(T_0^{\mathrm H}V_0)}{\operatorname{Tr}(T_0^{\mathrm H}T_0)},\quad T_0:=AV_0.\]
% block: frobenius-minimization-proposition page:155 type:proposition
\BT{The Frobenius norm minimization follows from the following fact:}{Frobenius 范数的极小化来自下列事实:}
\par\Needspace{3\baselineskip}
\begin{Proposition}\label{prop:frobenius-scalar-minimization}
\BT{Let $V,W\in\mathbb C^{m\times n}$. Then the following minimization problem:}{设 $V,W\in\mathbb C^{m\times n}$. 则下列极小化问题:}
\[\min_{\omega\in\mathbb C}\|V-\omega W\|_{\mathrm F}\]
\BT{can be solved by}{可由下式求解}
\[\omega=\frac{\operatorname{Tr}(W^{\mathrm H}V)}{\operatorname{Tr}(W^{\mathrm H}W)}.\]
\end{Proposition}
% block: frobenius-minimization-proof page:155 type:proof
\begin{Proof}
\BT{Let $V=[\bm v_1,\ldots,\bm v_n]$ and $W=[\bm w_1,\ldots,\bm w_n]$, then using}{设 $V=[\bm v_1,\ldots,\bm v_n]$, $W=[\bm w_1,\ldots,\bm w_n]$, 则使用}
\[\bm v:=\operatorname{vec}(V)=\begin{bmatrix}\bm v_1\\\vdots\\\bm v_n\end{bmatrix},\quad\bm w:=\operatorname{vec}(W)=\begin{bmatrix}\bm w_1\\\vdots\\\bm w_n\end{bmatrix}\]
\BT{yields $\|V-\omega W\|_{\mathrm F}=\|\bm v-\omega\bm w\|$. Thus from \eqref{eq:scalar-residual-minimizer}, the minimizer is}{得到 $\|V-\omega W\|_{\mathrm F}=\|\bm v-\omega\bm w\|$. 因而由 \eqref{eq:scalar-residual-minimizer}, 极小化参数为}
\[\omega=\frac{\bm w^{\mathrm H}\bm v}{\bm w^{\mathrm H}\bm w}=\frac{\operatorname{Tr}(W^{\mathrm H}V)}{\operatorname{Tr}(W^{\mathrm H}W)},\]
\BT{which concludes the proof.}{证明完毕.}
\par\nopagebreak[4]
\end{Proof}
% block: block-idr-next-residuals pages:155,156 type:ordinary-paragraph
\BT{Similar to the derivation of $\bm r_3$ and $\bm r_4$ in the IDR($s$) method, $R_3$ and $R_4$ are obtained as follows:}{与 IDR($s$) 方法中 $\bm r_3$ 和 $\bm r_4$ 的推导类似, 按下式得到 $R_3$ 和 $R_4$:}
\[\begin{aligned}V_1&=R_3-dR_1^{\blacksquare}C_1,&\operatorname{Im}(V_1)&\subset\mathcal G_0\cap\mathcal S,\\R_4&=(I-\omega_0A)V_1,&\operatorname{Im}(R_4)&\subset\mathcal G_1,\\V_2&=R_4-dR_2^{\blacksquare}C_2,&\operatorname{Im}(V_2)&\subset\mathcal G_0\cap\mathcal S,\\R_5&=(I-\omega_0A)V_2,&\operatorname{Im}(R_5)&\subset\mathcal G_1.\end{aligned}\]
\BT{From the above calculations, we have residual matrices $R_k$ ($k=0,\ldots,5$) such that}{由上述计算, 得到残差矩阵 $R_k$ ($k=0,\ldots,5$), 使得}
\[\begin{aligned}\operatorname{Im}(R_k)&\subset\mathcal G_0\quad\text{for }k=0,1,2,\\\operatorname{Im}(R_k)&\subset\mathcal G_1\quad\text{for }k=3,4,5.\end{aligned}\]
% block: block-idr-solution-extraction page:156 type:ordinary-paragraph
\BT{Next, we describe how to obtain approximate solution matrices $X_3,X_4,X_5$ from $X_0,X_1,X_2$. Let}{接着说明如何从 $X_0,X_1,X_2$ 得到近似解矩阵 $X_3,X_4,X_5$. 令}
\[\Delta X_k:=X_{k+1}-X_k,\qquad dX_k^{\blacksquare}:=[\Delta X_k,\Delta X_{k+1}].\]
\BT{Then $X_3$ is obtained as follows:}{则按下式得到 $X_3$:}
\[\begin{aligned}R_3&=(I-\omega_0A)V_0=R_2-dR_0^{\blacksquare}C_0-\omega_0AV_0\\&\Leftrightarrow B-AX_3=B-AX_2-dR_0^{\blacksquare}C_0-\omega_0AV_0\\&\Leftrightarrow X_3=X_2+A^{-1}dR_0^{\blacksquare}C_0+\omega_0V_0\\&\Leftrightarrow X_3=X_2-dX_0^{\blacksquare}C_0+\omega_0V_0.\end{aligned}\]
\BT{Similarly,}{类似地,}
\[\begin{aligned}X_4&=X_3-dX_1^{\blacksquare}C_1+\omega_0V_1,\\X_5&=X_4-dX_2^{\blacksquare}C_2+\omega_0V_2.\end{aligned}\]
\BT{We have described how to obtain $R_k$ and $X_k$ for $k=3,4,5$ from $R_k$ and $X_k$ for $k=0,1,2$ in the block IDR($s$) method with $s=2$. Following the above derivation, we have the block IDR($s$) method in Algorithm\ \ref{alg:block-idr}. Further, the preconditioned version of the block IDR($s$) method is shown in Algorithm\ \ref{alg:preconditioned-block-idr}. Note that the symbol ``$B(:,\ell)$'' represents the $\ell$th column of matrix $B$, and the symbol ``$\Delta X(:,im+1:(i+1)m)$'' represents the $(im+1)$th, $(im+2)$th, $\ldots$, $(i+1)m$th columns of matrix $\Delta X$.}{已说明在 $s=2$ 的块 IDR($s$) 方法中, 如何从 $k=0,1,2$ 的 $R_k$ 和 $X_k$ 得到 $k=3,4,5$ 的 $R_k$ 和 $X_k$. 按照上述推导, 得到 Algorithm\ \ref{alg:block-idr} 中的块 IDR($s$) 方法. 此外, 块 IDR($s$) 方法的预条件版本见 Algorithm\ \ref{alg:preconditioned-block-idr}. 注意, 符号 ``$B(:,\ell)$'' 表示矩阵 $B$ 的第 $\ell$ 列, 符号 ``$\Delta X(:,im+1:(i+1)m)$'' 表示矩阵 $\Delta X$ 的第 $(im+1)$, $(im+2)$, $\ldots$, $(i+1)m$ 列.}
% block: block-idr-cost-remark pages:156,157 type:remark
\par\Needspace{3\baselineskip}
\begin{Remark}\label{rem:block-idr-matvec-bound}
\BT{From Corollary\ \ref{cor:block-idr-dimension-step}, the dimension reduction is usually $sm$, i.e., $\dim(\mathcal G_j)=\dim(\mathcal G_{j-1})-sm$ and the block IDR($s$) method requires $m(s+1)$ matrix--vector multiplications. Thus, in exact precision arithmetic, the total number of matrix--vector multiplications for the block IDR($s$) method is at most $m(s+1)\times N/sm=N+N/s$ to obtain the exact solution matrix, which does not depend on the block size $m$. Note that if the IDR($s$) method is applied to $AX=B$ (i.e., $A\bm x_1=\bm b_1,A\bm x_2=\bm b_2,\ldots,A\bm x_m=\bm b_m$), then the required number of matrix--vector multiplications of the IDR($s$) method is at most $m(N+N/s)$, which is $m$ times larger than that of the block IDR($s$) method.}{由 Corollary\ \ref{cor:block-idr-dimension-step}, 维数的减少量通常为 $sm$, 即 $\dim(\mathcal G_j)=\dim(\mathcal G_{j-1})-sm$, 块 IDR($s$) 方法需要 $m(s+1)$ 次矩阵向量乘法. 因此, 在精确运算中, 块 IDR($s$) 方法得到精确解矩阵所需的矩阵向量乘法总次数至多为 $m(s+1)\times N/sm=N+N/s$, 不依赖于块大小 $m$. 注意, 若对 $AX=B$ 应用 IDR($s$) 方法 (即 $A\bm x_1=\bm b_1,A\bm x_2=\bm b_2,\ldots,A\bm x_m=\bm b_m$), 则 IDR($s$) 方法所需的矩阵向量乘法次数至多为 $m(N+N/s)$, 是块 IDR($s$) 方法的 $m$ 倍.}
\end{Remark}
% block: block-idr-bicgstab-remark page:157 type:remark
\BT{From the following remark, the block IDR($s$) method can be regarded as an extension of the block BiCGSTAB method.}{由下面的注, 块 IDR($s$) 方法可视为块 BiCGSTAB 方法的推广.}
\par\Needspace{3\baselineskip}
\begin{Remark}\label{rem:block-idr-bicgstab-equivalence}
\BT{The block IDR(1) method and the block BiCGSTAB method give the same residual vectors at the even steps.}{块 IDR(1) 方法和块 BiCGSTAB 方法在偶数步给出相同的残差向量.}
\end{Remark}
% block: block-idr-algorithm page:157 type:ordinary-paragraph
\begin{algorithm}[H]\caption{\BT{The block IDR($s$) method}{块 IDR($s$) 方法}}\label{alg:block-idr}\begin{algorithmic}[1]
\Statex\textbf{Input:} $X_0\in\mathbb C^{N\times s},\ R_0=B-AX_0,\ P\in\mathbb C^{n\times sm}$
\Statex\textbf{Output:} $X_i$
\For{$i=0$ \BT{to}{至} $s-1$}
\State $V=AR_i,\quad\omega=\frac{\operatorname{Tr}(V^{\mathrm H}R_i)}{\operatorname{Tr}(V^{\mathrm H}V)}$
\State $\Delta X(:,im+1:(i+1)m)=\omega R_i,\quad\Delta R(:,im+1:(i+1)m)=-\omega V$
\State $X_{i+1}=X_i+\Delta X(:,im+1:(i+1)m),\quad R_{i+1}=R_i+\Delta R(:,im+1:(i+1)m)$
\EndFor
\State $j=1,\quad i=s$
\State $M=P^{\mathrm H}\Delta R,\quad H=P^{\mathrm H}R_i$
\While{$\max_{\ell\in\{1,2,\ldots,m\}}\frac{\|R_i(:,\ell)\|_{\mathrm F}}{\|B(:,\ell)\|_{\mathrm F}}>\epsilon$}
\For{$k=0$ \BT{to}{至} $s$}
\State\BT{Solve $C$ from $MC=H$}{由 $MC=H$ 解出 $C$}
\State $Q=-\Delta RC,$
\State $V=R_i+Q$
\If{$k=0$}
\State $T=AV,\quad\omega=\frac{\operatorname{Tr}(T^{\mathrm H}V)}{\operatorname{Tr}(T^{\mathrm H}T)}$
\State $\Delta R(:,(j-1)m+1:jm)=Q-\omega T$
\State $\Delta X(:,(j-1)m+1:jm)=-\Delta XC+\omega V$
\Else
\State $\Delta X(:,(j-1)m+1:jm)=-\Delta XC+\omega V$
\State $\Delta R(:,(j-1)m+1:jm)=-A\Delta X(:,(j-1)m+1:jm)$
\EndIf
\State $R_{i+1}=R_i+\Delta R(:,(j-1)m+1:jm)$
\State $X_{i+1}=X_i+\Delta X(:,(j-1)m+1:jm)$
\State $\Delta M=P^{\mathrm H}\Delta R(:,(j-1)m+1:jm)$
\State $M(:,(j-1)m+1:jm)=\Delta M$
\State $H=H+\Delta M$
\State $i=i+1,\quad j=j+1$
\State $j=(j-1)\%s+1$ (\BT{$\%$: modulo operation, i.e.\ $a\%n=r$, where $a=mn+r$.}{$\%$: 取模运算, 即 $a\%n=r$, 其中 $a=mn+r$.})
\EndFor\EndWhile\end{algorithmic}\end{algorithm}
% block: preconditioned-block-idr-algorithm page:158 type:ordinary-paragraph
\begin{algorithm}[H]\caption{\BT{The preconditioned block IDR($s$) method}{预条件块 IDR($s$) 方法}}\label{alg:preconditioned-block-idr}\begin{algorithmic}[1]
\Statex\textbf{Input:} $X_0\in\mathbb C^{N\times s},\ R_0=B-AX_0,\ P\in\mathbb C^{n\times sm}$
\Statex\textbf{Input:} $X_0=KX_0$ ($K$: \BT{preconditioning matrix}{预条件矩阵})
\Statex\textbf{Output:} $X_i$
\For{$i=0$ \BT{to}{至} $s-1$}
\State $W=K^{-1}R_i$
\State $V=AW,\quad\omega=\frac{\operatorname{Tr}((K_1^{-1}V)^{\mathrm H}K_1^{-1}R_i)}{\operatorname{Tr}((K_1^{-1}V)^{\mathrm H}K_1^{-1}V)}$
\State $\Delta X(:,im+1:(i+1)m)=\omega R_i,\quad\Delta R(:,im+1:(i+1)m)=-\omega V$
\State $X_{i+1}=X_i+\Delta X(:,im+1:(i+1)m),\quad R_{i+1}=R_i+\Delta R(:,im+1:(i+1)m)$
\EndFor
\State $j=1,\quad i=s$
\State $M=P^{\mathrm H}\Delta R,\quad H=P^{\mathrm H}R_i$
\While{$\max_{\ell\in\{1,2,\ldots,m\}}\frac{\|R_i(:,\ell)\|_{\mathrm F}}{\|B(:,\ell)\|_{\mathrm F}}>\epsilon$}
\For{$k=0$ \BT{to}{至} $s$}
\State\BT{Solve $C$ from $MC=H$}{由 $MC=H$ 解出 $C$}
\State $Q=-\Delta RC$
\State $V=R_i+Q$
\If{$k=0$}
\State $T=AK^{-1}V,\quad\omega=\frac{\operatorname{Tr}((K_1^{-1}T)^{\mathrm H}K_1^{-1}V)}{\operatorname{Tr}((K_1^{-1}T)^{\mathrm H}K_1^{-1}T)}$
\State $\Delta R(:,(j-1)m+1:jm)=Q-\omega T$
\State $\Delta X(:,(j-1)m+1:jm)=-\Delta XC+\omega V$
\Else
\State $\Delta X(:,(j-1)m+1:jm)=-\Delta XC+\omega V$
\State $\Delta R(:,(j-1)m+1:jm)=-AK^{-1}\Delta X(:,(j-1)m+1:jm)$
\EndIf
\State $R_{i+1}=R_i+\Delta R(:,(j-1)m+1:jm)$
\State $X_{i+1}=X_i+\Delta X(:,(j-1)m+1:jm)$
\State $\Delta M=P^{\mathrm H}\Delta R(:,(j-1)m+1:jm)$
\State $M(:,(j-1)m+1:jm)=\Delta M$
\State $H=H+\Delta M$
\State $i=i+1,\quad j=j+1$
\State $j=(j-1)\%s+1$ (\BT{$\%$: modulo operation, i.e.\ $a\%n=r$, where $a=mn+r$.}{$\%$: 取模运算, 即 $a\%n=r$, 其中 $a=mn+r$.})
\EndFor\EndWhile
\State $X_i=K^{-1}X_i$
\end{algorithmic}\end{algorithm}
